% =====================================================================================================================
% DAN-208 (drafted 2026-11-08, agent "framework"): WCM as a framework for solving for circuits, not a fixed algorithm.
% Two versions in one file. Choose with \def\frameworkpart{main} or \def\frameworkpart{appendix} before \input:
%
%   \def\frameworkpart{main}% =====================================================================================================================
% DAN-208 (drafted 2026-11-08, agent "framework"): WCM as a framework for solving for circuits, not a fixed algorithm.
% Two versions in one file. Choose with \def\frameworkpart{main} or \def\frameworkpart{appendix} before \input:
%
%   \def\frameworkpart{main}% =====================================================================================================================
% DAN-208 (drafted 2026-11-08, agent "framework"): WCM as a framework for solving for circuits, not a fixed algorithm.
% Two versions in one file. Choose with \def\frameworkpart{main} or \def\frameworkpart{appendix} before \input:
%
%   \def\frameworkpart{main}% =====================================================================================================================
% DAN-208 (drafted 2026-11-08, agent "framework"): WCM as a framework for solving for circuits, not a fixed algorithm.
% Two versions in one file. Choose with \def\frameworkpart{main} or \def\frameworkpart{appendix} before \input:
%
%   \def\frameworkpart{main}\input{drafts/framework_section}        % ~1/3 page, Discussion (replaces the paragraph
%                                                                    % "Ablation methods are part of the method")
%   \def\frameworkpart{appendix}\input{drafts/framework_section}    % ~1.3 pages, its own appendix section
%
% New labels: sec:framework (main), app:framework, tab:framework (appendix). Every other \cref target exists in
% main.tex as of 2026-11-08. Every number below is already printed in main.tex (source comments give the place);
% the one exception, experiment 108, is marked NOT-YET-IN-PAPER where it appears.
% Vocabulary: paper/notes/vocabulary.md (no "endpoint", "floor", "drivers"; "fill" only as main.tex already uses it).
% Never-claim list respected: no "first"; nothing here claims novelty.
% =====================================================================================================================
\providecommand{\frameworkpart}{main}
\def\frameworkmain{main}
\def\frameworkappendix{appendix}
\ifx\frameworkpart\frameworkmain
% ---------------------------------------------------------------------------------------------------------------------
% (1) MAIN-TEXT VERSION. Proposed place: sec:discussion, in place of the paragraph "Ablation methods are part of the
% method" (main.tex lines 548-552), whose evidence it repeats and extends; the paragraph "Recommendations" (lines
% 564-566) can then point here. If it must sit in Method instead, it fits after "Pass criterion" (line 341), before
% \section{Experiments and Results}, but there it would cite results not yet shown.
% Numbers: 2-14% = abstract / run2026_alpha1 (sec:weighted); "one to three orders of magnitude" = app:single-method
% (exploratory); contrast sources = app:contrast-source (protocol v1); resample = app:resample.
% ---------------------------------------------------------------------------------------------------------------------
\subsection{Using WCM as a Framework}
\label{sec:framework}
WCM is better read as a framework for solving for circuits than as one fixed algorithm.
Its parts are interchangeable: an objective (the quantity reproduced, the ablation methods it is reproduced under, the regularisers, and their weights; \cref{eq:loss}), a parametrisation (a gate and a coefficient per node), data (training, contrast, and held-out contexts), and an optimiser.
The production recipe is one setting of these parts, each chosen by an experiment (\cref{tab:framework}), and changing a part changes the circuits found.
With every coefficient fixed at $1$ the same nodes pass for $2$--$14\%$ of targets (\cref{sec:weighted}); in exploratory fits, coefficients trained without zero ablation overshoot it by one to three orders of magnitude (\cref{app:single-method}); circuits trained against close contrast contexts are the most faithful on close backgrounds and the least on any other (\cref{app:contrast-source}); and circuits failed a resample test until training included a resample term (\cref{app:resample}).
When WCM fails an evaluation, we therefore suggest first asking whether training ever saw that evaluation's domain, and if not, adding the term, contexts, or setting that puts it there, while keeping the evaluation's own contexts and donors out of training so that it stays a held-out test (\cref{app:framework}).
A pass rate then describes a recipe, a dictionary, and an evaluation together, not WCM alone, which is why the ablation methods used for training and evaluation belong in the method statement~\citep{mueller2025mib, mahale2026explanation}.
%
\else\ifx\frameworkpart\frameworkappendix
% ---------------------------------------------------------------------------------------------------------------------
% (2) APPENDIX VERSION. Proposed place: a new appendix section directly before "Weighted Circuit Masking Design
% Ablations" (main.tex line 2003, app:membership-studies), so that the design ablations it summarises follow it; or
% directly after app:protocol / app:loss-details (before line 853) if Daniel prefers it next to tab:protocol.
% ---------------------------------------------------------------------------------------------------------------------
\section{Building with Weighted Circuit Masking}
\label{app:framework}
Weighted circuit masking fixes a problem, \labelcref{eq:p2}, and a way of relaxing it, but leaves most of the recipe open.
\Cref{tab:framework} lists its slots: what is reproduced, over which nodes, with which parametrisation, under which ablation methods and backgrounds, with which regularisers and optimiser, from which contexts, and how the result is scored.
The production choice in each slot was made by an experiment reported in these appendices, most of them on TuringLLM and carried over unchanged to the public models (\cref{tab:protocol}, Provenance); none is claimed to be the best setting, and several are known to trade one evaluation against another.

{\footnotesize
\setlength{\tabcolsep}{3pt}
\renewcommand{\arraystretch}{1.15}
\begin{longtable}{@{}p{0.13\linewidth}p{0.33\linewidth}p{0.33\linewidth}p{0.16\linewidth}@{}}
\caption{The slots of weighted circuit masking, the production choice in each, and the alternatives this paper examined or that the formulation admits.}
\label{tab:framework} \\
\toprule
\textbf{Slot} & \textbf{Production choice} & \textbf{Alternatives} & \textbf{Evidence} \\
\midrule
\endfirsthead
\toprule
\textbf{Slot} & \textbf{Production choice} & \textbf{Alternatives} & \textbf{Evidence} \\
\midrule
\endhead
\midrule
\multicolumn{4}{r@{}}{\emph{continued on the next page}} \\
\endfoot
\bottomrule
\endlastfoot
Target & the pre-activation $\preact$ of one SAE latent at its anchors; censored read under Z, C, A & an output metric, as in learned-mask circuits for outputs~\citep{bhaskar2024finding, caples2025scaling}; output- and latent-targeted circuits fail each other's objective (exploratory) & \cref{sec:ablmask}, \cref{app:secondary,app:logit-faith} \\
Nodes & every latent at every SAE site upstream of the target; SAE errors passed through; nodes recomputed & dense ReLU dictionaries (the method replicates; the random baseline needs its own design); members clamped (resample diagnostic) & \cref{sec:problem}, \cref{app:dense-relu,app:resample} \\
Parametrisation & gate $\sigma(\theta/T)$ and coefficient $\mathrm{softplus}(\psi)$ per node, $\alpha_0 = 1$ & $\alpha \equiv 1$ (unweighted circuit masking) & \cref{sec:weighted}, \cref{tab:run2026-alpha1} \\
Training ablations & Z (weight $1$), C and A ($0.25$ each) under evaluation's sparsity-preserving fill; resample D over the training contrast contexts ($1$) & one method alone; dense means; D in place of C; no D & \cref{app:single-method,app:respect-topk,app:resample} \\
Contrast background & $64$ close, silence-checked contexts per target ($48$ train, $16$ held out) & random, distant, or fused sources; the unchecked retrieval store & \cref{app:stores,app:contrast-source} \\
Regularisers & flat penalty $\lambda = 10^{-3}$ per open gate; closed-gate charge $\lambda(1-\sigma(\theta))|\alpha - 1|$; rank-keep $\gamma_{\mathrm{K}} = 10^{-1}$; one shared normaliser $\nu$ & other penalties $\lambda$ (sweep); off-target term, Top-K hinge, pull to rank $1$, firing margin; per-term normalisers & \cref{sec:lambda-transfer}, \cref{app:loss-details,app:rank-designs,app:membership-studies} \\
Optimiser & AdamW, $400$ steps, batch $4$; gate temperature annealed $1 \to 0.05$; learning-rate warm-up, then geometric decay & straight-through binarisation; $200$-step and reach-then-hold schedules; larger batches (exploratory) & \cref{tab:protocol}, \cref{app:membership-studies,app:framework-history} \\
Training contexts & $32$ strongest and $16$ mid-band activating contexts, deduplicated and split by whole groups of copies & a split without deduplication, which leaked copies into held-out sets & \cref{sec:setting}, \cref{app:heldout} \\
Evaluation & held-out strongest contexts; Z, A, C in $[0.75, 1.5]$ and necessity $\geq 0.9$; vacuous targets excluded and counted & reported, not required: mid-band contexts, sufficiency to induce, donor-background test, specificity, $\alpha = 1$ rescoring, random-circuit baselines, stable core & \cref{sec:eval}, \cref{app:protocol,app:nulls,app:underdetermination} \\
\end{longtable}}

\paragraph{Training in the evaluation's domain.}
Across these experiments, circuits did best on the evaluations their training resembled.
When a fit fails an evaluation, the failure is often a property of the recipe, and a slot of \cref{tab:framework} can put the fit in that evaluation's domain.
% app:resample (095: 72 targets, 0% -> 24%, CI +14 to +33; criterion 69% -> 62%, CI -17 to +3; 1.18x size, 1.45x time;
% D in place of C: 51%, CI -29 to -7), production paragraph (exp 102: 19 -> 39 and 48 -> 54 of 72, 1.33x size), T10
% (tab:run2026-tl-noR: +5.2, -0.6 to 11.6).
\emph{Real backgrounds.} Circuits trained only against zero and means failed a resample test almost always; adding a resample term over the training contrast contexts raised the resample pass rate on $72$ targets from $0\%$ to $24\%$, at $1.18\times$ the circuit size and $1.45\times$ the fitting time, and with rank-keep raised it took a $72$-target check panel from $19$ to $39$ targets passing (\cref{app:resample}).
Using it in place of the contrast-mean term gained as much on the resample test but cost the pass criterion, so terms should be added for a new evaluation rather than swapped for an old one; on the production sweep the term's effect on the pass criterion itself is $5.2$ points and not significant (\cref{tab:run2026-tl-noR}).
\emph{Ablation methods and fills.} A mask trained under one ablation method learns what that method rewards: without the zero-ablation term, coefficients overshoot zero ablation by orders of magnitude, and circuits trained against dense means can fail under the sparsity-preserving fill that evaluation uses (\cref{app:single-method,app:respect-topk}).
\emph{Top-K survival.} Reproducing the pre-activation does not keep the target in its site's Top-K; the censored read and rank-keep train for it, though neither removes amplifier circuits (\cref{app:loss-details,app:specificity}).
\emph{Other backgrounds.} Contrast faithfulness generalises furthest when training uses random or distant contrast contexts (\cref{app:contrast-source}).
\emph{Size.} Circuits shrink with $\lambda$ at a nearly constant pass rate up to a plateau's edge (\cref{sec:lambda-transfer}).
Two evaluations have no training term yet: sufficiency to induce, which no circuit is trained on (median $0.78$--$0.88$ on held-out contrast contexts, $0.35$ for TuringLLM attention outputs; \cref{app:production}), and transfer to weaker, mid-band contexts (\cref{sec:main-result}).
The framework's prescription for them is a term or a training set in their domain, which we have not tested.
% CANDIDATE EVIDENCE, NOT IN THE PAPER (main.tex line 503: the arena results are reported separately). If Daniel wants
% them here, each needs its own sentence and its own appendix or citation:
%  - arena A tuning (experiments/107-true-baselines/arena_a/TUNING.md): on GPT-2 transcoder + z-SAE substrate the
%    necessity hinge (NEC_W; off in production, docs/wcm-specification.md §2.3) was chosen at 0.3 by a pre-registered
%    9-target rule; NEC_W 1.0 raised nec_a (3.tc.3752 0.32 -> 0.95) but pushed A out of band (1.54). The censored read
%    is inert there by construction (k = d_sae).
%  - arena A scope a1 (z + z2 + tc) vs a2 (tc only) and the wikitext-smoke vs OpenWebText corpus test (DAN-208 text;
%    not traced to an output file in this pass).
%  - logit / maximise / behavioural objectives (arena C/D; memory notes say "don't quote yet").
A fit trained for an evaluation is, for that evaluation, scored on home turf (\cref{app:method-comparison}), so the evaluation stays meaningful only if its contexts and donors are never seen in training, as the donor-background test's donors are not (\cref{tab:protocol}), and the evaluations not trained for are reported alongside.

\paragraph{Choosing the sparsity penalty.}
% sec:lambda-transfer, tab:run2026-tl-sweep, tab:run2026-lambda-pilot.
The penalty sets circuit size and, past a point, the pass rate: on the TuringLLM sweep the pass rate holds at $89$--$93\%$ from $\lambda = 2.5 \times 10^{-5}$ to $5 \times 10^{-4}$ while each doubling of $\lambda$ shrinks circuits $1.4$--$1.9\times$, and it declines beyond (\cref{sec:lambda-transfer}).
We recommend fixing the selection rule before looking at pass rates, as our size-only pilot rule did on the public models, where it kept $\lambda = 10^{-3}$ on both (\cref{tab:run2026-lambda-pilot}), and reporting a curve over $\lambda$ rather than one point, since one setting that transfers is not the best setting for each model.

\paragraph{A new model or dictionary.}
Only the model and dictionary adapters change between our three models (\cref{sec:setting}); four things needed checking.
\emph{Precision}: TuringLLM runs in the bfloat16 its pipeline was validated in, but on the public models a bfloat16 data pass did not reproduce the float32 top-context sets, so they run in float32 (\cref{tab:protocol}); a precision check belongs before any fit, since node sets also move with rounding (\cref{app:underdetermination}).
\emph{The dictionary's sparsity}: $k$ sets the fill's budget and the censored read's cut; on a dictionary that keeps every latent ($k = d_{\mathrm{sae}}$) the censored read has no cut to model, and on dense ReLU dictionaries a random-circuit baseline with fitted coefficients stops discriminating, so necessity, not faithfulness, separates real from random circuits there (\cref{app:dense-relu}).
% The k = d_sae statement follows from the definition of tau^(K) (no K-th other latent); docs/wcm-specification.md §7.6.
\emph{Contrast retrieval}: contrast contexts need a pooled representation for every corpus sequence and a silence check, which rejects $14$--$51\%$ of the closest candidates on the targets we inspected (\cref{app:stores}).
\emph{Corpus copies}: duplicated windows silently leak into held-out sets unless pools are deduplicated and split by groups of copies (\cref{app:heldout}).

\paragraph{What to report.}
Because each slot changes the circuits, a result is interpretable only with its recipe: the ablation methods used in training and in evaluation; faithfulness under each method and necessity, on held-out contexts, by depth band and site kind, with intervals over targets; excluded and vacuous targets counted; circuit size with whether coefficients were fitted, and the $\alpha = 1$ rescoring; random-circuit baselines given the same fitting; and the reported-only evaluations (sufficiency to induce, the donor-background test, specificity, mid-band transfer) beside the pass rate (\cref{sec:discussion}).

\paragraph{Choices that matter, and what stays open.}
The ablation method changes not only scores but rankings: the two mean-ablation variants can order small node sets differently (\cref{app:respect-topk}), and in exploratory comparisons a method tended to win the scores of the ablation method it was trained or ranked under (\cref{app:method-comparison}), so comparisons between recipes need the full matrix of recipe against ablation method.
Node sets are not identified by any recipe: a refit with reordered training contexts shares a median Jaccard of $0.39$--$0.51$ with the original and its shared core alone passes for $44$--$72\%$ of targets (\cref{app:underdetermination}), so changes to a slot should be judged by sizes and scores, not by which latents moved.
Finally, the production choices were calibrated on TuringLLM targets that overlap its evaluation set (\cref{sec:setting}), and on the public models every setting except $\lambda$ carried over unchanged (\cref{tab:protocol}, Provenance), so on a new model they are defaults to be checked, not tuned values.

% ---------------------------------------------------------------------------------------------------------------------
% 2026-11-08 (Daniel, via coordinator): the development path as evidence for where WCM can go next. Every number is
% traced; sources in the comment before each step. Panels and protocols differ between steps, so numbers are compared
% only within a step (each step is a paired test against the recipe before it).
% ---------------------------------------------------------------------------------------------------------------------
\subsection{How the Production Recipe Was Reached}
\label{app:framework-history}
The production recipe is the end of a sequence of changes, each made because a run exposed a failure, and each tested against the recipe before it on the same targets, mostly under a selection rule stated before the comparison.
We list them in order as evidence of which kinds of change have paid off; the panels differ between steps, so numbers compare only within a step.

% app:stores (38% fallback), app:contrast-source; experiments/055-close-contrast, 071-contrast-sweep.
\paragraph{Contrast contexts built per target.}
The first contrast contexts came from a retrieval store whose index covered too few sequences, so $38\%$ of TuringLLM targets fell back to one shared list; WCM now selects contrast contexts per target at fit time, ranked by similarity to its strongest contexts and kept only if the target is silent (\cref{app:stores}), and a later comparison of contrast sources showed how much that choice shapes the C term (\cref{app:contrast-source}).
The lesson: inspect the data each training term actually sees.

% app:specificity, app:rank-designs; experiments/056-specificity.
\paragraph{Rank-keep against off-target lifts.}
Reproducing the pre-activation left some circuits restoring their target by lifting its whole neighbourhood; of the terms tried, rank-keep alone kept the gains without the alternatives' costs, and entered the first protocol at a small weight (\cref{app:rank-designs,app:specificity}).

% app:v1-run (15,046 targets, 46%, attention 14%, layers 8-11 22%); app:heldout (66.8%); app:resample (0.1% vs 43.1%,
% 747 targets); fig:failure-anatomy.
\paragraph{The first run at scale.}
Under that protocol, WCM was fitted to $15{,}046$ TuringLLM targets, and $46\%$ passed, $22\%$ at layers 8--11 and $14\%$ for attention outputs (\cref{app:v1-run}).
Analyses of the run exposed three problems: $66.8\%$ of held-out strongest contexts repeated a training context, which led to deduplicated records split by groups of copies (\cref{app:heldout}); with depth, failing circuits increasingly failed only the two mean ablations (\cref{fig:failure-anatomy}); and $0.1\%$ of circuits passed a resample test, against $43.1\%$ under the pass criterion (\cref{app:resample}).

% experiments/093-v2-validation/README.md (out/tables.md): 72 targets, 8 per band x kind, paired, pass floor 0.8, the
% v1 run's as-run held-out sets (leaky; same for both arms): 31 -> 50 of 72, +26.4 pp [15.3, 38.9]; L8-11 2 -> 11 of 24;
% attn 3 -> 12 of 24; median dZ 0.00; size ratio 0.91; fit time 1.27x. Pilot: experiments/090-sparse-fill-training.
% Also main.tex comment at app:respect-topk ("The same fill in training").
\paragraph{Training under the fill that evaluation uses.}
The mean-ablation failures were a mismatch: training replaced non-members with dense means, a code no Top-K dictionary emits, while evaluation re-imposed the Top-K (\cref{app:respect-topk}).
Training the A and C terms under evaluation's own fill and means raised the pass rate on $72$ stratified targets from $31$ to $50$ (paired difference $+26.4$ points, $95\%$ interval $15.3$--$38.9$), from $2$ to $11$ of $24$ at layers 8--11 and from $3$ to $12$ of $24$ for attention outputs, with zero-ablation faithfulness unchanged (median paired change $0.00$), circuits $0.91\times$ the size, and $1.27\times$ the fitting time.
The gain came from aligning training with evaluation, not from a new term.

% app:resample (095 numbers and the production paragraph: 102's 19 -> 39, 48 -> 54 of 72, 1.33x size);
% experiments/102-pass-improvement/README.md line 89 (diagnosis: 15 of 18 Z failures of the previous version restored
% the pre-activation and lost at the Top-K cut); tab:run2026-tl-noR.
\paragraph{A resample term, and a stronger rank-keep.}
To survive real backgrounds, training gained the resample term (above and \cref{app:resample}).
On the previous version's zero-ablation failures, $15$ of $18$ restored the target's pre-activation but lost it at the Top-K cut, so rank-keep was raised from $3 \times 10^{-3}$ to $10^{-1}$ with the resample term at full weight; chosen by a rule stated before the comparison, the pair took a $72$-target check panel from $48$ to $54$ targets passing and from $19$ to $39$ passing the resample test, at $1.33\times$ the circuit size (\cref{app:resample}).
In the production sweep, the arm trained without the resample term passes $5.2$ points less, not significantly (\cref{tab:run2026-tl-noR}).

% experiments/107-true-baselines/latent_pilot/README.md (9 targets enriched for failures: 3 -> 5 of 9, size ratio ~1.04;
% plain geometric: circuits 2-4x larger, sum(lr) 6.3 vs 20; constant lr makes lr/T grow 20x) and
% experiments/107-true-baselines/adaptive_lambda/README.md, "Check" table (arm "geo x 3.15 + censored term
% (latent_pilot)": 63 of 72 vs 54 for the previous version, +10 / -1, size ratio 1.04 on pass-in-both). VERIFY: whether
% that check arm included the 0.1 warm-up (run_tonight.sh stage 4 fits check_geo_x3.15_warmup_censor).
\paragraph{A learning-rate schedule and the censored read.}
Under a constant learning rate, the gate step grows twentyfold as the gate temperature anneals, and some fits degraded late in training.
A geometric decay alone stabilised them but left gates too little total step to close, so circuits grew $2$--$4\times$; raising its peak $3.15\times$ restored the summed step size, and a short warm-up was added for the early spikes the higher peak brought back (\cref{app:loss-details}).
With the censored read added, so that training also sees whether the target survives its site's Top-K, the check panel went from $54$ to $63$ of $72$ targets passing ($10$ gained, $1$ lost) at about the same circuit size (median ratio $1.04$).

% sec:lambda-transfer, tab:run2026-lambda-pilot; experiments/108-batch-size/out/analysis.md (16 TuringLLM targets;
% B16 at 400 steps: 15 / 16 vs 16 / 16, size 0.75x on pass-in-both, 99.9 vs 29.1 fit-minutes = 3.4x; pre-stated rule:
% +-1 target is noise); experiments/107-true-baselines/adaptive_lambda/README.md (every shrinking policy traded passes
% for size on the check panel; losses were necessity failures).
\paragraph{Settings left as they were.}
The penalty was carried to the public models by a size-only rule fixed before a pilot (\cref{sec:lambda-transfer}).
In exploratory tests, batches of $16$ for the same $400$ steps passed $15$ of $16$ TuringLLM targets against $16$, with circuits $0.75\times$ the size where both passed, at $3.4\times$ the fitting time, so the batch of $4$ was kept; and policies that adapted $\lambda$ per target traded passes for size on the check panel, losing passes mostly to necessity, which no training term covers.
Other rejected designs are in \cref{app:membership-studies}.

% Directions: app:production (induce), sec:main-result (mid-band, deep band, attention), app:specificity (amplifiers,
% v1: 30% flagged), app:rank-designs (off-target term: lifts cut 3-7x, but grew circuits), app:resample (anchors
% only; overfiring away from the anchor), adaptive_lambda README (necessity).
\paragraph{Directions not yet taken.}
The same path suggests where the recipe can go next.
Necessity, sufficiency to induce, and transfer to mid-band contexts are scored but not trained for; each could get a term or training contexts in its domain, as the resample test did.
Deep and attention-output targets still pass least (\cref{sec:main-result}), and in the schedule pilot the attention failures kept zero-ablation faithfulness near $0$ under every schedule, so the learning rate does not explain them.
Amplifier circuits survive rank-keep (\cref{app:specificity}); the off-target term cut lifts three- to sevenfold but grew circuits (\cref{app:rank-designs}), so a cheaper specificity term remains open.
And circuits explain the target at its anchors, not at every position (\cref{app:resample}); training on more positions is untested.
\fi\fi
        % ~1/3 page, Discussion (replaces the paragraph
%                                                                    % "Ablation methods are part of the method")
%   \def\frameworkpart{appendix}% =====================================================================================================================
% DAN-208 (drafted 2026-11-08, agent "framework"): WCM as a framework for solving for circuits, not a fixed algorithm.
% Two versions in one file. Choose with \def\frameworkpart{main} or \def\frameworkpart{appendix} before \input:
%
%   \def\frameworkpart{main}\input{drafts/framework_section}        % ~1/3 page, Discussion (replaces the paragraph
%                                                                    % "Ablation methods are part of the method")
%   \def\frameworkpart{appendix}\input{drafts/framework_section}    % ~1.3 pages, its own appendix section
%
% New labels: sec:framework (main), app:framework, tab:framework (appendix). Every other \cref target exists in
% main.tex as of 2026-11-08. Every number below is already printed in main.tex (source comments give the place);
% the one exception, experiment 108, is marked NOT-YET-IN-PAPER where it appears.
% Vocabulary: paper/notes/vocabulary.md (no "endpoint", "floor", "drivers"; "fill" only as main.tex already uses it).
% Never-claim list respected: no "first"; nothing here claims novelty.
% =====================================================================================================================
\providecommand{\frameworkpart}{main}
\def\frameworkmain{main}
\def\frameworkappendix{appendix}
\ifx\frameworkpart\frameworkmain
% ---------------------------------------------------------------------------------------------------------------------
% (1) MAIN-TEXT VERSION. Proposed place: sec:discussion, in place of the paragraph "Ablation methods are part of the
% method" (main.tex lines 548-552), whose evidence it repeats and extends; the paragraph "Recommendations" (lines
% 564-566) can then point here. If it must sit in Method instead, it fits after "Pass criterion" (line 341), before
% \section{Experiments and Results}, but there it would cite results not yet shown.
% Numbers: 2-14% = abstract / run2026_alpha1 (sec:weighted); "one to three orders of magnitude" = app:single-method
% (exploratory); contrast sources = app:contrast-source (protocol v1); resample = app:resample.
% ---------------------------------------------------------------------------------------------------------------------
\subsection{Using WCM as a Framework}
\label{sec:framework}
WCM is better read as a framework for solving for circuits than as one fixed algorithm.
Its parts are interchangeable: an objective (the quantity reproduced, the ablation methods it is reproduced under, the regularisers, and their weights; \cref{eq:loss}), a parametrisation (a gate and a coefficient per node), data (training, contrast, and held-out contexts), and an optimiser.
The production recipe is one setting of these parts, each chosen by an experiment (\cref{tab:framework}), and changing a part changes the circuits found.
With every coefficient fixed at $1$ the same nodes pass for $2$--$14\%$ of targets (\cref{sec:weighted}); in exploratory fits, coefficients trained without zero ablation overshoot it by one to three orders of magnitude (\cref{app:single-method}); circuits trained against close contrast contexts are the most faithful on close backgrounds and the least on any other (\cref{app:contrast-source}); and circuits failed a resample test until training included a resample term (\cref{app:resample}).
When WCM fails an evaluation, we therefore suggest first asking whether training ever saw that evaluation's domain, and if not, adding the term, contexts, or setting that puts it there, while keeping the evaluation's own contexts and donors out of training so that it stays a held-out test (\cref{app:framework}).
A pass rate then describes a recipe, a dictionary, and an evaluation together, not WCM alone, which is why the ablation methods used for training and evaluation belong in the method statement~\citep{mueller2025mib, mahale2026explanation}.
%
\else\ifx\frameworkpart\frameworkappendix
% ---------------------------------------------------------------------------------------------------------------------
% (2) APPENDIX VERSION. Proposed place: a new appendix section directly before "Weighted Circuit Masking Design
% Ablations" (main.tex line 2003, app:membership-studies), so that the design ablations it summarises follow it; or
% directly after app:protocol / app:loss-details (before line 853) if Daniel prefers it next to tab:protocol.
% ---------------------------------------------------------------------------------------------------------------------
\section{Building with Weighted Circuit Masking}
\label{app:framework}
Weighted circuit masking fixes a problem, \labelcref{eq:p2}, and a way of relaxing it, but leaves most of the recipe open.
\Cref{tab:framework} lists its slots: what is reproduced, over which nodes, with which parametrisation, under which ablation methods and backgrounds, with which regularisers and optimiser, from which contexts, and how the result is scored.
The production choice in each slot was made by an experiment reported in these appendices, most of them on TuringLLM and carried over unchanged to the public models (\cref{tab:protocol}, Provenance); none is claimed to be the best setting, and several are known to trade one evaluation against another.

{\footnotesize
\setlength{\tabcolsep}{3pt}
\renewcommand{\arraystretch}{1.15}
\begin{longtable}{@{}p{0.13\linewidth}p{0.33\linewidth}p{0.33\linewidth}p{0.16\linewidth}@{}}
\caption{The slots of weighted circuit masking, the production choice in each, and the alternatives this paper examined or that the formulation admits.}
\label{tab:framework} \\
\toprule
\textbf{Slot} & \textbf{Production choice} & \textbf{Alternatives} & \textbf{Evidence} \\
\midrule
\endfirsthead
\toprule
\textbf{Slot} & \textbf{Production choice} & \textbf{Alternatives} & \textbf{Evidence} \\
\midrule
\endhead
\midrule
\multicolumn{4}{r@{}}{\emph{continued on the next page}} \\
\endfoot
\bottomrule
\endlastfoot
Target & the pre-activation $\preact$ of one SAE latent at its anchors; censored read under Z, C, A & an output metric, as in learned-mask circuits for outputs~\citep{bhaskar2024finding, caples2025scaling}; output- and latent-targeted circuits fail each other's objective (exploratory) & \cref{sec:ablmask}, \cref{app:secondary,app:logit-faith} \\
Nodes & every latent at every SAE site upstream of the target; SAE errors passed through; nodes recomputed & dense ReLU dictionaries (the method replicates; the random baseline needs its own design); members clamped (resample diagnostic) & \cref{sec:problem}, \cref{app:dense-relu,app:resample} \\
Parametrisation & gate $\sigma(\theta/T)$ and coefficient $\mathrm{softplus}(\psi)$ per node, $\alpha_0 = 1$ & $\alpha \equiv 1$ (unweighted circuit masking) & \cref{sec:weighted}, \cref{tab:run2026-alpha1} \\
Training ablations & Z (weight $1$), C and A ($0.25$ each) under evaluation's sparsity-preserving fill; resample D over the training contrast contexts ($1$) & one method alone; dense means; D in place of C; no D & \cref{app:single-method,app:respect-topk,app:resample} \\
Contrast background & $64$ close, silence-checked contexts per target ($48$ train, $16$ held out) & random, distant, or fused sources; the unchecked retrieval store & \cref{app:stores,app:contrast-source} \\
Regularisers & flat penalty $\lambda = 10^{-3}$ per open gate; closed-gate charge $\lambda(1-\sigma(\theta))|\alpha - 1|$; rank-keep $\gamma_{\mathrm{K}} = 10^{-1}$; one shared normaliser $\nu$ & other penalties $\lambda$ (sweep); off-target term, Top-K hinge, pull to rank $1$, firing margin; per-term normalisers & \cref{sec:lambda-transfer}, \cref{app:loss-details,app:rank-designs,app:membership-studies} \\
Optimiser & AdamW, $400$ steps, batch $4$; gate temperature annealed $1 \to 0.05$; learning-rate warm-up, then geometric decay & straight-through binarisation; $200$-step and reach-then-hold schedules; larger batches (exploratory) & \cref{tab:protocol}, \cref{app:membership-studies,app:framework-history} \\
Training contexts & $32$ strongest and $16$ mid-band activating contexts, deduplicated and split by whole groups of copies & a split without deduplication, which leaked copies into held-out sets & \cref{sec:setting}, \cref{app:heldout} \\
Evaluation & held-out strongest contexts; Z, A, C in $[0.75, 1.5]$ and necessity $\geq 0.9$; vacuous targets excluded and counted & reported, not required: mid-band contexts, sufficiency to induce, donor-background test, specificity, $\alpha = 1$ rescoring, random-circuit baselines, stable core & \cref{sec:eval}, \cref{app:protocol,app:nulls,app:underdetermination} \\
\end{longtable}}

\paragraph{Training in the evaluation's domain.}
Across these experiments, circuits did best on the evaluations their training resembled.
When a fit fails an evaluation, the failure is often a property of the recipe, and a slot of \cref{tab:framework} can put the fit in that evaluation's domain.
% app:resample (095: 72 targets, 0% -> 24%, CI +14 to +33; criterion 69% -> 62%, CI -17 to +3; 1.18x size, 1.45x time;
% D in place of C: 51%, CI -29 to -7), production paragraph (exp 102: 19 -> 39 and 48 -> 54 of 72, 1.33x size), T10
% (tab:run2026-tl-noR: +5.2, -0.6 to 11.6).
\emph{Real backgrounds.} Circuits trained only against zero and means failed a resample test almost always; adding a resample term over the training contrast contexts raised the resample pass rate on $72$ targets from $0\%$ to $24\%$, at $1.18\times$ the circuit size and $1.45\times$ the fitting time, and with rank-keep raised it took a $72$-target check panel from $19$ to $39$ targets passing (\cref{app:resample}).
Using it in place of the contrast-mean term gained as much on the resample test but cost the pass criterion, so terms should be added for a new evaluation rather than swapped for an old one; on the production sweep the term's effect on the pass criterion itself is $5.2$ points and not significant (\cref{tab:run2026-tl-noR}).
\emph{Ablation methods and fills.} A mask trained under one ablation method learns what that method rewards: without the zero-ablation term, coefficients overshoot zero ablation by orders of magnitude, and circuits trained against dense means can fail under the sparsity-preserving fill that evaluation uses (\cref{app:single-method,app:respect-topk}).
\emph{Top-K survival.} Reproducing the pre-activation does not keep the target in its site's Top-K; the censored read and rank-keep train for it, though neither removes amplifier circuits (\cref{app:loss-details,app:specificity}).
\emph{Other backgrounds.} Contrast faithfulness generalises furthest when training uses random or distant contrast contexts (\cref{app:contrast-source}).
\emph{Size.} Circuits shrink with $\lambda$ at a nearly constant pass rate up to a plateau's edge (\cref{sec:lambda-transfer}).
Two evaluations have no training term yet: sufficiency to induce, which no circuit is trained on (median $0.78$--$0.88$ on held-out contrast contexts, $0.35$ for TuringLLM attention outputs; \cref{app:production}), and transfer to weaker, mid-band contexts (\cref{sec:main-result}).
The framework's prescription for them is a term or a training set in their domain, which we have not tested.
% CANDIDATE EVIDENCE, NOT IN THE PAPER (main.tex line 503: the arena results are reported separately). If Daniel wants
% them here, each needs its own sentence and its own appendix or citation:
%  - arena A tuning (experiments/107-true-baselines/arena_a/TUNING.md): on GPT-2 transcoder + z-SAE substrate the
%    necessity hinge (NEC_W; off in production, docs/wcm-specification.md §2.3) was chosen at 0.3 by a pre-registered
%    9-target rule; NEC_W 1.0 raised nec_a (3.tc.3752 0.32 -> 0.95) but pushed A out of band (1.54). The censored read
%    is inert there by construction (k = d_sae).
%  - arena A scope a1 (z + z2 + tc) vs a2 (tc only) and the wikitext-smoke vs OpenWebText corpus test (DAN-208 text;
%    not traced to an output file in this pass).
%  - logit / maximise / behavioural objectives (arena C/D; memory notes say "don't quote yet").
A fit trained for an evaluation is, for that evaluation, scored on home turf (\cref{app:method-comparison}), so the evaluation stays meaningful only if its contexts and donors are never seen in training, as the donor-background test's donors are not (\cref{tab:protocol}), and the evaluations not trained for are reported alongside.

\paragraph{Choosing the sparsity penalty.}
% sec:lambda-transfer, tab:run2026-tl-sweep, tab:run2026-lambda-pilot.
The penalty sets circuit size and, past a point, the pass rate: on the TuringLLM sweep the pass rate holds at $89$--$93\%$ from $\lambda = 2.5 \times 10^{-5}$ to $5 \times 10^{-4}$ while each doubling of $\lambda$ shrinks circuits $1.4$--$1.9\times$, and it declines beyond (\cref{sec:lambda-transfer}).
We recommend fixing the selection rule before looking at pass rates, as our size-only pilot rule did on the public models, where it kept $\lambda = 10^{-3}$ on both (\cref{tab:run2026-lambda-pilot}), and reporting a curve over $\lambda$ rather than one point, since one setting that transfers is not the best setting for each model.

\paragraph{A new model or dictionary.}
Only the model and dictionary adapters change between our three models (\cref{sec:setting}); four things needed checking.
\emph{Precision}: TuringLLM runs in the bfloat16 its pipeline was validated in, but on the public models a bfloat16 data pass did not reproduce the float32 top-context sets, so they run in float32 (\cref{tab:protocol}); a precision check belongs before any fit, since node sets also move with rounding (\cref{app:underdetermination}).
\emph{The dictionary's sparsity}: $k$ sets the fill's budget and the censored read's cut; on a dictionary that keeps every latent ($k = d_{\mathrm{sae}}$) the censored read has no cut to model, and on dense ReLU dictionaries a random-circuit baseline with fitted coefficients stops discriminating, so necessity, not faithfulness, separates real from random circuits there (\cref{app:dense-relu}).
% The k = d_sae statement follows from the definition of tau^(K) (no K-th other latent); docs/wcm-specification.md §7.6.
\emph{Contrast retrieval}: contrast contexts need a pooled representation for every corpus sequence and a silence check, which rejects $14$--$51\%$ of the closest candidates on the targets we inspected (\cref{app:stores}).
\emph{Corpus copies}: duplicated windows silently leak into held-out sets unless pools are deduplicated and split by groups of copies (\cref{app:heldout}).

\paragraph{What to report.}
Because each slot changes the circuits, a result is interpretable only with its recipe: the ablation methods used in training and in evaluation; faithfulness under each method and necessity, on held-out contexts, by depth band and site kind, with intervals over targets; excluded and vacuous targets counted; circuit size with whether coefficients were fitted, and the $\alpha = 1$ rescoring; random-circuit baselines given the same fitting; and the reported-only evaluations (sufficiency to induce, the donor-background test, specificity, mid-band transfer) beside the pass rate (\cref{sec:discussion}).

\paragraph{Choices that matter, and what stays open.}
The ablation method changes not only scores but rankings: the two mean-ablation variants can order small node sets differently (\cref{app:respect-topk}), and in exploratory comparisons a method tended to win the scores of the ablation method it was trained or ranked under (\cref{app:method-comparison}), so comparisons between recipes need the full matrix of recipe against ablation method.
Node sets are not identified by any recipe: a refit with reordered training contexts shares a median Jaccard of $0.39$--$0.51$ with the original and its shared core alone passes for $44$--$72\%$ of targets (\cref{app:underdetermination}), so changes to a slot should be judged by sizes and scores, not by which latents moved.
Finally, the production choices were calibrated on TuringLLM targets that overlap its evaluation set (\cref{sec:setting}), and on the public models every setting except $\lambda$ carried over unchanged (\cref{tab:protocol}, Provenance), so on a new model they are defaults to be checked, not tuned values.

% ---------------------------------------------------------------------------------------------------------------------
% 2026-11-08 (Daniel, via coordinator): the development path as evidence for where WCM can go next. Every number is
% traced; sources in the comment before each step. Panels and protocols differ between steps, so numbers are compared
% only within a step (each step is a paired test against the recipe before it).
% ---------------------------------------------------------------------------------------------------------------------
\subsection{How the Production Recipe Was Reached}
\label{app:framework-history}
The production recipe is the end of a sequence of changes, each made because a run exposed a failure, and each tested against the recipe before it on the same targets, mostly under a selection rule stated before the comparison.
We list them in order as evidence of which kinds of change have paid off; the panels differ between steps, so numbers compare only within a step.

% app:stores (38% fallback), app:contrast-source; experiments/055-close-contrast, 071-contrast-sweep.
\paragraph{Contrast contexts built per target.}
The first contrast contexts came from a retrieval store whose index covered too few sequences, so $38\%$ of TuringLLM targets fell back to one shared list; WCM now selects contrast contexts per target at fit time, ranked by similarity to its strongest contexts and kept only if the target is silent (\cref{app:stores}), and a later comparison of contrast sources showed how much that choice shapes the C term (\cref{app:contrast-source}).
The lesson: inspect the data each training term actually sees.

% app:specificity, app:rank-designs; experiments/056-specificity.
\paragraph{Rank-keep against off-target lifts.}
Reproducing the pre-activation left some circuits restoring their target by lifting its whole neighbourhood; of the terms tried, rank-keep alone kept the gains without the alternatives' costs, and entered the first protocol at a small weight (\cref{app:rank-designs,app:specificity}).

% app:v1-run (15,046 targets, 46%, attention 14%, layers 8-11 22%); app:heldout (66.8%); app:resample (0.1% vs 43.1%,
% 747 targets); fig:failure-anatomy.
\paragraph{The first run at scale.}
Under that protocol, WCM was fitted to $15{,}046$ TuringLLM targets, and $46\%$ passed, $22\%$ at layers 8--11 and $14\%$ for attention outputs (\cref{app:v1-run}).
Analyses of the run exposed three problems: $66.8\%$ of held-out strongest contexts repeated a training context, which led to deduplicated records split by groups of copies (\cref{app:heldout}); with depth, failing circuits increasingly failed only the two mean ablations (\cref{fig:failure-anatomy}); and $0.1\%$ of circuits passed a resample test, against $43.1\%$ under the pass criterion (\cref{app:resample}).

% experiments/093-v2-validation/README.md (out/tables.md): 72 targets, 8 per band x kind, paired, pass floor 0.8, the
% v1 run's as-run held-out sets (leaky; same for both arms): 31 -> 50 of 72, +26.4 pp [15.3, 38.9]; L8-11 2 -> 11 of 24;
% attn 3 -> 12 of 24; median dZ 0.00; size ratio 0.91; fit time 1.27x. Pilot: experiments/090-sparse-fill-training.
% Also main.tex comment at app:respect-topk ("The same fill in training").
\paragraph{Training under the fill that evaluation uses.}
The mean-ablation failures were a mismatch: training replaced non-members with dense means, a code no Top-K dictionary emits, while evaluation re-imposed the Top-K (\cref{app:respect-topk}).
Training the A and C terms under evaluation's own fill and means raised the pass rate on $72$ stratified targets from $31$ to $50$ (paired difference $+26.4$ points, $95\%$ interval $15.3$--$38.9$), from $2$ to $11$ of $24$ at layers 8--11 and from $3$ to $12$ of $24$ for attention outputs, with zero-ablation faithfulness unchanged (median paired change $0.00$), circuits $0.91\times$ the size, and $1.27\times$ the fitting time.
The gain came from aligning training with evaluation, not from a new term.

% app:resample (095 numbers and the production paragraph: 102's 19 -> 39, 48 -> 54 of 72, 1.33x size);
% experiments/102-pass-improvement/README.md line 89 (diagnosis: 15 of 18 Z failures of the previous version restored
% the pre-activation and lost at the Top-K cut); tab:run2026-tl-noR.
\paragraph{A resample term, and a stronger rank-keep.}
To survive real backgrounds, training gained the resample term (above and \cref{app:resample}).
On the previous version's zero-ablation failures, $15$ of $18$ restored the target's pre-activation but lost it at the Top-K cut, so rank-keep was raised from $3 \times 10^{-3}$ to $10^{-1}$ with the resample term at full weight; chosen by a rule stated before the comparison, the pair took a $72$-target check panel from $48$ to $54$ targets passing and from $19$ to $39$ passing the resample test, at $1.33\times$ the circuit size (\cref{app:resample}).
In the production sweep, the arm trained without the resample term passes $5.2$ points less, not significantly (\cref{tab:run2026-tl-noR}).

% experiments/107-true-baselines/latent_pilot/README.md (9 targets enriched for failures: 3 -> 5 of 9, size ratio ~1.04;
% plain geometric: circuits 2-4x larger, sum(lr) 6.3 vs 20; constant lr makes lr/T grow 20x) and
% experiments/107-true-baselines/adaptive_lambda/README.md, "Check" table (arm "geo x 3.15 + censored term
% (latent_pilot)": 63 of 72 vs 54 for the previous version, +10 / -1, size ratio 1.04 on pass-in-both). VERIFY: whether
% that check arm included the 0.1 warm-up (run_tonight.sh stage 4 fits check_geo_x3.15_warmup_censor).
\paragraph{A learning-rate schedule and the censored read.}
Under a constant learning rate, the gate step grows twentyfold as the gate temperature anneals, and some fits degraded late in training.
A geometric decay alone stabilised them but left gates too little total step to close, so circuits grew $2$--$4\times$; raising its peak $3.15\times$ restored the summed step size, and a short warm-up was added for the early spikes the higher peak brought back (\cref{app:loss-details}).
With the censored read added, so that training also sees whether the target survives its site's Top-K, the check panel went from $54$ to $63$ of $72$ targets passing ($10$ gained, $1$ lost) at about the same circuit size (median ratio $1.04$).

% sec:lambda-transfer, tab:run2026-lambda-pilot; experiments/108-batch-size/out/analysis.md (16 TuringLLM targets;
% B16 at 400 steps: 15 / 16 vs 16 / 16, size 0.75x on pass-in-both, 99.9 vs 29.1 fit-minutes = 3.4x; pre-stated rule:
% +-1 target is noise); experiments/107-true-baselines/adaptive_lambda/README.md (every shrinking policy traded passes
% for size on the check panel; losses were necessity failures).
\paragraph{Settings left as they were.}
The penalty was carried to the public models by a size-only rule fixed before a pilot (\cref{sec:lambda-transfer}).
In exploratory tests, batches of $16$ for the same $400$ steps passed $15$ of $16$ TuringLLM targets against $16$, with circuits $0.75\times$ the size where both passed, at $3.4\times$ the fitting time, so the batch of $4$ was kept; and policies that adapted $\lambda$ per target traded passes for size on the check panel, losing passes mostly to necessity, which no training term covers.
Other rejected designs are in \cref{app:membership-studies}.

% Directions: app:production (induce), sec:main-result (mid-band, deep band, attention), app:specificity (amplifiers,
% v1: 30% flagged), app:rank-designs (off-target term: lifts cut 3-7x, but grew circuits), app:resample (anchors
% only; overfiring away from the anchor), adaptive_lambda README (necessity).
\paragraph{Directions not yet taken.}
The same path suggests where the recipe can go next.
Necessity, sufficiency to induce, and transfer to mid-band contexts are scored but not trained for; each could get a term or training contexts in its domain, as the resample test did.
Deep and attention-output targets still pass least (\cref{sec:main-result}), and in the schedule pilot the attention failures kept zero-ablation faithfulness near $0$ under every schedule, so the learning rate does not explain them.
Amplifier circuits survive rank-keep (\cref{app:specificity}); the off-target term cut lifts three- to sevenfold but grew circuits (\cref{app:rank-designs}), so a cheaper specificity term remains open.
And circuits explain the target at its anchors, not at every position (\cref{app:resample}); training on more positions is untested.
\fi\fi
    % ~1.3 pages, its own appendix section
%
% New labels: sec:framework (main), app:framework, tab:framework (appendix). Every other \cref target exists in
% main.tex as of 2026-11-08. Every number below is already printed in main.tex (source comments give the place);
% the one exception, experiment 108, is marked NOT-YET-IN-PAPER where it appears.
% Vocabulary: paper/notes/vocabulary.md (no "endpoint", "floor", "drivers"; "fill" only as main.tex already uses it).
% Never-claim list respected: no "first"; nothing here claims novelty.
% =====================================================================================================================
\providecommand{\frameworkpart}{main}
\def\frameworkmain{main}
\def\frameworkappendix{appendix}
\ifx\frameworkpart\frameworkmain
% ---------------------------------------------------------------------------------------------------------------------
% (1) MAIN-TEXT VERSION. Proposed place: sec:discussion, in place of the paragraph "Ablation methods are part of the
% method" (main.tex lines 548-552), whose evidence it repeats and extends; the paragraph "Recommendations" (lines
% 564-566) can then point here. If it must sit in Method instead, it fits after "Pass criterion" (line 341), before
% \section{Experiments and Results}, but there it would cite results not yet shown.
% Numbers: 2-14% = abstract / run2026_alpha1 (sec:weighted); "one to three orders of magnitude" = app:single-method
% (exploratory); contrast sources = app:contrast-source (protocol v1); resample = app:resample.
% ---------------------------------------------------------------------------------------------------------------------
\subsection{Using WCM as a Framework}
\label{sec:framework}
WCM is better read as a framework for solving for circuits than as one fixed algorithm.
Its parts are interchangeable: an objective (the quantity reproduced, the ablation methods it is reproduced under, the regularisers, and their weights; \cref{eq:loss}), a parametrisation (a gate and a coefficient per node), data (training, contrast, and held-out contexts), and an optimiser.
The production recipe is one setting of these parts, each chosen by an experiment (\cref{tab:framework}), and changing a part changes the circuits found.
With every coefficient fixed at $1$ the same nodes pass for $2$--$14\%$ of targets (\cref{sec:weighted}); in exploratory fits, coefficients trained without zero ablation overshoot it by one to three orders of magnitude (\cref{app:single-method}); circuits trained against close contrast contexts are the most faithful on close backgrounds and the least on any other (\cref{app:contrast-source}); and circuits failed a resample test until training included a resample term (\cref{app:resample}).
When WCM fails an evaluation, we therefore suggest first asking whether training ever saw that evaluation's domain, and if not, adding the term, contexts, or setting that puts it there, while keeping the evaluation's own contexts and donors out of training so that it stays a held-out test (\cref{app:framework}).
A pass rate then describes a recipe, a dictionary, and an evaluation together, not WCM alone, which is why the ablation methods used for training and evaluation belong in the method statement~\citep{mueller2025mib, mahale2026explanation}.
%
\else\ifx\frameworkpart\frameworkappendix
% ---------------------------------------------------------------------------------------------------------------------
% (2) APPENDIX VERSION. Proposed place: a new appendix section directly before "Weighted Circuit Masking Design
% Ablations" (main.tex line 2003, app:membership-studies), so that the design ablations it summarises follow it; or
% directly after app:protocol / app:loss-details (before line 853) if Daniel prefers it next to tab:protocol.
% ---------------------------------------------------------------------------------------------------------------------
\section{Building with Weighted Circuit Masking}
\label{app:framework}
Weighted circuit masking fixes a problem, \labelcref{eq:p2}, and a way of relaxing it, but leaves most of the recipe open.
\Cref{tab:framework} lists its slots: what is reproduced, over which nodes, with which parametrisation, under which ablation methods and backgrounds, with which regularisers and optimiser, from which contexts, and how the result is scored.
The production choice in each slot was made by an experiment reported in these appendices, most of them on TuringLLM and carried over unchanged to the public models (\cref{tab:protocol}, Provenance); none is claimed to be the best setting, and several are known to trade one evaluation against another.

{\footnotesize
\setlength{\tabcolsep}{3pt}
\renewcommand{\arraystretch}{1.15}
\begin{longtable}{@{}p{0.13\linewidth}p{0.33\linewidth}p{0.33\linewidth}p{0.16\linewidth}@{}}
\caption{The slots of weighted circuit masking, the production choice in each, and the alternatives this paper examined or that the formulation admits.}
\label{tab:framework} \\
\toprule
\textbf{Slot} & \textbf{Production choice} & \textbf{Alternatives} & \textbf{Evidence} \\
\midrule
\endfirsthead
\toprule
\textbf{Slot} & \textbf{Production choice} & \textbf{Alternatives} & \textbf{Evidence} \\
\midrule
\endhead
\midrule
\multicolumn{4}{r@{}}{\emph{continued on the next page}} \\
\endfoot
\bottomrule
\endlastfoot
Target & the pre-activation $\preact$ of one SAE latent at its anchors; censored read under Z, C, A & an output metric, as in learned-mask circuits for outputs~\citep{bhaskar2024finding, caples2025scaling}; output- and latent-targeted circuits fail each other's objective (exploratory) & \cref{sec:ablmask}, \cref{app:secondary,app:logit-faith} \\
Nodes & every latent at every SAE site upstream of the target; SAE errors passed through; nodes recomputed & dense ReLU dictionaries (the method replicates; the random baseline needs its own design); members clamped (resample diagnostic) & \cref{sec:problem}, \cref{app:dense-relu,app:resample} \\
Parametrisation & gate $\sigma(\theta/T)$ and coefficient $\mathrm{softplus}(\psi)$ per node, $\alpha_0 = 1$ & $\alpha \equiv 1$ (unweighted circuit masking) & \cref{sec:weighted}, \cref{tab:run2026-alpha1} \\
Training ablations & Z (weight $1$), C and A ($0.25$ each) under evaluation's sparsity-preserving fill; resample D over the training contrast contexts ($1$) & one method alone; dense means; D in place of C; no D & \cref{app:single-method,app:respect-topk,app:resample} \\
Contrast background & $64$ close, silence-checked contexts per target ($48$ train, $16$ held out) & random, distant, or fused sources; the unchecked retrieval store & \cref{app:stores,app:contrast-source} \\
Regularisers & flat penalty $\lambda = 10^{-3}$ per open gate; closed-gate charge $\lambda(1-\sigma(\theta))|\alpha - 1|$; rank-keep $\gamma_{\mathrm{K}} = 10^{-1}$; one shared normaliser $\nu$ & other penalties $\lambda$ (sweep); off-target term, Top-K hinge, pull to rank $1$, firing margin; per-term normalisers & \cref{sec:lambda-transfer}, \cref{app:loss-details,app:rank-designs,app:membership-studies} \\
Optimiser & AdamW, $400$ steps, batch $4$; gate temperature annealed $1 \to 0.05$; learning-rate warm-up, then geometric decay & straight-through binarisation; $200$-step and reach-then-hold schedules; larger batches (exploratory) & \cref{tab:protocol}, \cref{app:membership-studies,app:framework-history} \\
Training contexts & $32$ strongest and $16$ mid-band activating contexts, deduplicated and split by whole groups of copies & a split without deduplication, which leaked copies into held-out sets & \cref{sec:setting}, \cref{app:heldout} \\
Evaluation & held-out strongest contexts; Z, A, C in $[0.75, 1.5]$ and necessity $\geq 0.9$; vacuous targets excluded and counted & reported, not required: mid-band contexts, sufficiency to induce, donor-background test, specificity, $\alpha = 1$ rescoring, random-circuit baselines, stable core & \cref{sec:eval}, \cref{app:protocol,app:nulls,app:underdetermination} \\
\end{longtable}}

\paragraph{Training in the evaluation's domain.}
Across these experiments, circuits did best on the evaluations their training resembled.
When a fit fails an evaluation, the failure is often a property of the recipe, and a slot of \cref{tab:framework} can put the fit in that evaluation's domain.
% app:resample (095: 72 targets, 0% -> 24%, CI +14 to +33; criterion 69% -> 62%, CI -17 to +3; 1.18x size, 1.45x time;
% D in place of C: 51%, CI -29 to -7), production paragraph (exp 102: 19 -> 39 and 48 -> 54 of 72, 1.33x size), T10
% (tab:run2026-tl-noR: +5.2, -0.6 to 11.6).
\emph{Real backgrounds.} Circuits trained only against zero and means failed a resample test almost always; adding a resample term over the training contrast contexts raised the resample pass rate on $72$ targets from $0\%$ to $24\%$, at $1.18\times$ the circuit size and $1.45\times$ the fitting time, and with rank-keep raised it took a $72$-target check panel from $19$ to $39$ targets passing (\cref{app:resample}).
Using it in place of the contrast-mean term gained as much on the resample test but cost the pass criterion, so terms should be added for a new evaluation rather than swapped for an old one; on the production sweep the term's effect on the pass criterion itself is $5.2$ points and not significant (\cref{tab:run2026-tl-noR}).
\emph{Ablation methods and fills.} A mask trained under one ablation method learns what that method rewards: without the zero-ablation term, coefficients overshoot zero ablation by orders of magnitude, and circuits trained against dense means can fail under the sparsity-preserving fill that evaluation uses (\cref{app:single-method,app:respect-topk}).
\emph{Top-K survival.} Reproducing the pre-activation does not keep the target in its site's Top-K; the censored read and rank-keep train for it, though neither removes amplifier circuits (\cref{app:loss-details,app:specificity}).
\emph{Other backgrounds.} Contrast faithfulness generalises furthest when training uses random or distant contrast contexts (\cref{app:contrast-source}).
\emph{Size.} Circuits shrink with $\lambda$ at a nearly constant pass rate up to a plateau's edge (\cref{sec:lambda-transfer}).
Two evaluations have no training term yet: sufficiency to induce, which no circuit is trained on (median $0.78$--$0.88$ on held-out contrast contexts, $0.35$ for TuringLLM attention outputs; \cref{app:production}), and transfer to weaker, mid-band contexts (\cref{sec:main-result}).
The framework's prescription for them is a term or a training set in their domain, which we have not tested.
% CANDIDATE EVIDENCE, NOT IN THE PAPER (main.tex line 503: the arena results are reported separately). If Daniel wants
% them here, each needs its own sentence and its own appendix or citation:
%  - arena A tuning (experiments/107-true-baselines/arena_a/TUNING.md): on GPT-2 transcoder + z-SAE substrate the
%    necessity hinge (NEC_W; off in production, docs/wcm-specification.md §2.3) was chosen at 0.3 by a pre-registered
%    9-target rule; NEC_W 1.0 raised nec_a (3.tc.3752 0.32 -> 0.95) but pushed A out of band (1.54). The censored read
%    is inert there by construction (k = d_sae).
%  - arena A scope a1 (z + z2 + tc) vs a2 (tc only) and the wikitext-smoke vs OpenWebText corpus test (DAN-208 text;
%    not traced to an output file in this pass).
%  - logit / maximise / behavioural objectives (arena C/D; memory notes say "don't quote yet").
A fit trained for an evaluation is, for that evaluation, scored on home turf (\cref{app:method-comparison}), so the evaluation stays meaningful only if its contexts and donors are never seen in training, as the donor-background test's donors are not (\cref{tab:protocol}), and the evaluations not trained for are reported alongside.

\paragraph{Choosing the sparsity penalty.}
% sec:lambda-transfer, tab:run2026-tl-sweep, tab:run2026-lambda-pilot.
The penalty sets circuit size and, past a point, the pass rate: on the TuringLLM sweep the pass rate holds at $89$--$93\%$ from $\lambda = 2.5 \times 10^{-5}$ to $5 \times 10^{-4}$ while each doubling of $\lambda$ shrinks circuits $1.4$--$1.9\times$, and it declines beyond (\cref{sec:lambda-transfer}).
We recommend fixing the selection rule before looking at pass rates, as our size-only pilot rule did on the public models, where it kept $\lambda = 10^{-3}$ on both (\cref{tab:run2026-lambda-pilot}), and reporting a curve over $\lambda$ rather than one point, since one setting that transfers is not the best setting for each model.

\paragraph{A new model or dictionary.}
Only the model and dictionary adapters change between our three models (\cref{sec:setting}); four things needed checking.
\emph{Precision}: TuringLLM runs in the bfloat16 its pipeline was validated in, but on the public models a bfloat16 data pass did not reproduce the float32 top-context sets, so they run in float32 (\cref{tab:protocol}); a precision check belongs before any fit, since node sets also move with rounding (\cref{app:underdetermination}).
\emph{The dictionary's sparsity}: $k$ sets the fill's budget and the censored read's cut; on a dictionary that keeps every latent ($k = d_{\mathrm{sae}}$) the censored read has no cut to model, and on dense ReLU dictionaries a random-circuit baseline with fitted coefficients stops discriminating, so necessity, not faithfulness, separates real from random circuits there (\cref{app:dense-relu}).
% The k = d_sae statement follows from the definition of tau^(K) (no K-th other latent); docs/wcm-specification.md §7.6.
\emph{Contrast retrieval}: contrast contexts need a pooled representation for every corpus sequence and a silence check, which rejects $14$--$51\%$ of the closest candidates on the targets we inspected (\cref{app:stores}).
\emph{Corpus copies}: duplicated windows silently leak into held-out sets unless pools are deduplicated and split by groups of copies (\cref{app:heldout}).

\paragraph{What to report.}
Because each slot changes the circuits, a result is interpretable only with its recipe: the ablation methods used in training and in evaluation; faithfulness under each method and necessity, on held-out contexts, by depth band and site kind, with intervals over targets; excluded and vacuous targets counted; circuit size with whether coefficients were fitted, and the $\alpha = 1$ rescoring; random-circuit baselines given the same fitting; and the reported-only evaluations (sufficiency to induce, the donor-background test, specificity, mid-band transfer) beside the pass rate (\cref{sec:discussion}).

\paragraph{Choices that matter, and what stays open.}
The ablation method changes not only scores but rankings: the two mean-ablation variants can order small node sets differently (\cref{app:respect-topk}), and in exploratory comparisons a method tended to win the scores of the ablation method it was trained or ranked under (\cref{app:method-comparison}), so comparisons between recipes need the full matrix of recipe against ablation method.
Node sets are not identified by any recipe: a refit with reordered training contexts shares a median Jaccard of $0.39$--$0.51$ with the original and its shared core alone passes for $44$--$72\%$ of targets (\cref{app:underdetermination}), so changes to a slot should be judged by sizes and scores, not by which latents moved.
Finally, the production choices were calibrated on TuringLLM targets that overlap its evaluation set (\cref{sec:setting}), and on the public models every setting except $\lambda$ carried over unchanged (\cref{tab:protocol}, Provenance), so on a new model they are defaults to be checked, not tuned values.

% ---------------------------------------------------------------------------------------------------------------------
% 2026-11-08 (Daniel, via coordinator): the development path as evidence for where WCM can go next. Every number is
% traced; sources in the comment before each step. Panels and protocols differ between steps, so numbers are compared
% only within a step (each step is a paired test against the recipe before it).
% ---------------------------------------------------------------------------------------------------------------------
\subsection{How the Production Recipe Was Reached}
\label{app:framework-history}
The production recipe is the end of a sequence of changes, each made because a run exposed a failure, and each tested against the recipe before it on the same targets, mostly under a selection rule stated before the comparison.
We list them in order as evidence of which kinds of change have paid off; the panels differ between steps, so numbers compare only within a step.

% app:stores (38% fallback), app:contrast-source; experiments/055-close-contrast, 071-contrast-sweep.
\paragraph{Contrast contexts built per target.}
The first contrast contexts came from a retrieval store whose index covered too few sequences, so $38\%$ of TuringLLM targets fell back to one shared list; WCM now selects contrast contexts per target at fit time, ranked by similarity to its strongest contexts and kept only if the target is silent (\cref{app:stores}), and a later comparison of contrast sources showed how much that choice shapes the C term (\cref{app:contrast-source}).
The lesson: inspect the data each training term actually sees.

% app:specificity, app:rank-designs; experiments/056-specificity.
\paragraph{Rank-keep against off-target lifts.}
Reproducing the pre-activation left some circuits restoring their target by lifting its whole neighbourhood; of the terms tried, rank-keep alone kept the gains without the alternatives' costs, and entered the first protocol at a small weight (\cref{app:rank-designs,app:specificity}).

% app:v1-run (15,046 targets, 46%, attention 14%, layers 8-11 22%); app:heldout (66.8%); app:resample (0.1% vs 43.1%,
% 747 targets); fig:failure-anatomy.
\paragraph{The first run at scale.}
Under that protocol, WCM was fitted to $15{,}046$ TuringLLM targets, and $46\%$ passed, $22\%$ at layers 8--11 and $14\%$ for attention outputs (\cref{app:v1-run}).
Analyses of the run exposed three problems: $66.8\%$ of held-out strongest contexts repeated a training context, which led to deduplicated records split by groups of copies (\cref{app:heldout}); with depth, failing circuits increasingly failed only the two mean ablations (\cref{fig:failure-anatomy}); and $0.1\%$ of circuits passed a resample test, against $43.1\%$ under the pass criterion (\cref{app:resample}).

% experiments/093-v2-validation/README.md (out/tables.md): 72 targets, 8 per band x kind, paired, pass floor 0.8, the
% v1 run's as-run held-out sets (leaky; same for both arms): 31 -> 50 of 72, +26.4 pp [15.3, 38.9]; L8-11 2 -> 11 of 24;
% attn 3 -> 12 of 24; median dZ 0.00; size ratio 0.91; fit time 1.27x. Pilot: experiments/090-sparse-fill-training.
% Also main.tex comment at app:respect-topk ("The same fill in training").
\paragraph{Training under the fill that evaluation uses.}
The mean-ablation failures were a mismatch: training replaced non-members with dense means, a code no Top-K dictionary emits, while evaluation re-imposed the Top-K (\cref{app:respect-topk}).
Training the A and C terms under evaluation's own fill and means raised the pass rate on $72$ stratified targets from $31$ to $50$ (paired difference $+26.4$ points, $95\%$ interval $15.3$--$38.9$), from $2$ to $11$ of $24$ at layers 8--11 and from $3$ to $12$ of $24$ for attention outputs, with zero-ablation faithfulness unchanged (median paired change $0.00$), circuits $0.91\times$ the size, and $1.27\times$ the fitting time.
The gain came from aligning training with evaluation, not from a new term.

% app:resample (095 numbers and the production paragraph: 102's 19 -> 39, 48 -> 54 of 72, 1.33x size);
% experiments/102-pass-improvement/README.md line 89 (diagnosis: 15 of 18 Z failures of the previous version restored
% the pre-activation and lost at the Top-K cut); tab:run2026-tl-noR.
\paragraph{A resample term, and a stronger rank-keep.}
To survive real backgrounds, training gained the resample term (above and \cref{app:resample}).
On the previous version's zero-ablation failures, $15$ of $18$ restored the target's pre-activation but lost it at the Top-K cut, so rank-keep was raised from $3 \times 10^{-3}$ to $10^{-1}$ with the resample term at full weight; chosen by a rule stated before the comparison, the pair took a $72$-target check panel from $48$ to $54$ targets passing and from $19$ to $39$ passing the resample test, at $1.33\times$ the circuit size (\cref{app:resample}).
In the production sweep, the arm trained without the resample term passes $5.2$ points less, not significantly (\cref{tab:run2026-tl-noR}).

% experiments/107-true-baselines/latent_pilot/README.md (9 targets enriched for failures: 3 -> 5 of 9, size ratio ~1.04;
% plain geometric: circuits 2-4x larger, sum(lr) 6.3 vs 20; constant lr makes lr/T grow 20x) and
% experiments/107-true-baselines/adaptive_lambda/README.md, "Check" table (arm "geo x 3.15 + censored term
% (latent_pilot)": 63 of 72 vs 54 for the previous version, +10 / -1, size ratio 1.04 on pass-in-both). VERIFY: whether
% that check arm included the 0.1 warm-up (run_tonight.sh stage 4 fits check_geo_x3.15_warmup_censor).
\paragraph{A learning-rate schedule and the censored read.}
Under a constant learning rate, the gate step grows twentyfold as the gate temperature anneals, and some fits degraded late in training.
A geometric decay alone stabilised them but left gates too little total step to close, so circuits grew $2$--$4\times$; raising its peak $3.15\times$ restored the summed step size, and a short warm-up was added for the early spikes the higher peak brought back (\cref{app:loss-details}).
With the censored read added, so that training also sees whether the target survives its site's Top-K, the check panel went from $54$ to $63$ of $72$ targets passing ($10$ gained, $1$ lost) at about the same circuit size (median ratio $1.04$).

% sec:lambda-transfer, tab:run2026-lambda-pilot; experiments/108-batch-size/out/analysis.md (16 TuringLLM targets;
% B16 at 400 steps: 15 / 16 vs 16 / 16, size 0.75x on pass-in-both, 99.9 vs 29.1 fit-minutes = 3.4x; pre-stated rule:
% +-1 target is noise); experiments/107-true-baselines/adaptive_lambda/README.md (every shrinking policy traded passes
% for size on the check panel; losses were necessity failures).
\paragraph{Settings left as they were.}
The penalty was carried to the public models by a size-only rule fixed before a pilot (\cref{sec:lambda-transfer}).
In exploratory tests, batches of $16$ for the same $400$ steps passed $15$ of $16$ TuringLLM targets against $16$, with circuits $0.75\times$ the size where both passed, at $3.4\times$ the fitting time, so the batch of $4$ was kept; and policies that adapted $\lambda$ per target traded passes for size on the check panel, losing passes mostly to necessity, which no training term covers.
Other rejected designs are in \cref{app:membership-studies}.

% Directions: app:production (induce), sec:main-result (mid-band, deep band, attention), app:specificity (amplifiers,
% v1: 30% flagged), app:rank-designs (off-target term: lifts cut 3-7x, but grew circuits), app:resample (anchors
% only; overfiring away from the anchor), adaptive_lambda README (necessity).
\paragraph{Directions not yet taken.}
The same path suggests where the recipe can go next.
Necessity, sufficiency to induce, and transfer to mid-band contexts are scored but not trained for; each could get a term or training contexts in its domain, as the resample test did.
Deep and attention-output targets still pass least (\cref{sec:main-result}), and in the schedule pilot the attention failures kept zero-ablation faithfulness near $0$ under every schedule, so the learning rate does not explain them.
Amplifier circuits survive rank-keep (\cref{app:specificity}); the off-target term cut lifts three- to sevenfold but grew circuits (\cref{app:rank-designs}), so a cheaper specificity term remains open.
And circuits explain the target at its anchors, not at every position (\cref{app:resample}); training on more positions is untested.
\fi\fi
        % ~1/3 page, Discussion (replaces the paragraph
%                                                                    % "Ablation methods are part of the method")
%   \def\frameworkpart{appendix}% =====================================================================================================================
% DAN-208 (drafted 2026-11-08, agent "framework"): WCM as a framework for solving for circuits, not a fixed algorithm.
% Two versions in one file. Choose with \def\frameworkpart{main} or \def\frameworkpart{appendix} before \input:
%
%   \def\frameworkpart{main}% =====================================================================================================================
% DAN-208 (drafted 2026-11-08, agent "framework"): WCM as a framework for solving for circuits, not a fixed algorithm.
% Two versions in one file. Choose with \def\frameworkpart{main} or \def\frameworkpart{appendix} before \input:
%
%   \def\frameworkpart{main}\input{drafts/framework_section}        % ~1/3 page, Discussion (replaces the paragraph
%                                                                    % "Ablation methods are part of the method")
%   \def\frameworkpart{appendix}\input{drafts/framework_section}    % ~1.3 pages, its own appendix section
%
% New labels: sec:framework (main), app:framework, tab:framework (appendix). Every other \cref target exists in
% main.tex as of 2026-11-08. Every number below is already printed in main.tex (source comments give the place);
% the one exception, experiment 108, is marked NOT-YET-IN-PAPER where it appears.
% Vocabulary: paper/notes/vocabulary.md (no "endpoint", "floor", "drivers"; "fill" only as main.tex already uses it).
% Never-claim list respected: no "first"; nothing here claims novelty.
% =====================================================================================================================
\providecommand{\frameworkpart}{main}
\def\frameworkmain{main}
\def\frameworkappendix{appendix}
\ifx\frameworkpart\frameworkmain
% ---------------------------------------------------------------------------------------------------------------------
% (1) MAIN-TEXT VERSION. Proposed place: sec:discussion, in place of the paragraph "Ablation methods are part of the
% method" (main.tex lines 548-552), whose evidence it repeats and extends; the paragraph "Recommendations" (lines
% 564-566) can then point here. If it must sit in Method instead, it fits after "Pass criterion" (line 341), before
% \section{Experiments and Results}, but there it would cite results not yet shown.
% Numbers: 2-14% = abstract / run2026_alpha1 (sec:weighted); "one to three orders of magnitude" = app:single-method
% (exploratory); contrast sources = app:contrast-source (protocol v1); resample = app:resample.
% ---------------------------------------------------------------------------------------------------------------------
\subsection{Using WCM as a Framework}
\label{sec:framework}
WCM is better read as a framework for solving for circuits than as one fixed algorithm.
Its parts are interchangeable: an objective (the quantity reproduced, the ablation methods it is reproduced under, the regularisers, and their weights; \cref{eq:loss}), a parametrisation (a gate and a coefficient per node), data (training, contrast, and held-out contexts), and an optimiser.
The production recipe is one setting of these parts, each chosen by an experiment (\cref{tab:framework}), and changing a part changes the circuits found.
With every coefficient fixed at $1$ the same nodes pass for $2$--$14\%$ of targets (\cref{sec:weighted}); in exploratory fits, coefficients trained without zero ablation overshoot it by one to three orders of magnitude (\cref{app:single-method}); circuits trained against close contrast contexts are the most faithful on close backgrounds and the least on any other (\cref{app:contrast-source}); and circuits failed a resample test until training included a resample term (\cref{app:resample}).
When WCM fails an evaluation, we therefore suggest first asking whether training ever saw that evaluation's domain, and if not, adding the term, contexts, or setting that puts it there, while keeping the evaluation's own contexts and donors out of training so that it stays a held-out test (\cref{app:framework}).
A pass rate then describes a recipe, a dictionary, and an evaluation together, not WCM alone, which is why the ablation methods used for training and evaluation belong in the method statement~\citep{mueller2025mib, mahale2026explanation}.
%
\else\ifx\frameworkpart\frameworkappendix
% ---------------------------------------------------------------------------------------------------------------------
% (2) APPENDIX VERSION. Proposed place: a new appendix section directly before "Weighted Circuit Masking Design
% Ablations" (main.tex line 2003, app:membership-studies), so that the design ablations it summarises follow it; or
% directly after app:protocol / app:loss-details (before line 853) if Daniel prefers it next to tab:protocol.
% ---------------------------------------------------------------------------------------------------------------------
\section{Building with Weighted Circuit Masking}
\label{app:framework}
Weighted circuit masking fixes a problem, \labelcref{eq:p2}, and a way of relaxing it, but leaves most of the recipe open.
\Cref{tab:framework} lists its slots: what is reproduced, over which nodes, with which parametrisation, under which ablation methods and backgrounds, with which regularisers and optimiser, from which contexts, and how the result is scored.
The production choice in each slot was made by an experiment reported in these appendices, most of them on TuringLLM and carried over unchanged to the public models (\cref{tab:protocol}, Provenance); none is claimed to be the best setting, and several are known to trade one evaluation against another.

{\footnotesize
\setlength{\tabcolsep}{3pt}
\renewcommand{\arraystretch}{1.15}
\begin{longtable}{@{}p{0.13\linewidth}p{0.33\linewidth}p{0.33\linewidth}p{0.16\linewidth}@{}}
\caption{The slots of weighted circuit masking, the production choice in each, and the alternatives this paper examined or that the formulation admits.}
\label{tab:framework} \\
\toprule
\textbf{Slot} & \textbf{Production choice} & \textbf{Alternatives} & \textbf{Evidence} \\
\midrule
\endfirsthead
\toprule
\textbf{Slot} & \textbf{Production choice} & \textbf{Alternatives} & \textbf{Evidence} \\
\midrule
\endhead
\midrule
\multicolumn{4}{r@{}}{\emph{continued on the next page}} \\
\endfoot
\bottomrule
\endlastfoot
Target & the pre-activation $\preact$ of one SAE latent at its anchors; censored read under Z, C, A & an output metric, as in learned-mask circuits for outputs~\citep{bhaskar2024finding, caples2025scaling}; output- and latent-targeted circuits fail each other's objective (exploratory) & \cref{sec:ablmask}, \cref{app:secondary,app:logit-faith} \\
Nodes & every latent at every SAE site upstream of the target; SAE errors passed through; nodes recomputed & dense ReLU dictionaries (the method replicates; the random baseline needs its own design); members clamped (resample diagnostic) & \cref{sec:problem}, \cref{app:dense-relu,app:resample} \\
Parametrisation & gate $\sigma(\theta/T)$ and coefficient $\mathrm{softplus}(\psi)$ per node, $\alpha_0 = 1$ & $\alpha \equiv 1$ (unweighted circuit masking) & \cref{sec:weighted}, \cref{tab:run2026-alpha1} \\
Training ablations & Z (weight $1$), C and A ($0.25$ each) under evaluation's sparsity-preserving fill; resample D over the training contrast contexts ($1$) & one method alone; dense means; D in place of C; no D & \cref{app:single-method,app:respect-topk,app:resample} \\
Contrast background & $64$ close, silence-checked contexts per target ($48$ train, $16$ held out) & random, distant, or fused sources; the unchecked retrieval store & \cref{app:stores,app:contrast-source} \\
Regularisers & flat penalty $\lambda = 10^{-3}$ per open gate; closed-gate charge $\lambda(1-\sigma(\theta))|\alpha - 1|$; rank-keep $\gamma_{\mathrm{K}} = 10^{-1}$; one shared normaliser $\nu$ & other penalties $\lambda$ (sweep); off-target term, Top-K hinge, pull to rank $1$, firing margin; per-term normalisers & \cref{sec:lambda-transfer}, \cref{app:loss-details,app:rank-designs,app:membership-studies} \\
Optimiser & AdamW, $400$ steps, batch $4$; gate temperature annealed $1 \to 0.05$; learning-rate warm-up, then geometric decay & straight-through binarisation; $200$-step and reach-then-hold schedules; larger batches (exploratory) & \cref{tab:protocol}, \cref{app:membership-studies,app:framework-history} \\
Training contexts & $32$ strongest and $16$ mid-band activating contexts, deduplicated and split by whole groups of copies & a split without deduplication, which leaked copies into held-out sets & \cref{sec:setting}, \cref{app:heldout} \\
Evaluation & held-out strongest contexts; Z, A, C in $[0.75, 1.5]$ and necessity $\geq 0.9$; vacuous targets excluded and counted & reported, not required: mid-band contexts, sufficiency to induce, donor-background test, specificity, $\alpha = 1$ rescoring, random-circuit baselines, stable core & \cref{sec:eval}, \cref{app:protocol,app:nulls,app:underdetermination} \\
\end{longtable}}

\paragraph{Training in the evaluation's domain.}
Across these experiments, circuits did best on the evaluations their training resembled.
When a fit fails an evaluation, the failure is often a property of the recipe, and a slot of \cref{tab:framework} can put the fit in that evaluation's domain.
% app:resample (095: 72 targets, 0% -> 24%, CI +14 to +33; criterion 69% -> 62%, CI -17 to +3; 1.18x size, 1.45x time;
% D in place of C: 51%, CI -29 to -7), production paragraph (exp 102: 19 -> 39 and 48 -> 54 of 72, 1.33x size), T10
% (tab:run2026-tl-noR: +5.2, -0.6 to 11.6).
\emph{Real backgrounds.} Circuits trained only against zero and means failed a resample test almost always; adding a resample term over the training contrast contexts raised the resample pass rate on $72$ targets from $0\%$ to $24\%$, at $1.18\times$ the circuit size and $1.45\times$ the fitting time, and with rank-keep raised it took a $72$-target check panel from $19$ to $39$ targets passing (\cref{app:resample}).
Using it in place of the contrast-mean term gained as much on the resample test but cost the pass criterion, so terms should be added for a new evaluation rather than swapped for an old one; on the production sweep the term's effect on the pass criterion itself is $5.2$ points and not significant (\cref{tab:run2026-tl-noR}).
\emph{Ablation methods and fills.} A mask trained under one ablation method learns what that method rewards: without the zero-ablation term, coefficients overshoot zero ablation by orders of magnitude, and circuits trained against dense means can fail under the sparsity-preserving fill that evaluation uses (\cref{app:single-method,app:respect-topk}).
\emph{Top-K survival.} Reproducing the pre-activation does not keep the target in its site's Top-K; the censored read and rank-keep train for it, though neither removes amplifier circuits (\cref{app:loss-details,app:specificity}).
\emph{Other backgrounds.} Contrast faithfulness generalises furthest when training uses random or distant contrast contexts (\cref{app:contrast-source}).
\emph{Size.} Circuits shrink with $\lambda$ at a nearly constant pass rate up to a plateau's edge (\cref{sec:lambda-transfer}).
Two evaluations have no training term yet: sufficiency to induce, which no circuit is trained on (median $0.78$--$0.88$ on held-out contrast contexts, $0.35$ for TuringLLM attention outputs; \cref{app:production}), and transfer to weaker, mid-band contexts (\cref{sec:main-result}).
The framework's prescription for them is a term or a training set in their domain, which we have not tested.
% CANDIDATE EVIDENCE, NOT IN THE PAPER (main.tex line 503: the arena results are reported separately). If Daniel wants
% them here, each needs its own sentence and its own appendix or citation:
%  - arena A tuning (experiments/107-true-baselines/arena_a/TUNING.md): on GPT-2 transcoder + z-SAE substrate the
%    necessity hinge (NEC_W; off in production, docs/wcm-specification.md §2.3) was chosen at 0.3 by a pre-registered
%    9-target rule; NEC_W 1.0 raised nec_a (3.tc.3752 0.32 -> 0.95) but pushed A out of band (1.54). The censored read
%    is inert there by construction (k = d_sae).
%  - arena A scope a1 (z + z2 + tc) vs a2 (tc only) and the wikitext-smoke vs OpenWebText corpus test (DAN-208 text;
%    not traced to an output file in this pass).
%  - logit / maximise / behavioural objectives (arena C/D; memory notes say "don't quote yet").
A fit trained for an evaluation is, for that evaluation, scored on home turf (\cref{app:method-comparison}), so the evaluation stays meaningful only if its contexts and donors are never seen in training, as the donor-background test's donors are not (\cref{tab:protocol}), and the evaluations not trained for are reported alongside.

\paragraph{Choosing the sparsity penalty.}
% sec:lambda-transfer, tab:run2026-tl-sweep, tab:run2026-lambda-pilot.
The penalty sets circuit size and, past a point, the pass rate: on the TuringLLM sweep the pass rate holds at $89$--$93\%$ from $\lambda = 2.5 \times 10^{-5}$ to $5 \times 10^{-4}$ while each doubling of $\lambda$ shrinks circuits $1.4$--$1.9\times$, and it declines beyond (\cref{sec:lambda-transfer}).
We recommend fixing the selection rule before looking at pass rates, as our size-only pilot rule did on the public models, where it kept $\lambda = 10^{-3}$ on both (\cref{tab:run2026-lambda-pilot}), and reporting a curve over $\lambda$ rather than one point, since one setting that transfers is not the best setting for each model.

\paragraph{A new model or dictionary.}
Only the model and dictionary adapters change between our three models (\cref{sec:setting}); four things needed checking.
\emph{Precision}: TuringLLM runs in the bfloat16 its pipeline was validated in, but on the public models a bfloat16 data pass did not reproduce the float32 top-context sets, so they run in float32 (\cref{tab:protocol}); a precision check belongs before any fit, since node sets also move with rounding (\cref{app:underdetermination}).
\emph{The dictionary's sparsity}: $k$ sets the fill's budget and the censored read's cut; on a dictionary that keeps every latent ($k = d_{\mathrm{sae}}$) the censored read has no cut to model, and on dense ReLU dictionaries a random-circuit baseline with fitted coefficients stops discriminating, so necessity, not faithfulness, separates real from random circuits there (\cref{app:dense-relu}).
% The k = d_sae statement follows from the definition of tau^(K) (no K-th other latent); docs/wcm-specification.md §7.6.
\emph{Contrast retrieval}: contrast contexts need a pooled representation for every corpus sequence and a silence check, which rejects $14$--$51\%$ of the closest candidates on the targets we inspected (\cref{app:stores}).
\emph{Corpus copies}: duplicated windows silently leak into held-out sets unless pools are deduplicated and split by groups of copies (\cref{app:heldout}).

\paragraph{What to report.}
Because each slot changes the circuits, a result is interpretable only with its recipe: the ablation methods used in training and in evaluation; faithfulness under each method and necessity, on held-out contexts, by depth band and site kind, with intervals over targets; excluded and vacuous targets counted; circuit size with whether coefficients were fitted, and the $\alpha = 1$ rescoring; random-circuit baselines given the same fitting; and the reported-only evaluations (sufficiency to induce, the donor-background test, specificity, mid-band transfer) beside the pass rate (\cref{sec:discussion}).

\paragraph{Choices that matter, and what stays open.}
The ablation method changes not only scores but rankings: the two mean-ablation variants can order small node sets differently (\cref{app:respect-topk}), and in exploratory comparisons a method tended to win the scores of the ablation method it was trained or ranked under (\cref{app:method-comparison}), so comparisons between recipes need the full matrix of recipe against ablation method.
Node sets are not identified by any recipe: a refit with reordered training contexts shares a median Jaccard of $0.39$--$0.51$ with the original and its shared core alone passes for $44$--$72\%$ of targets (\cref{app:underdetermination}), so changes to a slot should be judged by sizes and scores, not by which latents moved.
Finally, the production choices were calibrated on TuringLLM targets that overlap its evaluation set (\cref{sec:setting}), and on the public models every setting except $\lambda$ carried over unchanged (\cref{tab:protocol}, Provenance), so on a new model they are defaults to be checked, not tuned values.

% ---------------------------------------------------------------------------------------------------------------------
% 2026-11-08 (Daniel, via coordinator): the development path as evidence for where WCM can go next. Every number is
% traced; sources in the comment before each step. Panels and protocols differ between steps, so numbers are compared
% only within a step (each step is a paired test against the recipe before it).
% ---------------------------------------------------------------------------------------------------------------------
\subsection{How the Production Recipe Was Reached}
\label{app:framework-history}
The production recipe is the end of a sequence of changes, each made because a run exposed a failure, and each tested against the recipe before it on the same targets, mostly under a selection rule stated before the comparison.
We list them in order as evidence of which kinds of change have paid off; the panels differ between steps, so numbers compare only within a step.

% app:stores (38% fallback), app:contrast-source; experiments/055-close-contrast, 071-contrast-sweep.
\paragraph{Contrast contexts built per target.}
The first contrast contexts came from a retrieval store whose index covered too few sequences, so $38\%$ of TuringLLM targets fell back to one shared list; WCM now selects contrast contexts per target at fit time, ranked by similarity to its strongest contexts and kept only if the target is silent (\cref{app:stores}), and a later comparison of contrast sources showed how much that choice shapes the C term (\cref{app:contrast-source}).
The lesson: inspect the data each training term actually sees.

% app:specificity, app:rank-designs; experiments/056-specificity.
\paragraph{Rank-keep against off-target lifts.}
Reproducing the pre-activation left some circuits restoring their target by lifting its whole neighbourhood; of the terms tried, rank-keep alone kept the gains without the alternatives' costs, and entered the first protocol at a small weight (\cref{app:rank-designs,app:specificity}).

% app:v1-run (15,046 targets, 46%, attention 14%, layers 8-11 22%); app:heldout (66.8%); app:resample (0.1% vs 43.1%,
% 747 targets); fig:failure-anatomy.
\paragraph{The first run at scale.}
Under that protocol, WCM was fitted to $15{,}046$ TuringLLM targets, and $46\%$ passed, $22\%$ at layers 8--11 and $14\%$ for attention outputs (\cref{app:v1-run}).
Analyses of the run exposed three problems: $66.8\%$ of held-out strongest contexts repeated a training context, which led to deduplicated records split by groups of copies (\cref{app:heldout}); with depth, failing circuits increasingly failed only the two mean ablations (\cref{fig:failure-anatomy}); and $0.1\%$ of circuits passed a resample test, against $43.1\%$ under the pass criterion (\cref{app:resample}).

% experiments/093-v2-validation/README.md (out/tables.md): 72 targets, 8 per band x kind, paired, pass floor 0.8, the
% v1 run's as-run held-out sets (leaky; same for both arms): 31 -> 50 of 72, +26.4 pp [15.3, 38.9]; L8-11 2 -> 11 of 24;
% attn 3 -> 12 of 24; median dZ 0.00; size ratio 0.91; fit time 1.27x. Pilot: experiments/090-sparse-fill-training.
% Also main.tex comment at app:respect-topk ("The same fill in training").
\paragraph{Training under the fill that evaluation uses.}
The mean-ablation failures were a mismatch: training replaced non-members with dense means, a code no Top-K dictionary emits, while evaluation re-imposed the Top-K (\cref{app:respect-topk}).
Training the A and C terms under evaluation's own fill and means raised the pass rate on $72$ stratified targets from $31$ to $50$ (paired difference $+26.4$ points, $95\%$ interval $15.3$--$38.9$), from $2$ to $11$ of $24$ at layers 8--11 and from $3$ to $12$ of $24$ for attention outputs, with zero-ablation faithfulness unchanged (median paired change $0.00$), circuits $0.91\times$ the size, and $1.27\times$ the fitting time.
The gain came from aligning training with evaluation, not from a new term.

% app:resample (095 numbers and the production paragraph: 102's 19 -> 39, 48 -> 54 of 72, 1.33x size);
% experiments/102-pass-improvement/README.md line 89 (diagnosis: 15 of 18 Z failures of the previous version restored
% the pre-activation and lost at the Top-K cut); tab:run2026-tl-noR.
\paragraph{A resample term, and a stronger rank-keep.}
To survive real backgrounds, training gained the resample term (above and \cref{app:resample}).
On the previous version's zero-ablation failures, $15$ of $18$ restored the target's pre-activation but lost it at the Top-K cut, so rank-keep was raised from $3 \times 10^{-3}$ to $10^{-1}$ with the resample term at full weight; chosen by a rule stated before the comparison, the pair took a $72$-target check panel from $48$ to $54$ targets passing and from $19$ to $39$ passing the resample test, at $1.33\times$ the circuit size (\cref{app:resample}).
In the production sweep, the arm trained without the resample term passes $5.2$ points less, not significantly (\cref{tab:run2026-tl-noR}).

% experiments/107-true-baselines/latent_pilot/README.md (9 targets enriched for failures: 3 -> 5 of 9, size ratio ~1.04;
% plain geometric: circuits 2-4x larger, sum(lr) 6.3 vs 20; constant lr makes lr/T grow 20x) and
% experiments/107-true-baselines/adaptive_lambda/README.md, "Check" table (arm "geo x 3.15 + censored term
% (latent_pilot)": 63 of 72 vs 54 for the previous version, +10 / -1, size ratio 1.04 on pass-in-both). VERIFY: whether
% that check arm included the 0.1 warm-up (run_tonight.sh stage 4 fits check_geo_x3.15_warmup_censor).
\paragraph{A learning-rate schedule and the censored read.}
Under a constant learning rate, the gate step grows twentyfold as the gate temperature anneals, and some fits degraded late in training.
A geometric decay alone stabilised them but left gates too little total step to close, so circuits grew $2$--$4\times$; raising its peak $3.15\times$ restored the summed step size, and a short warm-up was added for the early spikes the higher peak brought back (\cref{app:loss-details}).
With the censored read added, so that training also sees whether the target survives its site's Top-K, the check panel went from $54$ to $63$ of $72$ targets passing ($10$ gained, $1$ lost) at about the same circuit size (median ratio $1.04$).

% sec:lambda-transfer, tab:run2026-lambda-pilot; experiments/108-batch-size/out/analysis.md (16 TuringLLM targets;
% B16 at 400 steps: 15 / 16 vs 16 / 16, size 0.75x on pass-in-both, 99.9 vs 29.1 fit-minutes = 3.4x; pre-stated rule:
% +-1 target is noise); experiments/107-true-baselines/adaptive_lambda/README.md (every shrinking policy traded passes
% for size on the check panel; losses were necessity failures).
\paragraph{Settings left as they were.}
The penalty was carried to the public models by a size-only rule fixed before a pilot (\cref{sec:lambda-transfer}).
In exploratory tests, batches of $16$ for the same $400$ steps passed $15$ of $16$ TuringLLM targets against $16$, with circuits $0.75\times$ the size where both passed, at $3.4\times$ the fitting time, so the batch of $4$ was kept; and policies that adapted $\lambda$ per target traded passes for size on the check panel, losing passes mostly to necessity, which no training term covers.
Other rejected designs are in \cref{app:membership-studies}.

% Directions: app:production (induce), sec:main-result (mid-band, deep band, attention), app:specificity (amplifiers,
% v1: 30% flagged), app:rank-designs (off-target term: lifts cut 3-7x, but grew circuits), app:resample (anchors
% only; overfiring away from the anchor), adaptive_lambda README (necessity).
\paragraph{Directions not yet taken.}
The same path suggests where the recipe can go next.
Necessity, sufficiency to induce, and transfer to mid-band contexts are scored but not trained for; each could get a term or training contexts in its domain, as the resample test did.
Deep and attention-output targets still pass least (\cref{sec:main-result}), and in the schedule pilot the attention failures kept zero-ablation faithfulness near $0$ under every schedule, so the learning rate does not explain them.
Amplifier circuits survive rank-keep (\cref{app:specificity}); the off-target term cut lifts three- to sevenfold but grew circuits (\cref{app:rank-designs}), so a cheaper specificity term remains open.
And circuits explain the target at its anchors, not at every position (\cref{app:resample}); training on more positions is untested.
\fi\fi
        % ~1/3 page, Discussion (replaces the paragraph
%                                                                    % "Ablation methods are part of the method")
%   \def\frameworkpart{appendix}% =====================================================================================================================
% DAN-208 (drafted 2026-11-08, agent "framework"): WCM as a framework for solving for circuits, not a fixed algorithm.
% Two versions in one file. Choose with \def\frameworkpart{main} or \def\frameworkpart{appendix} before \input:
%
%   \def\frameworkpart{main}\input{drafts/framework_section}        % ~1/3 page, Discussion (replaces the paragraph
%                                                                    % "Ablation methods are part of the method")
%   \def\frameworkpart{appendix}\input{drafts/framework_section}    % ~1.3 pages, its own appendix section
%
% New labels: sec:framework (main), app:framework, tab:framework (appendix). Every other \cref target exists in
% main.tex as of 2026-11-08. Every number below is already printed in main.tex (source comments give the place);
% the one exception, experiment 108, is marked NOT-YET-IN-PAPER where it appears.
% Vocabulary: paper/notes/vocabulary.md (no "endpoint", "floor", "drivers"; "fill" only as main.tex already uses it).
% Never-claim list respected: no "first"; nothing here claims novelty.
% =====================================================================================================================
\providecommand{\frameworkpart}{main}
\def\frameworkmain{main}
\def\frameworkappendix{appendix}
\ifx\frameworkpart\frameworkmain
% ---------------------------------------------------------------------------------------------------------------------
% (1) MAIN-TEXT VERSION. Proposed place: sec:discussion, in place of the paragraph "Ablation methods are part of the
% method" (main.tex lines 548-552), whose evidence it repeats and extends; the paragraph "Recommendations" (lines
% 564-566) can then point here. If it must sit in Method instead, it fits after "Pass criterion" (line 341), before
% \section{Experiments and Results}, but there it would cite results not yet shown.
% Numbers: 2-14% = abstract / run2026_alpha1 (sec:weighted); "one to three orders of magnitude" = app:single-method
% (exploratory); contrast sources = app:contrast-source (protocol v1); resample = app:resample.
% ---------------------------------------------------------------------------------------------------------------------
\subsection{Using WCM as a Framework}
\label{sec:framework}
WCM is better read as a framework for solving for circuits than as one fixed algorithm.
Its parts are interchangeable: an objective (the quantity reproduced, the ablation methods it is reproduced under, the regularisers, and their weights; \cref{eq:loss}), a parametrisation (a gate and a coefficient per node), data (training, contrast, and held-out contexts), and an optimiser.
The production recipe is one setting of these parts, each chosen by an experiment (\cref{tab:framework}), and changing a part changes the circuits found.
With every coefficient fixed at $1$ the same nodes pass for $2$--$14\%$ of targets (\cref{sec:weighted}); in exploratory fits, coefficients trained without zero ablation overshoot it by one to three orders of magnitude (\cref{app:single-method}); circuits trained against close contrast contexts are the most faithful on close backgrounds and the least on any other (\cref{app:contrast-source}); and circuits failed a resample test until training included a resample term (\cref{app:resample}).
When WCM fails an evaluation, we therefore suggest first asking whether training ever saw that evaluation's domain, and if not, adding the term, contexts, or setting that puts it there, while keeping the evaluation's own contexts and donors out of training so that it stays a held-out test (\cref{app:framework}).
A pass rate then describes a recipe, a dictionary, and an evaluation together, not WCM alone, which is why the ablation methods used for training and evaluation belong in the method statement~\citep{mueller2025mib, mahale2026explanation}.
%
\else\ifx\frameworkpart\frameworkappendix
% ---------------------------------------------------------------------------------------------------------------------
% (2) APPENDIX VERSION. Proposed place: a new appendix section directly before "Weighted Circuit Masking Design
% Ablations" (main.tex line 2003, app:membership-studies), so that the design ablations it summarises follow it; or
% directly after app:protocol / app:loss-details (before line 853) if Daniel prefers it next to tab:protocol.
% ---------------------------------------------------------------------------------------------------------------------
\section{Building with Weighted Circuit Masking}
\label{app:framework}
Weighted circuit masking fixes a problem, \labelcref{eq:p2}, and a way of relaxing it, but leaves most of the recipe open.
\Cref{tab:framework} lists its slots: what is reproduced, over which nodes, with which parametrisation, under which ablation methods and backgrounds, with which regularisers and optimiser, from which contexts, and how the result is scored.
The production choice in each slot was made by an experiment reported in these appendices, most of them on TuringLLM and carried over unchanged to the public models (\cref{tab:protocol}, Provenance); none is claimed to be the best setting, and several are known to trade one evaluation against another.

{\footnotesize
\setlength{\tabcolsep}{3pt}
\renewcommand{\arraystretch}{1.15}
\begin{longtable}{@{}p{0.13\linewidth}p{0.33\linewidth}p{0.33\linewidth}p{0.16\linewidth}@{}}
\caption{The slots of weighted circuit masking, the production choice in each, and the alternatives this paper examined or that the formulation admits.}
\label{tab:framework} \\
\toprule
\textbf{Slot} & \textbf{Production choice} & \textbf{Alternatives} & \textbf{Evidence} \\
\midrule
\endfirsthead
\toprule
\textbf{Slot} & \textbf{Production choice} & \textbf{Alternatives} & \textbf{Evidence} \\
\midrule
\endhead
\midrule
\multicolumn{4}{r@{}}{\emph{continued on the next page}} \\
\endfoot
\bottomrule
\endlastfoot
Target & the pre-activation $\preact$ of one SAE latent at its anchors; censored read under Z, C, A & an output metric, as in learned-mask circuits for outputs~\citep{bhaskar2024finding, caples2025scaling}; output- and latent-targeted circuits fail each other's objective (exploratory) & \cref{sec:ablmask}, \cref{app:secondary,app:logit-faith} \\
Nodes & every latent at every SAE site upstream of the target; SAE errors passed through; nodes recomputed & dense ReLU dictionaries (the method replicates; the random baseline needs its own design); members clamped (resample diagnostic) & \cref{sec:problem}, \cref{app:dense-relu,app:resample} \\
Parametrisation & gate $\sigma(\theta/T)$ and coefficient $\mathrm{softplus}(\psi)$ per node, $\alpha_0 = 1$ & $\alpha \equiv 1$ (unweighted circuit masking) & \cref{sec:weighted}, \cref{tab:run2026-alpha1} \\
Training ablations & Z (weight $1$), C and A ($0.25$ each) under evaluation's sparsity-preserving fill; resample D over the training contrast contexts ($1$) & one method alone; dense means; D in place of C; no D & \cref{app:single-method,app:respect-topk,app:resample} \\
Contrast background & $64$ close, silence-checked contexts per target ($48$ train, $16$ held out) & random, distant, or fused sources; the unchecked retrieval store & \cref{app:stores,app:contrast-source} \\
Regularisers & flat penalty $\lambda = 10^{-3}$ per open gate; closed-gate charge $\lambda(1-\sigma(\theta))|\alpha - 1|$; rank-keep $\gamma_{\mathrm{K}} = 10^{-1}$; one shared normaliser $\nu$ & other penalties $\lambda$ (sweep); off-target term, Top-K hinge, pull to rank $1$, firing margin; per-term normalisers & \cref{sec:lambda-transfer}, \cref{app:loss-details,app:rank-designs,app:membership-studies} \\
Optimiser & AdamW, $400$ steps, batch $4$; gate temperature annealed $1 \to 0.05$; learning-rate warm-up, then geometric decay & straight-through binarisation; $200$-step and reach-then-hold schedules; larger batches (exploratory) & \cref{tab:protocol}, \cref{app:membership-studies,app:framework-history} \\
Training contexts & $32$ strongest and $16$ mid-band activating contexts, deduplicated and split by whole groups of copies & a split without deduplication, which leaked copies into held-out sets & \cref{sec:setting}, \cref{app:heldout} \\
Evaluation & held-out strongest contexts; Z, A, C in $[0.75, 1.5]$ and necessity $\geq 0.9$; vacuous targets excluded and counted & reported, not required: mid-band contexts, sufficiency to induce, donor-background test, specificity, $\alpha = 1$ rescoring, random-circuit baselines, stable core & \cref{sec:eval}, \cref{app:protocol,app:nulls,app:underdetermination} \\
\end{longtable}}

\paragraph{Training in the evaluation's domain.}
Across these experiments, circuits did best on the evaluations their training resembled.
When a fit fails an evaluation, the failure is often a property of the recipe, and a slot of \cref{tab:framework} can put the fit in that evaluation's domain.
% app:resample (095: 72 targets, 0% -> 24%, CI +14 to +33; criterion 69% -> 62%, CI -17 to +3; 1.18x size, 1.45x time;
% D in place of C: 51%, CI -29 to -7), production paragraph (exp 102: 19 -> 39 and 48 -> 54 of 72, 1.33x size), T10
% (tab:run2026-tl-noR: +5.2, -0.6 to 11.6).
\emph{Real backgrounds.} Circuits trained only against zero and means failed a resample test almost always; adding a resample term over the training contrast contexts raised the resample pass rate on $72$ targets from $0\%$ to $24\%$, at $1.18\times$ the circuit size and $1.45\times$ the fitting time, and with rank-keep raised it took a $72$-target check panel from $19$ to $39$ targets passing (\cref{app:resample}).
Using it in place of the contrast-mean term gained as much on the resample test but cost the pass criterion, so terms should be added for a new evaluation rather than swapped for an old one; on the production sweep the term's effect on the pass criterion itself is $5.2$ points and not significant (\cref{tab:run2026-tl-noR}).
\emph{Ablation methods and fills.} A mask trained under one ablation method learns what that method rewards: without the zero-ablation term, coefficients overshoot zero ablation by orders of magnitude, and circuits trained against dense means can fail under the sparsity-preserving fill that evaluation uses (\cref{app:single-method,app:respect-topk}).
\emph{Top-K survival.} Reproducing the pre-activation does not keep the target in its site's Top-K; the censored read and rank-keep train for it, though neither removes amplifier circuits (\cref{app:loss-details,app:specificity}).
\emph{Other backgrounds.} Contrast faithfulness generalises furthest when training uses random or distant contrast contexts (\cref{app:contrast-source}).
\emph{Size.} Circuits shrink with $\lambda$ at a nearly constant pass rate up to a plateau's edge (\cref{sec:lambda-transfer}).
Two evaluations have no training term yet: sufficiency to induce, which no circuit is trained on (median $0.78$--$0.88$ on held-out contrast contexts, $0.35$ for TuringLLM attention outputs; \cref{app:production}), and transfer to weaker, mid-band contexts (\cref{sec:main-result}).
The framework's prescription for them is a term or a training set in their domain, which we have not tested.
% CANDIDATE EVIDENCE, NOT IN THE PAPER (main.tex line 503: the arena results are reported separately). If Daniel wants
% them here, each needs its own sentence and its own appendix or citation:
%  - arena A tuning (experiments/107-true-baselines/arena_a/TUNING.md): on GPT-2 transcoder + z-SAE substrate the
%    necessity hinge (NEC_W; off in production, docs/wcm-specification.md §2.3) was chosen at 0.3 by a pre-registered
%    9-target rule; NEC_W 1.0 raised nec_a (3.tc.3752 0.32 -> 0.95) but pushed A out of band (1.54). The censored read
%    is inert there by construction (k = d_sae).
%  - arena A scope a1 (z + z2 + tc) vs a2 (tc only) and the wikitext-smoke vs OpenWebText corpus test (DAN-208 text;
%    not traced to an output file in this pass).
%  - logit / maximise / behavioural objectives (arena C/D; memory notes say "don't quote yet").
A fit trained for an evaluation is, for that evaluation, scored on home turf (\cref{app:method-comparison}), so the evaluation stays meaningful only if its contexts and donors are never seen in training, as the donor-background test's donors are not (\cref{tab:protocol}), and the evaluations not trained for are reported alongside.

\paragraph{Choosing the sparsity penalty.}
% sec:lambda-transfer, tab:run2026-tl-sweep, tab:run2026-lambda-pilot.
The penalty sets circuit size and, past a point, the pass rate: on the TuringLLM sweep the pass rate holds at $89$--$93\%$ from $\lambda = 2.5 \times 10^{-5}$ to $5 \times 10^{-4}$ while each doubling of $\lambda$ shrinks circuits $1.4$--$1.9\times$, and it declines beyond (\cref{sec:lambda-transfer}).
We recommend fixing the selection rule before looking at pass rates, as our size-only pilot rule did on the public models, where it kept $\lambda = 10^{-3}$ on both (\cref{tab:run2026-lambda-pilot}), and reporting a curve over $\lambda$ rather than one point, since one setting that transfers is not the best setting for each model.

\paragraph{A new model or dictionary.}
Only the model and dictionary adapters change between our three models (\cref{sec:setting}); four things needed checking.
\emph{Precision}: TuringLLM runs in the bfloat16 its pipeline was validated in, but on the public models a bfloat16 data pass did not reproduce the float32 top-context sets, so they run in float32 (\cref{tab:protocol}); a precision check belongs before any fit, since node sets also move with rounding (\cref{app:underdetermination}).
\emph{The dictionary's sparsity}: $k$ sets the fill's budget and the censored read's cut; on a dictionary that keeps every latent ($k = d_{\mathrm{sae}}$) the censored read has no cut to model, and on dense ReLU dictionaries a random-circuit baseline with fitted coefficients stops discriminating, so necessity, not faithfulness, separates real from random circuits there (\cref{app:dense-relu}).
% The k = d_sae statement follows from the definition of tau^(K) (no K-th other latent); docs/wcm-specification.md §7.6.
\emph{Contrast retrieval}: contrast contexts need a pooled representation for every corpus sequence and a silence check, which rejects $14$--$51\%$ of the closest candidates on the targets we inspected (\cref{app:stores}).
\emph{Corpus copies}: duplicated windows silently leak into held-out sets unless pools are deduplicated and split by groups of copies (\cref{app:heldout}).

\paragraph{What to report.}
Because each slot changes the circuits, a result is interpretable only with its recipe: the ablation methods used in training and in evaluation; faithfulness under each method and necessity, on held-out contexts, by depth band and site kind, with intervals over targets; excluded and vacuous targets counted; circuit size with whether coefficients were fitted, and the $\alpha = 1$ rescoring; random-circuit baselines given the same fitting; and the reported-only evaluations (sufficiency to induce, the donor-background test, specificity, mid-band transfer) beside the pass rate (\cref{sec:discussion}).

\paragraph{Choices that matter, and what stays open.}
The ablation method changes not only scores but rankings: the two mean-ablation variants can order small node sets differently (\cref{app:respect-topk}), and in exploratory comparisons a method tended to win the scores of the ablation method it was trained or ranked under (\cref{app:method-comparison}), so comparisons between recipes need the full matrix of recipe against ablation method.
Node sets are not identified by any recipe: a refit with reordered training contexts shares a median Jaccard of $0.39$--$0.51$ with the original and its shared core alone passes for $44$--$72\%$ of targets (\cref{app:underdetermination}), so changes to a slot should be judged by sizes and scores, not by which latents moved.
Finally, the production choices were calibrated on TuringLLM targets that overlap its evaluation set (\cref{sec:setting}), and on the public models every setting except $\lambda$ carried over unchanged (\cref{tab:protocol}, Provenance), so on a new model they are defaults to be checked, not tuned values.

% ---------------------------------------------------------------------------------------------------------------------
% 2026-11-08 (Daniel, via coordinator): the development path as evidence for where WCM can go next. Every number is
% traced; sources in the comment before each step. Panels and protocols differ between steps, so numbers are compared
% only within a step (each step is a paired test against the recipe before it).
% ---------------------------------------------------------------------------------------------------------------------
\subsection{How the Production Recipe Was Reached}
\label{app:framework-history}
The production recipe is the end of a sequence of changes, each made because a run exposed a failure, and each tested against the recipe before it on the same targets, mostly under a selection rule stated before the comparison.
We list them in order as evidence of which kinds of change have paid off; the panels differ between steps, so numbers compare only within a step.

% app:stores (38% fallback), app:contrast-source; experiments/055-close-contrast, 071-contrast-sweep.
\paragraph{Contrast contexts built per target.}
The first contrast contexts came from a retrieval store whose index covered too few sequences, so $38\%$ of TuringLLM targets fell back to one shared list; WCM now selects contrast contexts per target at fit time, ranked by similarity to its strongest contexts and kept only if the target is silent (\cref{app:stores}), and a later comparison of contrast sources showed how much that choice shapes the C term (\cref{app:contrast-source}).
The lesson: inspect the data each training term actually sees.

% app:specificity, app:rank-designs; experiments/056-specificity.
\paragraph{Rank-keep against off-target lifts.}
Reproducing the pre-activation left some circuits restoring their target by lifting its whole neighbourhood; of the terms tried, rank-keep alone kept the gains without the alternatives' costs, and entered the first protocol at a small weight (\cref{app:rank-designs,app:specificity}).

% app:v1-run (15,046 targets, 46%, attention 14%, layers 8-11 22%); app:heldout (66.8%); app:resample (0.1% vs 43.1%,
% 747 targets); fig:failure-anatomy.
\paragraph{The first run at scale.}
Under that protocol, WCM was fitted to $15{,}046$ TuringLLM targets, and $46\%$ passed, $22\%$ at layers 8--11 and $14\%$ for attention outputs (\cref{app:v1-run}).
Analyses of the run exposed three problems: $66.8\%$ of held-out strongest contexts repeated a training context, which led to deduplicated records split by groups of copies (\cref{app:heldout}); with depth, failing circuits increasingly failed only the two mean ablations (\cref{fig:failure-anatomy}); and $0.1\%$ of circuits passed a resample test, against $43.1\%$ under the pass criterion (\cref{app:resample}).

% experiments/093-v2-validation/README.md (out/tables.md): 72 targets, 8 per band x kind, paired, pass floor 0.8, the
% v1 run's as-run held-out sets (leaky; same for both arms): 31 -> 50 of 72, +26.4 pp [15.3, 38.9]; L8-11 2 -> 11 of 24;
% attn 3 -> 12 of 24; median dZ 0.00; size ratio 0.91; fit time 1.27x. Pilot: experiments/090-sparse-fill-training.
% Also main.tex comment at app:respect-topk ("The same fill in training").
\paragraph{Training under the fill that evaluation uses.}
The mean-ablation failures were a mismatch: training replaced non-members with dense means, a code no Top-K dictionary emits, while evaluation re-imposed the Top-K (\cref{app:respect-topk}).
Training the A and C terms under evaluation's own fill and means raised the pass rate on $72$ stratified targets from $31$ to $50$ (paired difference $+26.4$ points, $95\%$ interval $15.3$--$38.9$), from $2$ to $11$ of $24$ at layers 8--11 and from $3$ to $12$ of $24$ for attention outputs, with zero-ablation faithfulness unchanged (median paired change $0.00$), circuits $0.91\times$ the size, and $1.27\times$ the fitting time.
The gain came from aligning training with evaluation, not from a new term.

% app:resample (095 numbers and the production paragraph: 102's 19 -> 39, 48 -> 54 of 72, 1.33x size);
% experiments/102-pass-improvement/README.md line 89 (diagnosis: 15 of 18 Z failures of the previous version restored
% the pre-activation and lost at the Top-K cut); tab:run2026-tl-noR.
\paragraph{A resample term, and a stronger rank-keep.}
To survive real backgrounds, training gained the resample term (above and \cref{app:resample}).
On the previous version's zero-ablation failures, $15$ of $18$ restored the target's pre-activation but lost it at the Top-K cut, so rank-keep was raised from $3 \times 10^{-3}$ to $10^{-1}$ with the resample term at full weight; chosen by a rule stated before the comparison, the pair took a $72$-target check panel from $48$ to $54$ targets passing and from $19$ to $39$ passing the resample test, at $1.33\times$ the circuit size (\cref{app:resample}).
In the production sweep, the arm trained without the resample term passes $5.2$ points less, not significantly (\cref{tab:run2026-tl-noR}).

% experiments/107-true-baselines/latent_pilot/README.md (9 targets enriched for failures: 3 -> 5 of 9, size ratio ~1.04;
% plain geometric: circuits 2-4x larger, sum(lr) 6.3 vs 20; constant lr makes lr/T grow 20x) and
% experiments/107-true-baselines/adaptive_lambda/README.md, "Check" table (arm "geo x 3.15 + censored term
% (latent_pilot)": 63 of 72 vs 54 for the previous version, +10 / -1, size ratio 1.04 on pass-in-both). VERIFY: whether
% that check arm included the 0.1 warm-up (run_tonight.sh stage 4 fits check_geo_x3.15_warmup_censor).
\paragraph{A learning-rate schedule and the censored read.}
Under a constant learning rate, the gate step grows twentyfold as the gate temperature anneals, and some fits degraded late in training.
A geometric decay alone stabilised them but left gates too little total step to close, so circuits grew $2$--$4\times$; raising its peak $3.15\times$ restored the summed step size, and a short warm-up was added for the early spikes the higher peak brought back (\cref{app:loss-details}).
With the censored read added, so that training also sees whether the target survives its site's Top-K, the check panel went from $54$ to $63$ of $72$ targets passing ($10$ gained, $1$ lost) at about the same circuit size (median ratio $1.04$).

% sec:lambda-transfer, tab:run2026-lambda-pilot; experiments/108-batch-size/out/analysis.md (16 TuringLLM targets;
% B16 at 400 steps: 15 / 16 vs 16 / 16, size 0.75x on pass-in-both, 99.9 vs 29.1 fit-minutes = 3.4x; pre-stated rule:
% +-1 target is noise); experiments/107-true-baselines/adaptive_lambda/README.md (every shrinking policy traded passes
% for size on the check panel; losses were necessity failures).
\paragraph{Settings left as they were.}
The penalty was carried to the public models by a size-only rule fixed before a pilot (\cref{sec:lambda-transfer}).
In exploratory tests, batches of $16$ for the same $400$ steps passed $15$ of $16$ TuringLLM targets against $16$, with circuits $0.75\times$ the size where both passed, at $3.4\times$ the fitting time, so the batch of $4$ was kept; and policies that adapted $\lambda$ per target traded passes for size on the check panel, losing passes mostly to necessity, which no training term covers.
Other rejected designs are in \cref{app:membership-studies}.

% Directions: app:production (induce), sec:main-result (mid-band, deep band, attention), app:specificity (amplifiers,
% v1: 30% flagged), app:rank-designs (off-target term: lifts cut 3-7x, but grew circuits), app:resample (anchors
% only; overfiring away from the anchor), adaptive_lambda README (necessity).
\paragraph{Directions not yet taken.}
The same path suggests where the recipe can go next.
Necessity, sufficiency to induce, and transfer to mid-band contexts are scored but not trained for; each could get a term or training contexts in its domain, as the resample test did.
Deep and attention-output targets still pass least (\cref{sec:main-result}), and in the schedule pilot the attention failures kept zero-ablation faithfulness near $0$ under every schedule, so the learning rate does not explain them.
Amplifier circuits survive rank-keep (\cref{app:specificity}); the off-target term cut lifts three- to sevenfold but grew circuits (\cref{app:rank-designs}), so a cheaper specificity term remains open.
And circuits explain the target at its anchors, not at every position (\cref{app:resample}); training on more positions is untested.
\fi\fi
    % ~1.3 pages, its own appendix section
%
% New labels: sec:framework (main), app:framework, tab:framework (appendix). Every other \cref target exists in
% main.tex as of 2026-11-08. Every number below is already printed in main.tex (source comments give the place);
% the one exception, experiment 108, is marked NOT-YET-IN-PAPER where it appears.
% Vocabulary: paper/notes/vocabulary.md (no "endpoint", "floor", "drivers"; "fill" only as main.tex already uses it).
% Never-claim list respected: no "first"; nothing here claims novelty.
% =====================================================================================================================
\providecommand{\frameworkpart}{main}
\def\frameworkmain{main}
\def\frameworkappendix{appendix}
\ifx\frameworkpart\frameworkmain
% ---------------------------------------------------------------------------------------------------------------------
% (1) MAIN-TEXT VERSION. Proposed place: sec:discussion, in place of the paragraph "Ablation methods are part of the
% method" (main.tex lines 548-552), whose evidence it repeats and extends; the paragraph "Recommendations" (lines
% 564-566) can then point here. If it must sit in Method instead, it fits after "Pass criterion" (line 341), before
% \section{Experiments and Results}, but there it would cite results not yet shown.
% Numbers: 2-14% = abstract / run2026_alpha1 (sec:weighted); "one to three orders of magnitude" = app:single-method
% (exploratory); contrast sources = app:contrast-source (protocol v1); resample = app:resample.
% ---------------------------------------------------------------------------------------------------------------------
\subsection{Using WCM as a Framework}
\label{sec:framework}
WCM is better read as a framework for solving for circuits than as one fixed algorithm.
Its parts are interchangeable: an objective (the quantity reproduced, the ablation methods it is reproduced under, the regularisers, and their weights; \cref{eq:loss}), a parametrisation (a gate and a coefficient per node), data (training, contrast, and held-out contexts), and an optimiser.
The production recipe is one setting of these parts, each chosen by an experiment (\cref{tab:framework}), and changing a part changes the circuits found.
With every coefficient fixed at $1$ the same nodes pass for $2$--$14\%$ of targets (\cref{sec:weighted}); in exploratory fits, coefficients trained without zero ablation overshoot it by one to three orders of magnitude (\cref{app:single-method}); circuits trained against close contrast contexts are the most faithful on close backgrounds and the least on any other (\cref{app:contrast-source}); and circuits failed a resample test until training included a resample term (\cref{app:resample}).
When WCM fails an evaluation, we therefore suggest first asking whether training ever saw that evaluation's domain, and if not, adding the term, contexts, or setting that puts it there, while keeping the evaluation's own contexts and donors out of training so that it stays a held-out test (\cref{app:framework}).
A pass rate then describes a recipe, a dictionary, and an evaluation together, not WCM alone, which is why the ablation methods used for training and evaluation belong in the method statement~\citep{mueller2025mib, mahale2026explanation}.
%
\else\ifx\frameworkpart\frameworkappendix
% ---------------------------------------------------------------------------------------------------------------------
% (2) APPENDIX VERSION. Proposed place: a new appendix section directly before "Weighted Circuit Masking Design
% Ablations" (main.tex line 2003, app:membership-studies), so that the design ablations it summarises follow it; or
% directly after app:protocol / app:loss-details (before line 853) if Daniel prefers it next to tab:protocol.
% ---------------------------------------------------------------------------------------------------------------------
\section{Building with Weighted Circuit Masking}
\label{app:framework}
Weighted circuit masking fixes a problem, \labelcref{eq:p2}, and a way of relaxing it, but leaves most of the recipe open.
\Cref{tab:framework} lists its slots: what is reproduced, over which nodes, with which parametrisation, under which ablation methods and backgrounds, with which regularisers and optimiser, from which contexts, and how the result is scored.
The production choice in each slot was made by an experiment reported in these appendices, most of them on TuringLLM and carried over unchanged to the public models (\cref{tab:protocol}, Provenance); none is claimed to be the best setting, and several are known to trade one evaluation against another.

{\footnotesize
\setlength{\tabcolsep}{3pt}
\renewcommand{\arraystretch}{1.15}
\begin{longtable}{@{}p{0.13\linewidth}p{0.33\linewidth}p{0.33\linewidth}p{0.16\linewidth}@{}}
\caption{The slots of weighted circuit masking, the production choice in each, and the alternatives this paper examined or that the formulation admits.}
\label{tab:framework} \\
\toprule
\textbf{Slot} & \textbf{Production choice} & \textbf{Alternatives} & \textbf{Evidence} \\
\midrule
\endfirsthead
\toprule
\textbf{Slot} & \textbf{Production choice} & \textbf{Alternatives} & \textbf{Evidence} \\
\midrule
\endhead
\midrule
\multicolumn{4}{r@{}}{\emph{continued on the next page}} \\
\endfoot
\bottomrule
\endlastfoot
Target & the pre-activation $\preact$ of one SAE latent at its anchors; censored read under Z, C, A & an output metric, as in learned-mask circuits for outputs~\citep{bhaskar2024finding, caples2025scaling}; output- and latent-targeted circuits fail each other's objective (exploratory) & \cref{sec:ablmask}, \cref{app:secondary,app:logit-faith} \\
Nodes & every latent at every SAE site upstream of the target; SAE errors passed through; nodes recomputed & dense ReLU dictionaries (the method replicates; the random baseline needs its own design); members clamped (resample diagnostic) & \cref{sec:problem}, \cref{app:dense-relu,app:resample} \\
Parametrisation & gate $\sigma(\theta/T)$ and coefficient $\mathrm{softplus}(\psi)$ per node, $\alpha_0 = 1$ & $\alpha \equiv 1$ (unweighted circuit masking) & \cref{sec:weighted}, \cref{tab:run2026-alpha1} \\
Training ablations & Z (weight $1$), C and A ($0.25$ each) under evaluation's sparsity-preserving fill; resample D over the training contrast contexts ($1$) & one method alone; dense means; D in place of C; no D & \cref{app:single-method,app:respect-topk,app:resample} \\
Contrast background & $64$ close, silence-checked contexts per target ($48$ train, $16$ held out) & random, distant, or fused sources; the unchecked retrieval store & \cref{app:stores,app:contrast-source} \\
Regularisers & flat penalty $\lambda = 10^{-3}$ per open gate; closed-gate charge $\lambda(1-\sigma(\theta))|\alpha - 1|$; rank-keep $\gamma_{\mathrm{K}} = 10^{-1}$; one shared normaliser $\nu$ & other penalties $\lambda$ (sweep); off-target term, Top-K hinge, pull to rank $1$, firing margin; per-term normalisers & \cref{sec:lambda-transfer}, \cref{app:loss-details,app:rank-designs,app:membership-studies} \\
Optimiser & AdamW, $400$ steps, batch $4$; gate temperature annealed $1 \to 0.05$; learning-rate warm-up, then geometric decay & straight-through binarisation; $200$-step and reach-then-hold schedules; larger batches (exploratory) & \cref{tab:protocol}, \cref{app:membership-studies,app:framework-history} \\
Training contexts & $32$ strongest and $16$ mid-band activating contexts, deduplicated and split by whole groups of copies & a split without deduplication, which leaked copies into held-out sets & \cref{sec:setting}, \cref{app:heldout} \\
Evaluation & held-out strongest contexts; Z, A, C in $[0.75, 1.5]$ and necessity $\geq 0.9$; vacuous targets excluded and counted & reported, not required: mid-band contexts, sufficiency to induce, donor-background test, specificity, $\alpha = 1$ rescoring, random-circuit baselines, stable core & \cref{sec:eval}, \cref{app:protocol,app:nulls,app:underdetermination} \\
\end{longtable}}

\paragraph{Training in the evaluation's domain.}
Across these experiments, circuits did best on the evaluations their training resembled.
When a fit fails an evaluation, the failure is often a property of the recipe, and a slot of \cref{tab:framework} can put the fit in that evaluation's domain.
% app:resample (095: 72 targets, 0% -> 24%, CI +14 to +33; criterion 69% -> 62%, CI -17 to +3; 1.18x size, 1.45x time;
% D in place of C: 51%, CI -29 to -7), production paragraph (exp 102: 19 -> 39 and 48 -> 54 of 72, 1.33x size), T10
% (tab:run2026-tl-noR: +5.2, -0.6 to 11.6).
\emph{Real backgrounds.} Circuits trained only against zero and means failed a resample test almost always; adding a resample term over the training contrast contexts raised the resample pass rate on $72$ targets from $0\%$ to $24\%$, at $1.18\times$ the circuit size and $1.45\times$ the fitting time, and with rank-keep raised it took a $72$-target check panel from $19$ to $39$ targets passing (\cref{app:resample}).
Using it in place of the contrast-mean term gained as much on the resample test but cost the pass criterion, so terms should be added for a new evaluation rather than swapped for an old one; on the production sweep the term's effect on the pass criterion itself is $5.2$ points and not significant (\cref{tab:run2026-tl-noR}).
\emph{Ablation methods and fills.} A mask trained under one ablation method learns what that method rewards: without the zero-ablation term, coefficients overshoot zero ablation by orders of magnitude, and circuits trained against dense means can fail under the sparsity-preserving fill that evaluation uses (\cref{app:single-method,app:respect-topk}).
\emph{Top-K survival.} Reproducing the pre-activation does not keep the target in its site's Top-K; the censored read and rank-keep train for it, though neither removes amplifier circuits (\cref{app:loss-details,app:specificity}).
\emph{Other backgrounds.} Contrast faithfulness generalises furthest when training uses random or distant contrast contexts (\cref{app:contrast-source}).
\emph{Size.} Circuits shrink with $\lambda$ at a nearly constant pass rate up to a plateau's edge (\cref{sec:lambda-transfer}).
Two evaluations have no training term yet: sufficiency to induce, which no circuit is trained on (median $0.78$--$0.88$ on held-out contrast contexts, $0.35$ for TuringLLM attention outputs; \cref{app:production}), and transfer to weaker, mid-band contexts (\cref{sec:main-result}).
The framework's prescription for them is a term or a training set in their domain, which we have not tested.
% CANDIDATE EVIDENCE, NOT IN THE PAPER (main.tex line 503: the arena results are reported separately). If Daniel wants
% them here, each needs its own sentence and its own appendix or citation:
%  - arena A tuning (experiments/107-true-baselines/arena_a/TUNING.md): on GPT-2 transcoder + z-SAE substrate the
%    necessity hinge (NEC_W; off in production, docs/wcm-specification.md §2.3) was chosen at 0.3 by a pre-registered
%    9-target rule; NEC_W 1.0 raised nec_a (3.tc.3752 0.32 -> 0.95) but pushed A out of band (1.54). The censored read
%    is inert there by construction (k = d_sae).
%  - arena A scope a1 (z + z2 + tc) vs a2 (tc only) and the wikitext-smoke vs OpenWebText corpus test (DAN-208 text;
%    not traced to an output file in this pass).
%  - logit / maximise / behavioural objectives (arena C/D; memory notes say "don't quote yet").
A fit trained for an evaluation is, for that evaluation, scored on home turf (\cref{app:method-comparison}), so the evaluation stays meaningful only if its contexts and donors are never seen in training, as the donor-background test's donors are not (\cref{tab:protocol}), and the evaluations not trained for are reported alongside.

\paragraph{Choosing the sparsity penalty.}
% sec:lambda-transfer, tab:run2026-tl-sweep, tab:run2026-lambda-pilot.
The penalty sets circuit size and, past a point, the pass rate: on the TuringLLM sweep the pass rate holds at $89$--$93\%$ from $\lambda = 2.5 \times 10^{-5}$ to $5 \times 10^{-4}$ while each doubling of $\lambda$ shrinks circuits $1.4$--$1.9\times$, and it declines beyond (\cref{sec:lambda-transfer}).
We recommend fixing the selection rule before looking at pass rates, as our size-only pilot rule did on the public models, where it kept $\lambda = 10^{-3}$ on both (\cref{tab:run2026-lambda-pilot}), and reporting a curve over $\lambda$ rather than one point, since one setting that transfers is not the best setting for each model.

\paragraph{A new model or dictionary.}
Only the model and dictionary adapters change between our three models (\cref{sec:setting}); four things needed checking.
\emph{Precision}: TuringLLM runs in the bfloat16 its pipeline was validated in, but on the public models a bfloat16 data pass did not reproduce the float32 top-context sets, so they run in float32 (\cref{tab:protocol}); a precision check belongs before any fit, since node sets also move with rounding (\cref{app:underdetermination}).
\emph{The dictionary's sparsity}: $k$ sets the fill's budget and the censored read's cut; on a dictionary that keeps every latent ($k = d_{\mathrm{sae}}$) the censored read has no cut to model, and on dense ReLU dictionaries a random-circuit baseline with fitted coefficients stops discriminating, so necessity, not faithfulness, separates real from random circuits there (\cref{app:dense-relu}).
% The k = d_sae statement follows from the definition of tau^(K) (no K-th other latent); docs/wcm-specification.md §7.6.
\emph{Contrast retrieval}: contrast contexts need a pooled representation for every corpus sequence and a silence check, which rejects $14$--$51\%$ of the closest candidates on the targets we inspected (\cref{app:stores}).
\emph{Corpus copies}: duplicated windows silently leak into held-out sets unless pools are deduplicated and split by groups of copies (\cref{app:heldout}).

\paragraph{What to report.}
Because each slot changes the circuits, a result is interpretable only with its recipe: the ablation methods used in training and in evaluation; faithfulness under each method and necessity, on held-out contexts, by depth band and site kind, with intervals over targets; excluded and vacuous targets counted; circuit size with whether coefficients were fitted, and the $\alpha = 1$ rescoring; random-circuit baselines given the same fitting; and the reported-only evaluations (sufficiency to induce, the donor-background test, specificity, mid-band transfer) beside the pass rate (\cref{sec:discussion}).

\paragraph{Choices that matter, and what stays open.}
The ablation method changes not only scores but rankings: the two mean-ablation variants can order small node sets differently (\cref{app:respect-topk}), and in exploratory comparisons a method tended to win the scores of the ablation method it was trained or ranked under (\cref{app:method-comparison}), so comparisons between recipes need the full matrix of recipe against ablation method.
Node sets are not identified by any recipe: a refit with reordered training contexts shares a median Jaccard of $0.39$--$0.51$ with the original and its shared core alone passes for $44$--$72\%$ of targets (\cref{app:underdetermination}), so changes to a slot should be judged by sizes and scores, not by which latents moved.
Finally, the production choices were calibrated on TuringLLM targets that overlap its evaluation set (\cref{sec:setting}), and on the public models every setting except $\lambda$ carried over unchanged (\cref{tab:protocol}, Provenance), so on a new model they are defaults to be checked, not tuned values.

% ---------------------------------------------------------------------------------------------------------------------
% 2026-11-08 (Daniel, via coordinator): the development path as evidence for where WCM can go next. Every number is
% traced; sources in the comment before each step. Panels and protocols differ between steps, so numbers are compared
% only within a step (each step is a paired test against the recipe before it).
% ---------------------------------------------------------------------------------------------------------------------
\subsection{How the Production Recipe Was Reached}
\label{app:framework-history}
The production recipe is the end of a sequence of changes, each made because a run exposed a failure, and each tested against the recipe before it on the same targets, mostly under a selection rule stated before the comparison.
We list them in order as evidence of which kinds of change have paid off; the panels differ between steps, so numbers compare only within a step.

% app:stores (38% fallback), app:contrast-source; experiments/055-close-contrast, 071-contrast-sweep.
\paragraph{Contrast contexts built per target.}
The first contrast contexts came from a retrieval store whose index covered too few sequences, so $38\%$ of TuringLLM targets fell back to one shared list; WCM now selects contrast contexts per target at fit time, ranked by similarity to its strongest contexts and kept only if the target is silent (\cref{app:stores}), and a later comparison of contrast sources showed how much that choice shapes the C term (\cref{app:contrast-source}).
The lesson: inspect the data each training term actually sees.

% app:specificity, app:rank-designs; experiments/056-specificity.
\paragraph{Rank-keep against off-target lifts.}
Reproducing the pre-activation left some circuits restoring their target by lifting its whole neighbourhood; of the terms tried, rank-keep alone kept the gains without the alternatives' costs, and entered the first protocol at a small weight (\cref{app:rank-designs,app:specificity}).

% app:v1-run (15,046 targets, 46%, attention 14%, layers 8-11 22%); app:heldout (66.8%); app:resample (0.1% vs 43.1%,
% 747 targets); fig:failure-anatomy.
\paragraph{The first run at scale.}
Under that protocol, WCM was fitted to $15{,}046$ TuringLLM targets, and $46\%$ passed, $22\%$ at layers 8--11 and $14\%$ for attention outputs (\cref{app:v1-run}).
Analyses of the run exposed three problems: $66.8\%$ of held-out strongest contexts repeated a training context, which led to deduplicated records split by groups of copies (\cref{app:heldout}); with depth, failing circuits increasingly failed only the two mean ablations (\cref{fig:failure-anatomy}); and $0.1\%$ of circuits passed a resample test, against $43.1\%$ under the pass criterion (\cref{app:resample}).

% experiments/093-v2-validation/README.md (out/tables.md): 72 targets, 8 per band x kind, paired, pass floor 0.8, the
% v1 run's as-run held-out sets (leaky; same for both arms): 31 -> 50 of 72, +26.4 pp [15.3, 38.9]; L8-11 2 -> 11 of 24;
% attn 3 -> 12 of 24; median dZ 0.00; size ratio 0.91; fit time 1.27x. Pilot: experiments/090-sparse-fill-training.
% Also main.tex comment at app:respect-topk ("The same fill in training").
\paragraph{Training under the fill that evaluation uses.}
The mean-ablation failures were a mismatch: training replaced non-members with dense means, a code no Top-K dictionary emits, while evaluation re-imposed the Top-K (\cref{app:respect-topk}).
Training the A and C terms under evaluation's own fill and means raised the pass rate on $72$ stratified targets from $31$ to $50$ (paired difference $+26.4$ points, $95\%$ interval $15.3$--$38.9$), from $2$ to $11$ of $24$ at layers 8--11 and from $3$ to $12$ of $24$ for attention outputs, with zero-ablation faithfulness unchanged (median paired change $0.00$), circuits $0.91\times$ the size, and $1.27\times$ the fitting time.
The gain came from aligning training with evaluation, not from a new term.

% app:resample (095 numbers and the production paragraph: 102's 19 -> 39, 48 -> 54 of 72, 1.33x size);
% experiments/102-pass-improvement/README.md line 89 (diagnosis: 15 of 18 Z failures of the previous version restored
% the pre-activation and lost at the Top-K cut); tab:run2026-tl-noR.
\paragraph{A resample term, and a stronger rank-keep.}
To survive real backgrounds, training gained the resample term (above and \cref{app:resample}).
On the previous version's zero-ablation failures, $15$ of $18$ restored the target's pre-activation but lost it at the Top-K cut, so rank-keep was raised from $3 \times 10^{-3}$ to $10^{-1}$ with the resample term at full weight; chosen by a rule stated before the comparison, the pair took a $72$-target check panel from $48$ to $54$ targets passing and from $19$ to $39$ passing the resample test, at $1.33\times$ the circuit size (\cref{app:resample}).
In the production sweep, the arm trained without the resample term passes $5.2$ points less, not significantly (\cref{tab:run2026-tl-noR}).

% experiments/107-true-baselines/latent_pilot/README.md (9 targets enriched for failures: 3 -> 5 of 9, size ratio ~1.04;
% plain geometric: circuits 2-4x larger, sum(lr) 6.3 vs 20; constant lr makes lr/T grow 20x) and
% experiments/107-true-baselines/adaptive_lambda/README.md, "Check" table (arm "geo x 3.15 + censored term
% (latent_pilot)": 63 of 72 vs 54 for the previous version, +10 / -1, size ratio 1.04 on pass-in-both). VERIFY: whether
% that check arm included the 0.1 warm-up (run_tonight.sh stage 4 fits check_geo_x3.15_warmup_censor).
\paragraph{A learning-rate schedule and the censored read.}
Under a constant learning rate, the gate step grows twentyfold as the gate temperature anneals, and some fits degraded late in training.
A geometric decay alone stabilised them but left gates too little total step to close, so circuits grew $2$--$4\times$; raising its peak $3.15\times$ restored the summed step size, and a short warm-up was added for the early spikes the higher peak brought back (\cref{app:loss-details}).
With the censored read added, so that training also sees whether the target survives its site's Top-K, the check panel went from $54$ to $63$ of $72$ targets passing ($10$ gained, $1$ lost) at about the same circuit size (median ratio $1.04$).

% sec:lambda-transfer, tab:run2026-lambda-pilot; experiments/108-batch-size/out/analysis.md (16 TuringLLM targets;
% B16 at 400 steps: 15 / 16 vs 16 / 16, size 0.75x on pass-in-both, 99.9 vs 29.1 fit-minutes = 3.4x; pre-stated rule:
% +-1 target is noise); experiments/107-true-baselines/adaptive_lambda/README.md (every shrinking policy traded passes
% for size on the check panel; losses were necessity failures).
\paragraph{Settings left as they were.}
The penalty was carried to the public models by a size-only rule fixed before a pilot (\cref{sec:lambda-transfer}).
In exploratory tests, batches of $16$ for the same $400$ steps passed $15$ of $16$ TuringLLM targets against $16$, with circuits $0.75\times$ the size where both passed, at $3.4\times$ the fitting time, so the batch of $4$ was kept; and policies that adapted $\lambda$ per target traded passes for size on the check panel, losing passes mostly to necessity, which no training term covers.
Other rejected designs are in \cref{app:membership-studies}.

% Directions: app:production (induce), sec:main-result (mid-band, deep band, attention), app:specificity (amplifiers,
% v1: 30% flagged), app:rank-designs (off-target term: lifts cut 3-7x, but grew circuits), app:resample (anchors
% only; overfiring away from the anchor), adaptive_lambda README (necessity).
\paragraph{Directions not yet taken.}
The same path suggests where the recipe can go next.
Necessity, sufficiency to induce, and transfer to mid-band contexts are scored but not trained for; each could get a term or training contexts in its domain, as the resample test did.
Deep and attention-output targets still pass least (\cref{sec:main-result}), and in the schedule pilot the attention failures kept zero-ablation faithfulness near $0$ under every schedule, so the learning rate does not explain them.
Amplifier circuits survive rank-keep (\cref{app:specificity}); the off-target term cut lifts three- to sevenfold but grew circuits (\cref{app:rank-designs}), so a cheaper specificity term remains open.
And circuits explain the target at its anchors, not at every position (\cref{app:resample}); training on more positions is untested.
\fi\fi
    % ~1.3 pages, its own appendix section
%
% New labels: sec:framework (main), app:framework, tab:framework (appendix). Every other \cref target exists in
% main.tex as of 2026-11-08. Every number below is already printed in main.tex (source comments give the place);
% the one exception, experiment 108, is marked NOT-YET-IN-PAPER where it appears.
% Vocabulary: paper/notes/vocabulary.md (no "endpoint", "floor", "drivers"; "fill" only as main.tex already uses it).
% Never-claim list respected: no "first"; nothing here claims novelty.
% =====================================================================================================================
\providecommand{\frameworkpart}{main}
\def\frameworkmain{main}
\def\frameworkappendix{appendix}
\ifx\frameworkpart\frameworkmain
% ---------------------------------------------------------------------------------------------------------------------
% (1) MAIN-TEXT VERSION. Proposed place: sec:discussion, in place of the paragraph "Ablation methods are part of the
% method" (main.tex lines 548-552), whose evidence it repeats and extends; the paragraph "Recommendations" (lines
% 564-566) can then point here. If it must sit in Method instead, it fits after "Pass criterion" (line 341), before
% \section{Experiments and Results}, but there it would cite results not yet shown.
% Numbers: 2-14% = abstract / run2026_alpha1 (sec:weighted); "one to three orders of magnitude" = app:single-method
% (exploratory); contrast sources = app:contrast-source (protocol v1); resample = app:resample.
% ---------------------------------------------------------------------------------------------------------------------
\subsection{Using WCM as a Framework}
\label{sec:framework}
WCM is better read as a framework for solving for circuits than as one fixed algorithm.
Its parts are interchangeable: an objective (the quantity reproduced, the ablation methods it is reproduced under, the regularisers, and their weights; \cref{eq:loss}), a parametrisation (a gate and a coefficient per node), data (training, contrast, and held-out contexts), and an optimiser.
The production recipe is one setting of these parts, each chosen by an experiment (\cref{tab:framework}), and changing a part changes the circuits found.
With every coefficient fixed at $1$ the same nodes pass for $2$--$14\%$ of targets (\cref{sec:weighted}); in exploratory fits, coefficients trained without zero ablation overshoot it by one to three orders of magnitude (\cref{app:single-method}); circuits trained against close contrast contexts are the most faithful on close backgrounds and the least on any other (\cref{app:contrast-source}); and circuits failed a resample test until training included a resample term (\cref{app:resample}).
When WCM fails an evaluation, we therefore suggest first asking whether training ever saw that evaluation's domain, and if not, adding the term, contexts, or setting that puts it there, while keeping the evaluation's own contexts and donors out of training so that it stays a held-out test (\cref{app:framework}).
A pass rate then describes a recipe, a dictionary, and an evaluation together, not WCM alone, which is why the ablation methods used for training and evaluation belong in the method statement~\citep{mueller2025mib, mahale2026explanation}.
%
\else\ifx\frameworkpart\frameworkappendix
% ---------------------------------------------------------------------------------------------------------------------
% (2) APPENDIX VERSION. Proposed place: a new appendix section directly before "Weighted Circuit Masking Design
% Ablations" (main.tex line 2003, app:membership-studies), so that the design ablations it summarises follow it; or
% directly after app:protocol / app:loss-details (before line 853) if Daniel prefers it next to tab:protocol.
% ---------------------------------------------------------------------------------------------------------------------
\section{Building with Weighted Circuit Masking}
\label{app:framework}
Weighted circuit masking fixes a problem, \labelcref{eq:p2}, and a way of relaxing it, but leaves most of the recipe open.
\Cref{tab:framework} lists its slots: what is reproduced, over which nodes, with which parametrisation, under which ablation methods and backgrounds, with which regularisers and optimiser, from which contexts, and how the result is scored.
The production choice in each slot was made by an experiment reported in these appendices, most of them on TuringLLM and carried over unchanged to the public models (\cref{tab:protocol}, Provenance); none is claimed to be the best setting, and several are known to trade one evaluation against another.

{\footnotesize
\setlength{\tabcolsep}{3pt}
\renewcommand{\arraystretch}{1.15}
\begin{longtable}{@{}p{0.13\linewidth}p{0.33\linewidth}p{0.33\linewidth}p{0.16\linewidth}@{}}
\caption{The slots of weighted circuit masking, the production choice in each, and the alternatives this paper examined or that the formulation admits.}
\label{tab:framework} \\
\toprule
\textbf{Slot} & \textbf{Production choice} & \textbf{Alternatives} & \textbf{Evidence} \\
\midrule
\endfirsthead
\toprule
\textbf{Slot} & \textbf{Production choice} & \textbf{Alternatives} & \textbf{Evidence} \\
\midrule
\endhead
\midrule
\multicolumn{4}{r@{}}{\emph{continued on the next page}} \\
\endfoot
\bottomrule
\endlastfoot
Target & the pre-activation $\preact$ of one SAE latent at its anchors; censored read under Z, C, A & an output metric, as in learned-mask circuits for outputs~\citep{bhaskar2024finding, caples2025scaling}; output- and latent-targeted circuits fail each other's objective (exploratory) & \cref{sec:ablmask}, \cref{app:secondary,app:logit-faith} \\
Nodes & every latent at every SAE site upstream of the target; SAE errors passed through; nodes recomputed & dense ReLU dictionaries (the method replicates; the random baseline needs its own design); members clamped (resample diagnostic) & \cref{sec:problem}, \cref{app:dense-relu,app:resample} \\
Parametrisation & gate $\sigma(\theta/T)$ and coefficient $\mathrm{softplus}(\psi)$ per node, $\alpha_0 = 1$ & $\alpha \equiv 1$ (unweighted circuit masking) & \cref{sec:weighted}, \cref{tab:run2026-alpha1} \\
Training ablations & Z (weight $1$), C and A ($0.25$ each) under evaluation's sparsity-preserving fill; resample D over the training contrast contexts ($1$) & one method alone; dense means; D in place of C; no D & \cref{app:single-method,app:respect-topk,app:resample} \\
Contrast background & $64$ close, silence-checked contexts per target ($48$ train, $16$ held out) & random, distant, or fused sources; the unchecked retrieval store & \cref{app:stores,app:contrast-source} \\
Regularisers & flat penalty $\lambda = 10^{-3}$ per open gate; closed-gate charge $\lambda(1-\sigma(\theta))|\alpha - 1|$; rank-keep $\gamma_{\mathrm{K}} = 10^{-1}$; one shared normaliser $\nu$ & other penalties $\lambda$ (sweep); off-target term, Top-K hinge, pull to rank $1$, firing margin; per-term normalisers & \cref{sec:lambda-transfer}, \cref{app:loss-details,app:rank-designs,app:membership-studies} \\
Optimiser & AdamW, $400$ steps, batch $4$; gate temperature annealed $1 \to 0.05$; learning-rate warm-up, then geometric decay & straight-through binarisation; $200$-step and reach-then-hold schedules; larger batches (exploratory) & \cref{tab:protocol}, \cref{app:membership-studies,app:framework-history} \\
Training contexts & $32$ strongest and $16$ mid-band activating contexts, deduplicated and split by whole groups of copies & a split without deduplication, which leaked copies into held-out sets & \cref{sec:setting}, \cref{app:heldout} \\
Evaluation & held-out strongest contexts; Z, A, C in $[0.75, 1.5]$ and necessity $\geq 0.9$; vacuous targets excluded and counted & reported, not required: mid-band contexts, sufficiency to induce, donor-background test, specificity, $\alpha = 1$ rescoring, random-circuit baselines, stable core & \cref{sec:eval}, \cref{app:protocol,app:nulls,app:underdetermination} \\
\end{longtable}}

\paragraph{Training in the evaluation's domain.}
Across these experiments, circuits did best on the evaluations their training resembled.
When a fit fails an evaluation, the failure is often a property of the recipe, and a slot of \cref{tab:framework} can put the fit in that evaluation's domain.
% app:resample (095: 72 targets, 0% -> 24%, CI +14 to +33; criterion 69% -> 62%, CI -17 to +3; 1.18x size, 1.45x time;
% D in place of C: 51%, CI -29 to -7), production paragraph (exp 102: 19 -> 39 and 48 -> 54 of 72, 1.33x size), T10
% (tab:run2026-tl-noR: +5.2, -0.6 to 11.6).
\emph{Real backgrounds.} Circuits trained only against zero and means failed a resample test almost always; adding a resample term over the training contrast contexts raised the resample pass rate on $72$ targets from $0\%$ to $24\%$, at $1.18\times$ the circuit size and $1.45\times$ the fitting time, and with rank-keep raised it took a $72$-target check panel from $19$ to $39$ targets passing (\cref{app:resample}).
Using it in place of the contrast-mean term gained as much on the resample test but cost the pass criterion, so terms should be added for a new evaluation rather than swapped for an old one; on the production sweep the term's effect on the pass criterion itself is $5.2$ points and not significant (\cref{tab:run2026-tl-noR}).
\emph{Ablation methods and fills.} A mask trained under one ablation method learns what that method rewards: without the zero-ablation term, coefficients overshoot zero ablation by orders of magnitude, and circuits trained against dense means can fail under the sparsity-preserving fill that evaluation uses (\cref{app:single-method,app:respect-topk}).
\emph{Top-K survival.} Reproducing the pre-activation does not keep the target in its site's Top-K; the censored read and rank-keep train for it, though neither removes amplifier circuits (\cref{app:loss-details,app:specificity}).
\emph{Other backgrounds.} Contrast faithfulness generalises furthest when training uses random or distant contrast contexts (\cref{app:contrast-source}).
\emph{Size.} Circuits shrink with $\lambda$ at a nearly constant pass rate up to a plateau's edge (\cref{sec:lambda-transfer}).
Two evaluations have no training term yet: sufficiency to induce, which no circuit is trained on (median $0.78$--$0.88$ on held-out contrast contexts, $0.35$ for TuringLLM attention outputs; \cref{app:production}), and transfer to weaker, mid-band contexts (\cref{sec:main-result}).
The framework's prescription for them is a term or a training set in their domain, which we have not tested.
% CANDIDATE EVIDENCE, NOT IN THE PAPER (main.tex line 503: the arena results are reported separately). If Daniel wants
% them here, each needs its own sentence and its own appendix or citation:
%  - arena A tuning (experiments/107-true-baselines/arena_a/TUNING.md): on GPT-2 transcoder + z-SAE substrate the
%    necessity hinge (NEC_W; off in production, docs/wcm-specification.md §2.3) was chosen at 0.3 by a pre-registered
%    9-target rule; NEC_W 1.0 raised nec_a (3.tc.3752 0.32 -> 0.95) but pushed A out of band (1.54). The censored read
%    is inert there by construction (k = d_sae).
%  - arena A scope a1 (z + z2 + tc) vs a2 (tc only) and the wikitext-smoke vs OpenWebText corpus test (DAN-208 text;
%    not traced to an output file in this pass).
%  - logit / maximise / behavioural objectives (arena C/D; memory notes say "don't quote yet").
A fit trained for an evaluation is, for that evaluation, scored on home turf (\cref{app:method-comparison}), so the evaluation stays meaningful only if its contexts and donors are never seen in training, as the donor-background test's donors are not (\cref{tab:protocol}), and the evaluations not trained for are reported alongside.

\paragraph{Choosing the sparsity penalty.}
% sec:lambda-transfer, tab:run2026-tl-sweep, tab:run2026-lambda-pilot.
The penalty sets circuit size and, past a point, the pass rate: on the TuringLLM sweep the pass rate holds at $89$--$93\%$ from $\lambda = 2.5 \times 10^{-5}$ to $5 \times 10^{-4}$ while each doubling of $\lambda$ shrinks circuits $1.4$--$1.9\times$, and it declines beyond (\cref{sec:lambda-transfer}).
We recommend fixing the selection rule before looking at pass rates, as our size-only pilot rule did on the public models, where it kept $\lambda = 10^{-3}$ on both (\cref{tab:run2026-lambda-pilot}), and reporting a curve over $\lambda$ rather than one point, since one setting that transfers is not the best setting for each model.

\paragraph{A new model or dictionary.}
Only the model and dictionary adapters change between our three models (\cref{sec:setting}); four things needed checking.
\emph{Precision}: TuringLLM runs in the bfloat16 its pipeline was validated in, but on the public models a bfloat16 data pass did not reproduce the float32 top-context sets, so they run in float32 (\cref{tab:protocol}); a precision check belongs before any fit, since node sets also move with rounding (\cref{app:underdetermination}).
\emph{The dictionary's sparsity}: $k$ sets the fill's budget and the censored read's cut; on a dictionary that keeps every latent ($k = d_{\mathrm{sae}}$) the censored read has no cut to model, and on dense ReLU dictionaries a random-circuit baseline with fitted coefficients stops discriminating, so necessity, not faithfulness, separates real from random circuits there (\cref{app:dense-relu}).
% The k = d_sae statement follows from the definition of tau^(K) (no K-th other latent); docs/wcm-specification.md §7.6.
\emph{Contrast retrieval}: contrast contexts need a pooled representation for every corpus sequence and a silence check, which rejects $14$--$51\%$ of the closest candidates on the targets we inspected (\cref{app:stores}).
\emph{Corpus copies}: duplicated windows silently leak into held-out sets unless pools are deduplicated and split by groups of copies (\cref{app:heldout}).

\paragraph{What to report.}
Because each slot changes the circuits, a result is interpretable only with its recipe: the ablation methods used in training and in evaluation; faithfulness under each method and necessity, on held-out contexts, by depth band and site kind, with intervals over targets; excluded and vacuous targets counted; circuit size with whether coefficients were fitted, and the $\alpha = 1$ rescoring; random-circuit baselines given the same fitting; and the reported-only evaluations (sufficiency to induce, the donor-background test, specificity, mid-band transfer) beside the pass rate (\cref{sec:discussion}).

\paragraph{Choices that matter, and what stays open.}
The ablation method changes not only scores but rankings: the two mean-ablation variants can order small node sets differently (\cref{app:respect-topk}), and in exploratory comparisons a method tended to win the scores of the ablation method it was trained or ranked under (\cref{app:method-comparison}), so comparisons between recipes need the full matrix of recipe against ablation method.
Node sets are not identified by any recipe: a refit with reordered training contexts shares a median Jaccard of $0.39$--$0.51$ with the original and its shared core alone passes for $44$--$72\%$ of targets (\cref{app:underdetermination}), so changes to a slot should be judged by sizes and scores, not by which latents moved.
Finally, the production choices were calibrated on TuringLLM targets that overlap its evaluation set (\cref{sec:setting}), and on the public models every setting except $\lambda$ carried over unchanged (\cref{tab:protocol}, Provenance), so on a new model they are defaults to be checked, not tuned values.

% ---------------------------------------------------------------------------------------------------------------------
% 2026-11-08 (Daniel, via coordinator): the development path as evidence for where WCM can go next. Every number is
% traced; sources in the comment before each step. Panels and protocols differ between steps, so numbers are compared
% only within a step (each step is a paired test against the recipe before it).
% ---------------------------------------------------------------------------------------------------------------------
\subsection{How the Production Recipe Was Reached}
\label{app:framework-history}
The production recipe is the end of a sequence of changes, each made because a run exposed a failure, and each tested against the recipe before it on the same targets, mostly under a selection rule stated before the comparison.
We list them in order as evidence of which kinds of change have paid off; the panels differ between steps, so numbers compare only within a step.

% app:stores (38% fallback), app:contrast-source; experiments/055-close-contrast, 071-contrast-sweep.
\paragraph{Contrast contexts built per target.}
The first contrast contexts came from a retrieval store whose index covered too few sequences, so $38\%$ of TuringLLM targets fell back to one shared list; WCM now selects contrast contexts per target at fit time, ranked by similarity to its strongest contexts and kept only if the target is silent (\cref{app:stores}), and a later comparison of contrast sources showed how much that choice shapes the C term (\cref{app:contrast-source}).
The lesson: inspect the data each training term actually sees.

% app:specificity, app:rank-designs; experiments/056-specificity.
\paragraph{Rank-keep against off-target lifts.}
Reproducing the pre-activation left some circuits restoring their target by lifting its whole neighbourhood; of the terms tried, rank-keep alone kept the gains without the alternatives' costs, and entered the first protocol at a small weight (\cref{app:rank-designs,app:specificity}).

% app:v1-run (15,046 targets, 46%, attention 14%, layers 8-11 22%); app:heldout (66.8%); app:resample (0.1% vs 43.1%,
% 747 targets); fig:failure-anatomy.
\paragraph{The first run at scale.}
Under that protocol, WCM was fitted to $15{,}046$ TuringLLM targets, and $46\%$ passed, $22\%$ at layers 8--11 and $14\%$ for attention outputs (\cref{app:v1-run}).
Analyses of the run exposed three problems: $66.8\%$ of held-out strongest contexts repeated a training context, which led to deduplicated records split by groups of copies (\cref{app:heldout}); with depth, failing circuits increasingly failed only the two mean ablations (\cref{fig:failure-anatomy}); and $0.1\%$ of circuits passed a resample test, against $43.1\%$ under the pass criterion (\cref{app:resample}).

% experiments/093-v2-validation/README.md (out/tables.md): 72 targets, 8 per band x kind, paired, pass floor 0.8, the
% v1 run's as-run held-out sets (leaky; same for both arms): 31 -> 50 of 72, +26.4 pp [15.3, 38.9]; L8-11 2 -> 11 of 24;
% attn 3 -> 12 of 24; median dZ 0.00; size ratio 0.91; fit time 1.27x. Pilot: experiments/090-sparse-fill-training.
% Also main.tex comment at app:respect-topk ("The same fill in training").
\paragraph{Training under the fill that evaluation uses.}
The mean-ablation failures were a mismatch: training replaced non-members with dense means, a code no Top-K dictionary emits, while evaluation re-imposed the Top-K (\cref{app:respect-topk}).
Training the A and C terms under evaluation's own fill and means raised the pass rate on $72$ stratified targets from $31$ to $50$ (paired difference $+26.4$ points, $95\%$ interval $15.3$--$38.9$), from $2$ to $11$ of $24$ at layers 8--11 and from $3$ to $12$ of $24$ for attention outputs, with zero-ablation faithfulness unchanged (median paired change $0.00$), circuits $0.91\times$ the size, and $1.27\times$ the fitting time.
The gain came from aligning training with evaluation, not from a new term.

% app:resample (095 numbers and the production paragraph: 102's 19 -> 39, 48 -> 54 of 72, 1.33x size);
% experiments/102-pass-improvement/README.md line 89 (diagnosis: 15 of 18 Z failures of the previous version restored
% the pre-activation and lost at the Top-K cut); tab:run2026-tl-noR.
\paragraph{A resample term, and a stronger rank-keep.}
To survive real backgrounds, training gained the resample term (above and \cref{app:resample}).
On the previous version's zero-ablation failures, $15$ of $18$ restored the target's pre-activation but lost it at the Top-K cut, so rank-keep was raised from $3 \times 10^{-3}$ to $10^{-1}$ with the resample term at full weight; chosen by a rule stated before the comparison, the pair took a $72$-target check panel from $48$ to $54$ targets passing and from $19$ to $39$ passing the resample test, at $1.33\times$ the circuit size (\cref{app:resample}).
In the production sweep, the arm trained without the resample term passes $5.2$ points less, not significantly (\cref{tab:run2026-tl-noR}).

% experiments/107-true-baselines/latent_pilot/README.md (9 targets enriched for failures: 3 -> 5 of 9, size ratio ~1.04;
% plain geometric: circuits 2-4x larger, sum(lr) 6.3 vs 20; constant lr makes lr/T grow 20x) and
% experiments/107-true-baselines/adaptive_lambda/README.md, "Check" table (arm "geo x 3.15 + censored term
% (latent_pilot)": 63 of 72 vs 54 for the previous version, +10 / -1, size ratio 1.04 on pass-in-both). VERIFY: whether
% that check arm included the 0.1 warm-up (run_tonight.sh stage 4 fits check_geo_x3.15_warmup_censor).
\paragraph{A learning-rate schedule and the censored read.}
Under a constant learning rate, the gate step grows twentyfold as the gate temperature anneals, and some fits degraded late in training.
A geometric decay alone stabilised them but left gates too little total step to close, so circuits grew $2$--$4\times$; raising its peak $3.15\times$ restored the summed step size, and a short warm-up was added for the early spikes the higher peak brought back (\cref{app:loss-details}).
With the censored read added, so that training also sees whether the target survives its site's Top-K, the check panel went from $54$ to $63$ of $72$ targets passing ($10$ gained, $1$ lost) at about the same circuit size (median ratio $1.04$).

% sec:lambda-transfer, tab:run2026-lambda-pilot; experiments/108-batch-size/out/analysis.md (16 TuringLLM targets;
% B16 at 400 steps: 15 / 16 vs 16 / 16, size 0.75x on pass-in-both, 99.9 vs 29.1 fit-minutes = 3.4x; pre-stated rule:
% +-1 target is noise); experiments/107-true-baselines/adaptive_lambda/README.md (every shrinking policy traded passes
% for size on the check panel; losses were necessity failures).
\paragraph{Settings left as they were.}
The penalty was carried to the public models by a size-only rule fixed before a pilot (\cref{sec:lambda-transfer}).
In exploratory tests, batches of $16$ for the same $400$ steps passed $15$ of $16$ TuringLLM targets against $16$, with circuits $0.75\times$ the size where both passed, at $3.4\times$ the fitting time, so the batch of $4$ was kept; and policies that adapted $\lambda$ per target traded passes for size on the check panel, losing passes mostly to necessity, which no training term covers.
Other rejected designs are in \cref{app:membership-studies}.

% Directions: app:production (induce), sec:main-result (mid-band, deep band, attention), app:specificity (amplifiers,
% v1: 30% flagged), app:rank-designs (off-target term: lifts cut 3-7x, but grew circuits), app:resample (anchors
% only; overfiring away from the anchor), adaptive_lambda README (necessity).
\paragraph{Directions not yet taken.}
The same path suggests where the recipe can go next.
Necessity, sufficiency to induce, and transfer to mid-band contexts are scored but not trained for; each could get a term or training contexts in its domain, as the resample test did.
Deep and attention-output targets still pass least (\cref{sec:main-result}), and in the schedule pilot the attention failures kept zero-ablation faithfulness near $0$ under every schedule, so the learning rate does not explain them.
Amplifier circuits survive rank-keep (\cref{app:specificity}); the off-target term cut lifts three- to sevenfold but grew circuits (\cref{app:rank-designs}), so a cheaper specificity term remains open.
And circuits explain the target at its anchors, not at every position (\cref{app:resample}); training on more positions is untested.
\fi\fi
        % ~1/3 page, Discussion (replaces the paragraph
%                                                                    % "Ablation methods are part of the method")
%   \def\frameworkpart{appendix}% =====================================================================================================================
% DAN-208 (drafted 2026-11-08, agent "framework"): WCM as a framework for solving for circuits, not a fixed algorithm.
% Two versions in one file. Choose with \def\frameworkpart{main} or \def\frameworkpart{appendix} before \input:
%
%   \def\frameworkpart{main}% =====================================================================================================================
% DAN-208 (drafted 2026-11-08, agent "framework"): WCM as a framework for solving for circuits, not a fixed algorithm.
% Two versions in one file. Choose with \def\frameworkpart{main} or \def\frameworkpart{appendix} before \input:
%
%   \def\frameworkpart{main}% =====================================================================================================================
% DAN-208 (drafted 2026-11-08, agent "framework"): WCM as a framework for solving for circuits, not a fixed algorithm.
% Two versions in one file. Choose with \def\frameworkpart{main} or \def\frameworkpart{appendix} before \input:
%
%   \def\frameworkpart{main}\input{drafts/framework_section}        % ~1/3 page, Discussion (replaces the paragraph
%                                                                    % "Ablation methods are part of the method")
%   \def\frameworkpart{appendix}\input{drafts/framework_section}    % ~1.3 pages, its own appendix section
%
% New labels: sec:framework (main), app:framework, tab:framework (appendix). Every other \cref target exists in
% main.tex as of 2026-11-08. Every number below is already printed in main.tex (source comments give the place);
% the one exception, experiment 108, is marked NOT-YET-IN-PAPER where it appears.
% Vocabulary: paper/notes/vocabulary.md (no "endpoint", "floor", "drivers"; "fill" only as main.tex already uses it).
% Never-claim list respected: no "first"; nothing here claims novelty.
% =====================================================================================================================
\providecommand{\frameworkpart}{main}
\def\frameworkmain{main}
\def\frameworkappendix{appendix}
\ifx\frameworkpart\frameworkmain
% ---------------------------------------------------------------------------------------------------------------------
% (1) MAIN-TEXT VERSION. Proposed place: sec:discussion, in place of the paragraph "Ablation methods are part of the
% method" (main.tex lines 548-552), whose evidence it repeats and extends; the paragraph "Recommendations" (lines
% 564-566) can then point here. If it must sit in Method instead, it fits after "Pass criterion" (line 341), before
% \section{Experiments and Results}, but there it would cite results not yet shown.
% Numbers: 2-14% = abstract / run2026_alpha1 (sec:weighted); "one to three orders of magnitude" = app:single-method
% (exploratory); contrast sources = app:contrast-source (protocol v1); resample = app:resample.
% ---------------------------------------------------------------------------------------------------------------------
\subsection{Using WCM as a Framework}
\label{sec:framework}
WCM is better read as a framework for solving for circuits than as one fixed algorithm.
Its parts are interchangeable: an objective (the quantity reproduced, the ablation methods it is reproduced under, the regularisers, and their weights; \cref{eq:loss}), a parametrisation (a gate and a coefficient per node), data (training, contrast, and held-out contexts), and an optimiser.
The production recipe is one setting of these parts, each chosen by an experiment (\cref{tab:framework}), and changing a part changes the circuits found.
With every coefficient fixed at $1$ the same nodes pass for $2$--$14\%$ of targets (\cref{sec:weighted}); in exploratory fits, coefficients trained without zero ablation overshoot it by one to three orders of magnitude (\cref{app:single-method}); circuits trained against close contrast contexts are the most faithful on close backgrounds and the least on any other (\cref{app:contrast-source}); and circuits failed a resample test until training included a resample term (\cref{app:resample}).
When WCM fails an evaluation, we therefore suggest first asking whether training ever saw that evaluation's domain, and if not, adding the term, contexts, or setting that puts it there, while keeping the evaluation's own contexts and donors out of training so that it stays a held-out test (\cref{app:framework}).
A pass rate then describes a recipe, a dictionary, and an evaluation together, not WCM alone, which is why the ablation methods used for training and evaluation belong in the method statement~\citep{mueller2025mib, mahale2026explanation}.
%
\else\ifx\frameworkpart\frameworkappendix
% ---------------------------------------------------------------------------------------------------------------------
% (2) APPENDIX VERSION. Proposed place: a new appendix section directly before "Weighted Circuit Masking Design
% Ablations" (main.tex line 2003, app:membership-studies), so that the design ablations it summarises follow it; or
% directly after app:protocol / app:loss-details (before line 853) if Daniel prefers it next to tab:protocol.
% ---------------------------------------------------------------------------------------------------------------------
\section{Building with Weighted Circuit Masking}
\label{app:framework}
Weighted circuit masking fixes a problem, \labelcref{eq:p2}, and a way of relaxing it, but leaves most of the recipe open.
\Cref{tab:framework} lists its slots: what is reproduced, over which nodes, with which parametrisation, under which ablation methods and backgrounds, with which regularisers and optimiser, from which contexts, and how the result is scored.
The production choice in each slot was made by an experiment reported in these appendices, most of them on TuringLLM and carried over unchanged to the public models (\cref{tab:protocol}, Provenance); none is claimed to be the best setting, and several are known to trade one evaluation against another.

{\footnotesize
\setlength{\tabcolsep}{3pt}
\renewcommand{\arraystretch}{1.15}
\begin{longtable}{@{}p{0.13\linewidth}p{0.33\linewidth}p{0.33\linewidth}p{0.16\linewidth}@{}}
\caption{The slots of weighted circuit masking, the production choice in each, and the alternatives this paper examined or that the formulation admits.}
\label{tab:framework} \\
\toprule
\textbf{Slot} & \textbf{Production choice} & \textbf{Alternatives} & \textbf{Evidence} \\
\midrule
\endfirsthead
\toprule
\textbf{Slot} & \textbf{Production choice} & \textbf{Alternatives} & \textbf{Evidence} \\
\midrule
\endhead
\midrule
\multicolumn{4}{r@{}}{\emph{continued on the next page}} \\
\endfoot
\bottomrule
\endlastfoot
Target & the pre-activation $\preact$ of one SAE latent at its anchors; censored read under Z, C, A & an output metric, as in learned-mask circuits for outputs~\citep{bhaskar2024finding, caples2025scaling}; output- and latent-targeted circuits fail each other's objective (exploratory) & \cref{sec:ablmask}, \cref{app:secondary,app:logit-faith} \\
Nodes & every latent at every SAE site upstream of the target; SAE errors passed through; nodes recomputed & dense ReLU dictionaries (the method replicates; the random baseline needs its own design); members clamped (resample diagnostic) & \cref{sec:problem}, \cref{app:dense-relu,app:resample} \\
Parametrisation & gate $\sigma(\theta/T)$ and coefficient $\mathrm{softplus}(\psi)$ per node, $\alpha_0 = 1$ & $\alpha \equiv 1$ (unweighted circuit masking) & \cref{sec:weighted}, \cref{tab:run2026-alpha1} \\
Training ablations & Z (weight $1$), C and A ($0.25$ each) under evaluation's sparsity-preserving fill; resample D over the training contrast contexts ($1$) & one method alone; dense means; D in place of C; no D & \cref{app:single-method,app:respect-topk,app:resample} \\
Contrast background & $64$ close, silence-checked contexts per target ($48$ train, $16$ held out) & random, distant, or fused sources; the unchecked retrieval store & \cref{app:stores,app:contrast-source} \\
Regularisers & flat penalty $\lambda = 10^{-3}$ per open gate; closed-gate charge $\lambda(1-\sigma(\theta))|\alpha - 1|$; rank-keep $\gamma_{\mathrm{K}} = 10^{-1}$; one shared normaliser $\nu$ & other penalties $\lambda$ (sweep); off-target term, Top-K hinge, pull to rank $1$, firing margin; per-term normalisers & \cref{sec:lambda-transfer}, \cref{app:loss-details,app:rank-designs,app:membership-studies} \\
Optimiser & AdamW, $400$ steps, batch $4$; gate temperature annealed $1 \to 0.05$; learning-rate warm-up, then geometric decay & straight-through binarisation; $200$-step and reach-then-hold schedules; larger batches (exploratory) & \cref{tab:protocol}, \cref{app:membership-studies,app:framework-history} \\
Training contexts & $32$ strongest and $16$ mid-band activating contexts, deduplicated and split by whole groups of copies & a split without deduplication, which leaked copies into held-out sets & \cref{sec:setting}, \cref{app:heldout} \\
Evaluation & held-out strongest contexts; Z, A, C in $[0.75, 1.5]$ and necessity $\geq 0.9$; vacuous targets excluded and counted & reported, not required: mid-band contexts, sufficiency to induce, donor-background test, specificity, $\alpha = 1$ rescoring, random-circuit baselines, stable core & \cref{sec:eval}, \cref{app:protocol,app:nulls,app:underdetermination} \\
\end{longtable}}

\paragraph{Training in the evaluation's domain.}
Across these experiments, circuits did best on the evaluations their training resembled.
When a fit fails an evaluation, the failure is often a property of the recipe, and a slot of \cref{tab:framework} can put the fit in that evaluation's domain.
% app:resample (095: 72 targets, 0% -> 24%, CI +14 to +33; criterion 69% -> 62%, CI -17 to +3; 1.18x size, 1.45x time;
% D in place of C: 51%, CI -29 to -7), production paragraph (exp 102: 19 -> 39 and 48 -> 54 of 72, 1.33x size), T10
% (tab:run2026-tl-noR: +5.2, -0.6 to 11.6).
\emph{Real backgrounds.} Circuits trained only against zero and means failed a resample test almost always; adding a resample term over the training contrast contexts raised the resample pass rate on $72$ targets from $0\%$ to $24\%$, at $1.18\times$ the circuit size and $1.45\times$ the fitting time, and with rank-keep raised it took a $72$-target check panel from $19$ to $39$ targets passing (\cref{app:resample}).
Using it in place of the contrast-mean term gained as much on the resample test but cost the pass criterion, so terms should be added for a new evaluation rather than swapped for an old one; on the production sweep the term's effect on the pass criterion itself is $5.2$ points and not significant (\cref{tab:run2026-tl-noR}).
\emph{Ablation methods and fills.} A mask trained under one ablation method learns what that method rewards: without the zero-ablation term, coefficients overshoot zero ablation by orders of magnitude, and circuits trained against dense means can fail under the sparsity-preserving fill that evaluation uses (\cref{app:single-method,app:respect-topk}).
\emph{Top-K survival.} Reproducing the pre-activation does not keep the target in its site's Top-K; the censored read and rank-keep train for it, though neither removes amplifier circuits (\cref{app:loss-details,app:specificity}).
\emph{Other backgrounds.} Contrast faithfulness generalises furthest when training uses random or distant contrast contexts (\cref{app:contrast-source}).
\emph{Size.} Circuits shrink with $\lambda$ at a nearly constant pass rate up to a plateau's edge (\cref{sec:lambda-transfer}).
Two evaluations have no training term yet: sufficiency to induce, which no circuit is trained on (median $0.78$--$0.88$ on held-out contrast contexts, $0.35$ for TuringLLM attention outputs; \cref{app:production}), and transfer to weaker, mid-band contexts (\cref{sec:main-result}).
The framework's prescription for them is a term or a training set in their domain, which we have not tested.
% CANDIDATE EVIDENCE, NOT IN THE PAPER (main.tex line 503: the arena results are reported separately). If Daniel wants
% them here, each needs its own sentence and its own appendix or citation:
%  - arena A tuning (experiments/107-true-baselines/arena_a/TUNING.md): on GPT-2 transcoder + z-SAE substrate the
%    necessity hinge (NEC_W; off in production, docs/wcm-specification.md §2.3) was chosen at 0.3 by a pre-registered
%    9-target rule; NEC_W 1.0 raised nec_a (3.tc.3752 0.32 -> 0.95) but pushed A out of band (1.54). The censored read
%    is inert there by construction (k = d_sae).
%  - arena A scope a1 (z + z2 + tc) vs a2 (tc only) and the wikitext-smoke vs OpenWebText corpus test (DAN-208 text;
%    not traced to an output file in this pass).
%  - logit / maximise / behavioural objectives (arena C/D; memory notes say "don't quote yet").
A fit trained for an evaluation is, for that evaluation, scored on home turf (\cref{app:method-comparison}), so the evaluation stays meaningful only if its contexts and donors are never seen in training, as the donor-background test's donors are not (\cref{tab:protocol}), and the evaluations not trained for are reported alongside.

\paragraph{Choosing the sparsity penalty.}
% sec:lambda-transfer, tab:run2026-tl-sweep, tab:run2026-lambda-pilot.
The penalty sets circuit size and, past a point, the pass rate: on the TuringLLM sweep the pass rate holds at $89$--$93\%$ from $\lambda = 2.5 \times 10^{-5}$ to $5 \times 10^{-4}$ while each doubling of $\lambda$ shrinks circuits $1.4$--$1.9\times$, and it declines beyond (\cref{sec:lambda-transfer}).
We recommend fixing the selection rule before looking at pass rates, as our size-only pilot rule did on the public models, where it kept $\lambda = 10^{-3}$ on both (\cref{tab:run2026-lambda-pilot}), and reporting a curve over $\lambda$ rather than one point, since one setting that transfers is not the best setting for each model.

\paragraph{A new model or dictionary.}
Only the model and dictionary adapters change between our three models (\cref{sec:setting}); four things needed checking.
\emph{Precision}: TuringLLM runs in the bfloat16 its pipeline was validated in, but on the public models a bfloat16 data pass did not reproduce the float32 top-context sets, so they run in float32 (\cref{tab:protocol}); a precision check belongs before any fit, since node sets also move with rounding (\cref{app:underdetermination}).
\emph{The dictionary's sparsity}: $k$ sets the fill's budget and the censored read's cut; on a dictionary that keeps every latent ($k = d_{\mathrm{sae}}$) the censored read has no cut to model, and on dense ReLU dictionaries a random-circuit baseline with fitted coefficients stops discriminating, so necessity, not faithfulness, separates real from random circuits there (\cref{app:dense-relu}).
% The k = d_sae statement follows from the definition of tau^(K) (no K-th other latent); docs/wcm-specification.md §7.6.
\emph{Contrast retrieval}: contrast contexts need a pooled representation for every corpus sequence and a silence check, which rejects $14$--$51\%$ of the closest candidates on the targets we inspected (\cref{app:stores}).
\emph{Corpus copies}: duplicated windows silently leak into held-out sets unless pools are deduplicated and split by groups of copies (\cref{app:heldout}).

\paragraph{What to report.}
Because each slot changes the circuits, a result is interpretable only with its recipe: the ablation methods used in training and in evaluation; faithfulness under each method and necessity, on held-out contexts, by depth band and site kind, with intervals over targets; excluded and vacuous targets counted; circuit size with whether coefficients were fitted, and the $\alpha = 1$ rescoring; random-circuit baselines given the same fitting; and the reported-only evaluations (sufficiency to induce, the donor-background test, specificity, mid-band transfer) beside the pass rate (\cref{sec:discussion}).

\paragraph{Choices that matter, and what stays open.}
The ablation method changes not only scores but rankings: the two mean-ablation variants can order small node sets differently (\cref{app:respect-topk}), and in exploratory comparisons a method tended to win the scores of the ablation method it was trained or ranked under (\cref{app:method-comparison}), so comparisons between recipes need the full matrix of recipe against ablation method.
Node sets are not identified by any recipe: a refit with reordered training contexts shares a median Jaccard of $0.39$--$0.51$ with the original and its shared core alone passes for $44$--$72\%$ of targets (\cref{app:underdetermination}), so changes to a slot should be judged by sizes and scores, not by which latents moved.
Finally, the production choices were calibrated on TuringLLM targets that overlap its evaluation set (\cref{sec:setting}), and on the public models every setting except $\lambda$ carried over unchanged (\cref{tab:protocol}, Provenance), so on a new model they are defaults to be checked, not tuned values.

% ---------------------------------------------------------------------------------------------------------------------
% 2026-11-08 (Daniel, via coordinator): the development path as evidence for where WCM can go next. Every number is
% traced; sources in the comment before each step. Panels and protocols differ between steps, so numbers are compared
% only within a step (each step is a paired test against the recipe before it).
% ---------------------------------------------------------------------------------------------------------------------
\subsection{How the Production Recipe Was Reached}
\label{app:framework-history}
The production recipe is the end of a sequence of changes, each made because a run exposed a failure, and each tested against the recipe before it on the same targets, mostly under a selection rule stated before the comparison.
We list them in order as evidence of which kinds of change have paid off; the panels differ between steps, so numbers compare only within a step.

% app:stores (38% fallback), app:contrast-source; experiments/055-close-contrast, 071-contrast-sweep.
\paragraph{Contrast contexts built per target.}
The first contrast contexts came from a retrieval store whose index covered too few sequences, so $38\%$ of TuringLLM targets fell back to one shared list; WCM now selects contrast contexts per target at fit time, ranked by similarity to its strongest contexts and kept only if the target is silent (\cref{app:stores}), and a later comparison of contrast sources showed how much that choice shapes the C term (\cref{app:contrast-source}).
The lesson: inspect the data each training term actually sees.

% app:specificity, app:rank-designs; experiments/056-specificity.
\paragraph{Rank-keep against off-target lifts.}
Reproducing the pre-activation left some circuits restoring their target by lifting its whole neighbourhood; of the terms tried, rank-keep alone kept the gains without the alternatives' costs, and entered the first protocol at a small weight (\cref{app:rank-designs,app:specificity}).

% app:v1-run (15,046 targets, 46%, attention 14%, layers 8-11 22%); app:heldout (66.8%); app:resample (0.1% vs 43.1%,
% 747 targets); fig:failure-anatomy.
\paragraph{The first run at scale.}
Under that protocol, WCM was fitted to $15{,}046$ TuringLLM targets, and $46\%$ passed, $22\%$ at layers 8--11 and $14\%$ for attention outputs (\cref{app:v1-run}).
Analyses of the run exposed three problems: $66.8\%$ of held-out strongest contexts repeated a training context, which led to deduplicated records split by groups of copies (\cref{app:heldout}); with depth, failing circuits increasingly failed only the two mean ablations (\cref{fig:failure-anatomy}); and $0.1\%$ of circuits passed a resample test, against $43.1\%$ under the pass criterion (\cref{app:resample}).

% experiments/093-v2-validation/README.md (out/tables.md): 72 targets, 8 per band x kind, paired, pass floor 0.8, the
% v1 run's as-run held-out sets (leaky; same for both arms): 31 -> 50 of 72, +26.4 pp [15.3, 38.9]; L8-11 2 -> 11 of 24;
% attn 3 -> 12 of 24; median dZ 0.00; size ratio 0.91; fit time 1.27x. Pilot: experiments/090-sparse-fill-training.
% Also main.tex comment at app:respect-topk ("The same fill in training").
\paragraph{Training under the fill that evaluation uses.}
The mean-ablation failures were a mismatch: training replaced non-members with dense means, a code no Top-K dictionary emits, while evaluation re-imposed the Top-K (\cref{app:respect-topk}).
Training the A and C terms under evaluation's own fill and means raised the pass rate on $72$ stratified targets from $31$ to $50$ (paired difference $+26.4$ points, $95\%$ interval $15.3$--$38.9$), from $2$ to $11$ of $24$ at layers 8--11 and from $3$ to $12$ of $24$ for attention outputs, with zero-ablation faithfulness unchanged (median paired change $0.00$), circuits $0.91\times$ the size, and $1.27\times$ the fitting time.
The gain came from aligning training with evaluation, not from a new term.

% app:resample (095 numbers and the production paragraph: 102's 19 -> 39, 48 -> 54 of 72, 1.33x size);
% experiments/102-pass-improvement/README.md line 89 (diagnosis: 15 of 18 Z failures of the previous version restored
% the pre-activation and lost at the Top-K cut); tab:run2026-tl-noR.
\paragraph{A resample term, and a stronger rank-keep.}
To survive real backgrounds, training gained the resample term (above and \cref{app:resample}).
On the previous version's zero-ablation failures, $15$ of $18$ restored the target's pre-activation but lost it at the Top-K cut, so rank-keep was raised from $3 \times 10^{-3}$ to $10^{-1}$ with the resample term at full weight; chosen by a rule stated before the comparison, the pair took a $72$-target check panel from $48$ to $54$ targets passing and from $19$ to $39$ passing the resample test, at $1.33\times$ the circuit size (\cref{app:resample}).
In the production sweep, the arm trained without the resample term passes $5.2$ points less, not significantly (\cref{tab:run2026-tl-noR}).

% experiments/107-true-baselines/latent_pilot/README.md (9 targets enriched for failures: 3 -> 5 of 9, size ratio ~1.04;
% plain geometric: circuits 2-4x larger, sum(lr) 6.3 vs 20; constant lr makes lr/T grow 20x) and
% experiments/107-true-baselines/adaptive_lambda/README.md, "Check" table (arm "geo x 3.15 + censored term
% (latent_pilot)": 63 of 72 vs 54 for the previous version, +10 / -1, size ratio 1.04 on pass-in-both). VERIFY: whether
% that check arm included the 0.1 warm-up (run_tonight.sh stage 4 fits check_geo_x3.15_warmup_censor).
\paragraph{A learning-rate schedule and the censored read.}
Under a constant learning rate, the gate step grows twentyfold as the gate temperature anneals, and some fits degraded late in training.
A geometric decay alone stabilised them but left gates too little total step to close, so circuits grew $2$--$4\times$; raising its peak $3.15\times$ restored the summed step size, and a short warm-up was added for the early spikes the higher peak brought back (\cref{app:loss-details}).
With the censored read added, so that training also sees whether the target survives its site's Top-K, the check panel went from $54$ to $63$ of $72$ targets passing ($10$ gained, $1$ lost) at about the same circuit size (median ratio $1.04$).

% sec:lambda-transfer, tab:run2026-lambda-pilot; experiments/108-batch-size/out/analysis.md (16 TuringLLM targets;
% B16 at 400 steps: 15 / 16 vs 16 / 16, size 0.75x on pass-in-both, 99.9 vs 29.1 fit-minutes = 3.4x; pre-stated rule:
% +-1 target is noise); experiments/107-true-baselines/adaptive_lambda/README.md (every shrinking policy traded passes
% for size on the check panel; losses were necessity failures).
\paragraph{Settings left as they were.}
The penalty was carried to the public models by a size-only rule fixed before a pilot (\cref{sec:lambda-transfer}).
In exploratory tests, batches of $16$ for the same $400$ steps passed $15$ of $16$ TuringLLM targets against $16$, with circuits $0.75\times$ the size where both passed, at $3.4\times$ the fitting time, so the batch of $4$ was kept; and policies that adapted $\lambda$ per target traded passes for size on the check panel, losing passes mostly to necessity, which no training term covers.
Other rejected designs are in \cref{app:membership-studies}.

% Directions: app:production (induce), sec:main-result (mid-band, deep band, attention), app:specificity (amplifiers,
% v1: 30% flagged), app:rank-designs (off-target term: lifts cut 3-7x, but grew circuits), app:resample (anchors
% only; overfiring away from the anchor), adaptive_lambda README (necessity).
\paragraph{Directions not yet taken.}
The same path suggests where the recipe can go next.
Necessity, sufficiency to induce, and transfer to mid-band contexts are scored but not trained for; each could get a term or training contexts in its domain, as the resample test did.
Deep and attention-output targets still pass least (\cref{sec:main-result}), and in the schedule pilot the attention failures kept zero-ablation faithfulness near $0$ under every schedule, so the learning rate does not explain them.
Amplifier circuits survive rank-keep (\cref{app:specificity}); the off-target term cut lifts three- to sevenfold but grew circuits (\cref{app:rank-designs}), so a cheaper specificity term remains open.
And circuits explain the target at its anchors, not at every position (\cref{app:resample}); training on more positions is untested.
\fi\fi
        % ~1/3 page, Discussion (replaces the paragraph
%                                                                    % "Ablation methods are part of the method")
%   \def\frameworkpart{appendix}% =====================================================================================================================
% DAN-208 (drafted 2026-11-08, agent "framework"): WCM as a framework for solving for circuits, not a fixed algorithm.
% Two versions in one file. Choose with \def\frameworkpart{main} or \def\frameworkpart{appendix} before \input:
%
%   \def\frameworkpart{main}\input{drafts/framework_section}        % ~1/3 page, Discussion (replaces the paragraph
%                                                                    % "Ablation methods are part of the method")
%   \def\frameworkpart{appendix}\input{drafts/framework_section}    % ~1.3 pages, its own appendix section
%
% New labels: sec:framework (main), app:framework, tab:framework (appendix). Every other \cref target exists in
% main.tex as of 2026-11-08. Every number below is already printed in main.tex (source comments give the place);
% the one exception, experiment 108, is marked NOT-YET-IN-PAPER where it appears.
% Vocabulary: paper/notes/vocabulary.md (no "endpoint", "floor", "drivers"; "fill" only as main.tex already uses it).
% Never-claim list respected: no "first"; nothing here claims novelty.
% =====================================================================================================================
\providecommand{\frameworkpart}{main}
\def\frameworkmain{main}
\def\frameworkappendix{appendix}
\ifx\frameworkpart\frameworkmain
% ---------------------------------------------------------------------------------------------------------------------
% (1) MAIN-TEXT VERSION. Proposed place: sec:discussion, in place of the paragraph "Ablation methods are part of the
% method" (main.tex lines 548-552), whose evidence it repeats and extends; the paragraph "Recommendations" (lines
% 564-566) can then point here. If it must sit in Method instead, it fits after "Pass criterion" (line 341), before
% \section{Experiments and Results}, but there it would cite results not yet shown.
% Numbers: 2-14% = abstract / run2026_alpha1 (sec:weighted); "one to three orders of magnitude" = app:single-method
% (exploratory); contrast sources = app:contrast-source (protocol v1); resample = app:resample.
% ---------------------------------------------------------------------------------------------------------------------
\subsection{Using WCM as a Framework}
\label{sec:framework}
WCM is better read as a framework for solving for circuits than as one fixed algorithm.
Its parts are interchangeable: an objective (the quantity reproduced, the ablation methods it is reproduced under, the regularisers, and their weights; \cref{eq:loss}), a parametrisation (a gate and a coefficient per node), data (training, contrast, and held-out contexts), and an optimiser.
The production recipe is one setting of these parts, each chosen by an experiment (\cref{tab:framework}), and changing a part changes the circuits found.
With every coefficient fixed at $1$ the same nodes pass for $2$--$14\%$ of targets (\cref{sec:weighted}); in exploratory fits, coefficients trained without zero ablation overshoot it by one to three orders of magnitude (\cref{app:single-method}); circuits trained against close contrast contexts are the most faithful on close backgrounds and the least on any other (\cref{app:contrast-source}); and circuits failed a resample test until training included a resample term (\cref{app:resample}).
When WCM fails an evaluation, we therefore suggest first asking whether training ever saw that evaluation's domain, and if not, adding the term, contexts, or setting that puts it there, while keeping the evaluation's own contexts and donors out of training so that it stays a held-out test (\cref{app:framework}).
A pass rate then describes a recipe, a dictionary, and an evaluation together, not WCM alone, which is why the ablation methods used for training and evaluation belong in the method statement~\citep{mueller2025mib, mahale2026explanation}.
%
\else\ifx\frameworkpart\frameworkappendix
% ---------------------------------------------------------------------------------------------------------------------
% (2) APPENDIX VERSION. Proposed place: a new appendix section directly before "Weighted Circuit Masking Design
% Ablations" (main.tex line 2003, app:membership-studies), so that the design ablations it summarises follow it; or
% directly after app:protocol / app:loss-details (before line 853) if Daniel prefers it next to tab:protocol.
% ---------------------------------------------------------------------------------------------------------------------
\section{Building with Weighted Circuit Masking}
\label{app:framework}
Weighted circuit masking fixes a problem, \labelcref{eq:p2}, and a way of relaxing it, but leaves most of the recipe open.
\Cref{tab:framework} lists its slots: what is reproduced, over which nodes, with which parametrisation, under which ablation methods and backgrounds, with which regularisers and optimiser, from which contexts, and how the result is scored.
The production choice in each slot was made by an experiment reported in these appendices, most of them on TuringLLM and carried over unchanged to the public models (\cref{tab:protocol}, Provenance); none is claimed to be the best setting, and several are known to trade one evaluation against another.

{\footnotesize
\setlength{\tabcolsep}{3pt}
\renewcommand{\arraystretch}{1.15}
\begin{longtable}{@{}p{0.13\linewidth}p{0.33\linewidth}p{0.33\linewidth}p{0.16\linewidth}@{}}
\caption{The slots of weighted circuit masking, the production choice in each, and the alternatives this paper examined or that the formulation admits.}
\label{tab:framework} \\
\toprule
\textbf{Slot} & \textbf{Production choice} & \textbf{Alternatives} & \textbf{Evidence} \\
\midrule
\endfirsthead
\toprule
\textbf{Slot} & \textbf{Production choice} & \textbf{Alternatives} & \textbf{Evidence} \\
\midrule
\endhead
\midrule
\multicolumn{4}{r@{}}{\emph{continued on the next page}} \\
\endfoot
\bottomrule
\endlastfoot
Target & the pre-activation $\preact$ of one SAE latent at its anchors; censored read under Z, C, A & an output metric, as in learned-mask circuits for outputs~\citep{bhaskar2024finding, caples2025scaling}; output- and latent-targeted circuits fail each other's objective (exploratory) & \cref{sec:ablmask}, \cref{app:secondary,app:logit-faith} \\
Nodes & every latent at every SAE site upstream of the target; SAE errors passed through; nodes recomputed & dense ReLU dictionaries (the method replicates; the random baseline needs its own design); members clamped (resample diagnostic) & \cref{sec:problem}, \cref{app:dense-relu,app:resample} \\
Parametrisation & gate $\sigma(\theta/T)$ and coefficient $\mathrm{softplus}(\psi)$ per node, $\alpha_0 = 1$ & $\alpha \equiv 1$ (unweighted circuit masking) & \cref{sec:weighted}, \cref{tab:run2026-alpha1} \\
Training ablations & Z (weight $1$), C and A ($0.25$ each) under evaluation's sparsity-preserving fill; resample D over the training contrast contexts ($1$) & one method alone; dense means; D in place of C; no D & \cref{app:single-method,app:respect-topk,app:resample} \\
Contrast background & $64$ close, silence-checked contexts per target ($48$ train, $16$ held out) & random, distant, or fused sources; the unchecked retrieval store & \cref{app:stores,app:contrast-source} \\
Regularisers & flat penalty $\lambda = 10^{-3}$ per open gate; closed-gate charge $\lambda(1-\sigma(\theta))|\alpha - 1|$; rank-keep $\gamma_{\mathrm{K}} = 10^{-1}$; one shared normaliser $\nu$ & other penalties $\lambda$ (sweep); off-target term, Top-K hinge, pull to rank $1$, firing margin; per-term normalisers & \cref{sec:lambda-transfer}, \cref{app:loss-details,app:rank-designs,app:membership-studies} \\
Optimiser & AdamW, $400$ steps, batch $4$; gate temperature annealed $1 \to 0.05$; learning-rate warm-up, then geometric decay & straight-through binarisation; $200$-step and reach-then-hold schedules; larger batches (exploratory) & \cref{tab:protocol}, \cref{app:membership-studies,app:framework-history} \\
Training contexts & $32$ strongest and $16$ mid-band activating contexts, deduplicated and split by whole groups of copies & a split without deduplication, which leaked copies into held-out sets & \cref{sec:setting}, \cref{app:heldout} \\
Evaluation & held-out strongest contexts; Z, A, C in $[0.75, 1.5]$ and necessity $\geq 0.9$; vacuous targets excluded and counted & reported, not required: mid-band contexts, sufficiency to induce, donor-background test, specificity, $\alpha = 1$ rescoring, random-circuit baselines, stable core & \cref{sec:eval}, \cref{app:protocol,app:nulls,app:underdetermination} \\
\end{longtable}}

\paragraph{Training in the evaluation's domain.}
Across these experiments, circuits did best on the evaluations their training resembled.
When a fit fails an evaluation, the failure is often a property of the recipe, and a slot of \cref{tab:framework} can put the fit in that evaluation's domain.
% app:resample (095: 72 targets, 0% -> 24%, CI +14 to +33; criterion 69% -> 62%, CI -17 to +3; 1.18x size, 1.45x time;
% D in place of C: 51%, CI -29 to -7), production paragraph (exp 102: 19 -> 39 and 48 -> 54 of 72, 1.33x size), T10
% (tab:run2026-tl-noR: +5.2, -0.6 to 11.6).
\emph{Real backgrounds.} Circuits trained only against zero and means failed a resample test almost always; adding a resample term over the training contrast contexts raised the resample pass rate on $72$ targets from $0\%$ to $24\%$, at $1.18\times$ the circuit size and $1.45\times$ the fitting time, and with rank-keep raised it took a $72$-target check panel from $19$ to $39$ targets passing (\cref{app:resample}).
Using it in place of the contrast-mean term gained as much on the resample test but cost the pass criterion, so terms should be added for a new evaluation rather than swapped for an old one; on the production sweep the term's effect on the pass criterion itself is $5.2$ points and not significant (\cref{tab:run2026-tl-noR}).
\emph{Ablation methods and fills.} A mask trained under one ablation method learns what that method rewards: without the zero-ablation term, coefficients overshoot zero ablation by orders of magnitude, and circuits trained against dense means can fail under the sparsity-preserving fill that evaluation uses (\cref{app:single-method,app:respect-topk}).
\emph{Top-K survival.} Reproducing the pre-activation does not keep the target in its site's Top-K; the censored read and rank-keep train for it, though neither removes amplifier circuits (\cref{app:loss-details,app:specificity}).
\emph{Other backgrounds.} Contrast faithfulness generalises furthest when training uses random or distant contrast contexts (\cref{app:contrast-source}).
\emph{Size.} Circuits shrink with $\lambda$ at a nearly constant pass rate up to a plateau's edge (\cref{sec:lambda-transfer}).
Two evaluations have no training term yet: sufficiency to induce, which no circuit is trained on (median $0.78$--$0.88$ on held-out contrast contexts, $0.35$ for TuringLLM attention outputs; \cref{app:production}), and transfer to weaker, mid-band contexts (\cref{sec:main-result}).
The framework's prescription for them is a term or a training set in their domain, which we have not tested.
% CANDIDATE EVIDENCE, NOT IN THE PAPER (main.tex line 503: the arena results are reported separately). If Daniel wants
% them here, each needs its own sentence and its own appendix or citation:
%  - arena A tuning (experiments/107-true-baselines/arena_a/TUNING.md): on GPT-2 transcoder + z-SAE substrate the
%    necessity hinge (NEC_W; off in production, docs/wcm-specification.md §2.3) was chosen at 0.3 by a pre-registered
%    9-target rule; NEC_W 1.0 raised nec_a (3.tc.3752 0.32 -> 0.95) but pushed A out of band (1.54). The censored read
%    is inert there by construction (k = d_sae).
%  - arena A scope a1 (z + z2 + tc) vs a2 (tc only) and the wikitext-smoke vs OpenWebText corpus test (DAN-208 text;
%    not traced to an output file in this pass).
%  - logit / maximise / behavioural objectives (arena C/D; memory notes say "don't quote yet").
A fit trained for an evaluation is, for that evaluation, scored on home turf (\cref{app:method-comparison}), so the evaluation stays meaningful only if its contexts and donors are never seen in training, as the donor-background test's donors are not (\cref{tab:protocol}), and the evaluations not trained for are reported alongside.

\paragraph{Choosing the sparsity penalty.}
% sec:lambda-transfer, tab:run2026-tl-sweep, tab:run2026-lambda-pilot.
The penalty sets circuit size and, past a point, the pass rate: on the TuringLLM sweep the pass rate holds at $89$--$93\%$ from $\lambda = 2.5 \times 10^{-5}$ to $5 \times 10^{-4}$ while each doubling of $\lambda$ shrinks circuits $1.4$--$1.9\times$, and it declines beyond (\cref{sec:lambda-transfer}).
We recommend fixing the selection rule before looking at pass rates, as our size-only pilot rule did on the public models, where it kept $\lambda = 10^{-3}$ on both (\cref{tab:run2026-lambda-pilot}), and reporting a curve over $\lambda$ rather than one point, since one setting that transfers is not the best setting for each model.

\paragraph{A new model or dictionary.}
Only the model and dictionary adapters change between our three models (\cref{sec:setting}); four things needed checking.
\emph{Precision}: TuringLLM runs in the bfloat16 its pipeline was validated in, but on the public models a bfloat16 data pass did not reproduce the float32 top-context sets, so they run in float32 (\cref{tab:protocol}); a precision check belongs before any fit, since node sets also move with rounding (\cref{app:underdetermination}).
\emph{The dictionary's sparsity}: $k$ sets the fill's budget and the censored read's cut; on a dictionary that keeps every latent ($k = d_{\mathrm{sae}}$) the censored read has no cut to model, and on dense ReLU dictionaries a random-circuit baseline with fitted coefficients stops discriminating, so necessity, not faithfulness, separates real from random circuits there (\cref{app:dense-relu}).
% The k = d_sae statement follows from the definition of tau^(K) (no K-th other latent); docs/wcm-specification.md §7.6.
\emph{Contrast retrieval}: contrast contexts need a pooled representation for every corpus sequence and a silence check, which rejects $14$--$51\%$ of the closest candidates on the targets we inspected (\cref{app:stores}).
\emph{Corpus copies}: duplicated windows silently leak into held-out sets unless pools are deduplicated and split by groups of copies (\cref{app:heldout}).

\paragraph{What to report.}
Because each slot changes the circuits, a result is interpretable only with its recipe: the ablation methods used in training and in evaluation; faithfulness under each method and necessity, on held-out contexts, by depth band and site kind, with intervals over targets; excluded and vacuous targets counted; circuit size with whether coefficients were fitted, and the $\alpha = 1$ rescoring; random-circuit baselines given the same fitting; and the reported-only evaluations (sufficiency to induce, the donor-background test, specificity, mid-band transfer) beside the pass rate (\cref{sec:discussion}).

\paragraph{Choices that matter, and what stays open.}
The ablation method changes not only scores but rankings: the two mean-ablation variants can order small node sets differently (\cref{app:respect-topk}), and in exploratory comparisons a method tended to win the scores of the ablation method it was trained or ranked under (\cref{app:method-comparison}), so comparisons between recipes need the full matrix of recipe against ablation method.
Node sets are not identified by any recipe: a refit with reordered training contexts shares a median Jaccard of $0.39$--$0.51$ with the original and its shared core alone passes for $44$--$72\%$ of targets (\cref{app:underdetermination}), so changes to a slot should be judged by sizes and scores, not by which latents moved.
Finally, the production choices were calibrated on TuringLLM targets that overlap its evaluation set (\cref{sec:setting}), and on the public models every setting except $\lambda$ carried over unchanged (\cref{tab:protocol}, Provenance), so on a new model they are defaults to be checked, not tuned values.

% ---------------------------------------------------------------------------------------------------------------------
% 2026-11-08 (Daniel, via coordinator): the development path as evidence for where WCM can go next. Every number is
% traced; sources in the comment before each step. Panels and protocols differ between steps, so numbers are compared
% only within a step (each step is a paired test against the recipe before it).
% ---------------------------------------------------------------------------------------------------------------------
\subsection{How the Production Recipe Was Reached}
\label{app:framework-history}
The production recipe is the end of a sequence of changes, each made because a run exposed a failure, and each tested against the recipe before it on the same targets, mostly under a selection rule stated before the comparison.
We list them in order as evidence of which kinds of change have paid off; the panels differ between steps, so numbers compare only within a step.

% app:stores (38% fallback), app:contrast-source; experiments/055-close-contrast, 071-contrast-sweep.
\paragraph{Contrast contexts built per target.}
The first contrast contexts came from a retrieval store whose index covered too few sequences, so $38\%$ of TuringLLM targets fell back to one shared list; WCM now selects contrast contexts per target at fit time, ranked by similarity to its strongest contexts and kept only if the target is silent (\cref{app:stores}), and a later comparison of contrast sources showed how much that choice shapes the C term (\cref{app:contrast-source}).
The lesson: inspect the data each training term actually sees.

% app:specificity, app:rank-designs; experiments/056-specificity.
\paragraph{Rank-keep against off-target lifts.}
Reproducing the pre-activation left some circuits restoring their target by lifting its whole neighbourhood; of the terms tried, rank-keep alone kept the gains without the alternatives' costs, and entered the first protocol at a small weight (\cref{app:rank-designs,app:specificity}).

% app:v1-run (15,046 targets, 46%, attention 14%, layers 8-11 22%); app:heldout (66.8%); app:resample (0.1% vs 43.1%,
% 747 targets); fig:failure-anatomy.
\paragraph{The first run at scale.}
Under that protocol, WCM was fitted to $15{,}046$ TuringLLM targets, and $46\%$ passed, $22\%$ at layers 8--11 and $14\%$ for attention outputs (\cref{app:v1-run}).
Analyses of the run exposed three problems: $66.8\%$ of held-out strongest contexts repeated a training context, which led to deduplicated records split by groups of copies (\cref{app:heldout}); with depth, failing circuits increasingly failed only the two mean ablations (\cref{fig:failure-anatomy}); and $0.1\%$ of circuits passed a resample test, against $43.1\%$ under the pass criterion (\cref{app:resample}).

% experiments/093-v2-validation/README.md (out/tables.md): 72 targets, 8 per band x kind, paired, pass floor 0.8, the
% v1 run's as-run held-out sets (leaky; same for both arms): 31 -> 50 of 72, +26.4 pp [15.3, 38.9]; L8-11 2 -> 11 of 24;
% attn 3 -> 12 of 24; median dZ 0.00; size ratio 0.91; fit time 1.27x. Pilot: experiments/090-sparse-fill-training.
% Also main.tex comment at app:respect-topk ("The same fill in training").
\paragraph{Training under the fill that evaluation uses.}
The mean-ablation failures were a mismatch: training replaced non-members with dense means, a code no Top-K dictionary emits, while evaluation re-imposed the Top-K (\cref{app:respect-topk}).
Training the A and C terms under evaluation's own fill and means raised the pass rate on $72$ stratified targets from $31$ to $50$ (paired difference $+26.4$ points, $95\%$ interval $15.3$--$38.9$), from $2$ to $11$ of $24$ at layers 8--11 and from $3$ to $12$ of $24$ for attention outputs, with zero-ablation faithfulness unchanged (median paired change $0.00$), circuits $0.91\times$ the size, and $1.27\times$ the fitting time.
The gain came from aligning training with evaluation, not from a new term.

% app:resample (095 numbers and the production paragraph: 102's 19 -> 39, 48 -> 54 of 72, 1.33x size);
% experiments/102-pass-improvement/README.md line 89 (diagnosis: 15 of 18 Z failures of the previous version restored
% the pre-activation and lost at the Top-K cut); tab:run2026-tl-noR.
\paragraph{A resample term, and a stronger rank-keep.}
To survive real backgrounds, training gained the resample term (above and \cref{app:resample}).
On the previous version's zero-ablation failures, $15$ of $18$ restored the target's pre-activation but lost it at the Top-K cut, so rank-keep was raised from $3 \times 10^{-3}$ to $10^{-1}$ with the resample term at full weight; chosen by a rule stated before the comparison, the pair took a $72$-target check panel from $48$ to $54$ targets passing and from $19$ to $39$ passing the resample test, at $1.33\times$ the circuit size (\cref{app:resample}).
In the production sweep, the arm trained without the resample term passes $5.2$ points less, not significantly (\cref{tab:run2026-tl-noR}).

% experiments/107-true-baselines/latent_pilot/README.md (9 targets enriched for failures: 3 -> 5 of 9, size ratio ~1.04;
% plain geometric: circuits 2-4x larger, sum(lr) 6.3 vs 20; constant lr makes lr/T grow 20x) and
% experiments/107-true-baselines/adaptive_lambda/README.md, "Check" table (arm "geo x 3.15 + censored term
% (latent_pilot)": 63 of 72 vs 54 for the previous version, +10 / -1, size ratio 1.04 on pass-in-both). VERIFY: whether
% that check arm included the 0.1 warm-up (run_tonight.sh stage 4 fits check_geo_x3.15_warmup_censor).
\paragraph{A learning-rate schedule and the censored read.}
Under a constant learning rate, the gate step grows twentyfold as the gate temperature anneals, and some fits degraded late in training.
A geometric decay alone stabilised them but left gates too little total step to close, so circuits grew $2$--$4\times$; raising its peak $3.15\times$ restored the summed step size, and a short warm-up was added for the early spikes the higher peak brought back (\cref{app:loss-details}).
With the censored read added, so that training also sees whether the target survives its site's Top-K, the check panel went from $54$ to $63$ of $72$ targets passing ($10$ gained, $1$ lost) at about the same circuit size (median ratio $1.04$).

% sec:lambda-transfer, tab:run2026-lambda-pilot; experiments/108-batch-size/out/analysis.md (16 TuringLLM targets;
% B16 at 400 steps: 15 / 16 vs 16 / 16, size 0.75x on pass-in-both, 99.9 vs 29.1 fit-minutes = 3.4x; pre-stated rule:
% +-1 target is noise); experiments/107-true-baselines/adaptive_lambda/README.md (every shrinking policy traded passes
% for size on the check panel; losses were necessity failures).
\paragraph{Settings left as they were.}
The penalty was carried to the public models by a size-only rule fixed before a pilot (\cref{sec:lambda-transfer}).
In exploratory tests, batches of $16$ for the same $400$ steps passed $15$ of $16$ TuringLLM targets against $16$, with circuits $0.75\times$ the size where both passed, at $3.4\times$ the fitting time, so the batch of $4$ was kept; and policies that adapted $\lambda$ per target traded passes for size on the check panel, losing passes mostly to necessity, which no training term covers.
Other rejected designs are in \cref{app:membership-studies}.

% Directions: app:production (induce), sec:main-result (mid-band, deep band, attention), app:specificity (amplifiers,
% v1: 30% flagged), app:rank-designs (off-target term: lifts cut 3-7x, but grew circuits), app:resample (anchors
% only; overfiring away from the anchor), adaptive_lambda README (necessity).
\paragraph{Directions not yet taken.}
The same path suggests where the recipe can go next.
Necessity, sufficiency to induce, and transfer to mid-band contexts are scored but not trained for; each could get a term or training contexts in its domain, as the resample test did.
Deep and attention-output targets still pass least (\cref{sec:main-result}), and in the schedule pilot the attention failures kept zero-ablation faithfulness near $0$ under every schedule, so the learning rate does not explain them.
Amplifier circuits survive rank-keep (\cref{app:specificity}); the off-target term cut lifts three- to sevenfold but grew circuits (\cref{app:rank-designs}), so a cheaper specificity term remains open.
And circuits explain the target at its anchors, not at every position (\cref{app:resample}); training on more positions is untested.
\fi\fi
    % ~1.3 pages, its own appendix section
%
% New labels: sec:framework (main), app:framework, tab:framework (appendix). Every other \cref target exists in
% main.tex as of 2026-11-08. Every number below is already printed in main.tex (source comments give the place);
% the one exception, experiment 108, is marked NOT-YET-IN-PAPER where it appears.
% Vocabulary: paper/notes/vocabulary.md (no "endpoint", "floor", "drivers"; "fill" only as main.tex already uses it).
% Never-claim list respected: no "first"; nothing here claims novelty.
% =====================================================================================================================
\providecommand{\frameworkpart}{main}
\def\frameworkmain{main}
\def\frameworkappendix{appendix}
\ifx\frameworkpart\frameworkmain
% ---------------------------------------------------------------------------------------------------------------------
% (1) MAIN-TEXT VERSION. Proposed place: sec:discussion, in place of the paragraph "Ablation methods are part of the
% method" (main.tex lines 548-552), whose evidence it repeats and extends; the paragraph "Recommendations" (lines
% 564-566) can then point here. If it must sit in Method instead, it fits after "Pass criterion" (line 341), before
% \section{Experiments and Results}, but there it would cite results not yet shown.
% Numbers: 2-14% = abstract / run2026_alpha1 (sec:weighted); "one to three orders of magnitude" = app:single-method
% (exploratory); contrast sources = app:contrast-source (protocol v1); resample = app:resample.
% ---------------------------------------------------------------------------------------------------------------------
\subsection{Using WCM as a Framework}
\label{sec:framework}
WCM is better read as a framework for solving for circuits than as one fixed algorithm.
Its parts are interchangeable: an objective (the quantity reproduced, the ablation methods it is reproduced under, the regularisers, and their weights; \cref{eq:loss}), a parametrisation (a gate and a coefficient per node), data (training, contrast, and held-out contexts), and an optimiser.
The production recipe is one setting of these parts, each chosen by an experiment (\cref{tab:framework}), and changing a part changes the circuits found.
With every coefficient fixed at $1$ the same nodes pass for $2$--$14\%$ of targets (\cref{sec:weighted}); in exploratory fits, coefficients trained without zero ablation overshoot it by one to three orders of magnitude (\cref{app:single-method}); circuits trained against close contrast contexts are the most faithful on close backgrounds and the least on any other (\cref{app:contrast-source}); and circuits failed a resample test until training included a resample term (\cref{app:resample}).
When WCM fails an evaluation, we therefore suggest first asking whether training ever saw that evaluation's domain, and if not, adding the term, contexts, or setting that puts it there, while keeping the evaluation's own contexts and donors out of training so that it stays a held-out test (\cref{app:framework}).
A pass rate then describes a recipe, a dictionary, and an evaluation together, not WCM alone, which is why the ablation methods used for training and evaluation belong in the method statement~\citep{mueller2025mib, mahale2026explanation}.
%
\else\ifx\frameworkpart\frameworkappendix
% ---------------------------------------------------------------------------------------------------------------------
% (2) APPENDIX VERSION. Proposed place: a new appendix section directly before "Weighted Circuit Masking Design
% Ablations" (main.tex line 2003, app:membership-studies), so that the design ablations it summarises follow it; or
% directly after app:protocol / app:loss-details (before line 853) if Daniel prefers it next to tab:protocol.
% ---------------------------------------------------------------------------------------------------------------------
\section{Building with Weighted Circuit Masking}
\label{app:framework}
Weighted circuit masking fixes a problem, \labelcref{eq:p2}, and a way of relaxing it, but leaves most of the recipe open.
\Cref{tab:framework} lists its slots: what is reproduced, over which nodes, with which parametrisation, under which ablation methods and backgrounds, with which regularisers and optimiser, from which contexts, and how the result is scored.
The production choice in each slot was made by an experiment reported in these appendices, most of them on TuringLLM and carried over unchanged to the public models (\cref{tab:protocol}, Provenance); none is claimed to be the best setting, and several are known to trade one evaluation against another.

{\footnotesize
\setlength{\tabcolsep}{3pt}
\renewcommand{\arraystretch}{1.15}
\begin{longtable}{@{}p{0.13\linewidth}p{0.33\linewidth}p{0.33\linewidth}p{0.16\linewidth}@{}}
\caption{The slots of weighted circuit masking, the production choice in each, and the alternatives this paper examined or that the formulation admits.}
\label{tab:framework} \\
\toprule
\textbf{Slot} & \textbf{Production choice} & \textbf{Alternatives} & \textbf{Evidence} \\
\midrule
\endfirsthead
\toprule
\textbf{Slot} & \textbf{Production choice} & \textbf{Alternatives} & \textbf{Evidence} \\
\midrule
\endhead
\midrule
\multicolumn{4}{r@{}}{\emph{continued on the next page}} \\
\endfoot
\bottomrule
\endlastfoot
Target & the pre-activation $\preact$ of one SAE latent at its anchors; censored read under Z, C, A & an output metric, as in learned-mask circuits for outputs~\citep{bhaskar2024finding, caples2025scaling}; output- and latent-targeted circuits fail each other's objective (exploratory) & \cref{sec:ablmask}, \cref{app:secondary,app:logit-faith} \\
Nodes & every latent at every SAE site upstream of the target; SAE errors passed through; nodes recomputed & dense ReLU dictionaries (the method replicates; the random baseline needs its own design); members clamped (resample diagnostic) & \cref{sec:problem}, \cref{app:dense-relu,app:resample} \\
Parametrisation & gate $\sigma(\theta/T)$ and coefficient $\mathrm{softplus}(\psi)$ per node, $\alpha_0 = 1$ & $\alpha \equiv 1$ (unweighted circuit masking) & \cref{sec:weighted}, \cref{tab:run2026-alpha1} \\
Training ablations & Z (weight $1$), C and A ($0.25$ each) under evaluation's sparsity-preserving fill; resample D over the training contrast contexts ($1$) & one method alone; dense means; D in place of C; no D & \cref{app:single-method,app:respect-topk,app:resample} \\
Contrast background & $64$ close, silence-checked contexts per target ($48$ train, $16$ held out) & random, distant, or fused sources; the unchecked retrieval store & \cref{app:stores,app:contrast-source} \\
Regularisers & flat penalty $\lambda = 10^{-3}$ per open gate; closed-gate charge $\lambda(1-\sigma(\theta))|\alpha - 1|$; rank-keep $\gamma_{\mathrm{K}} = 10^{-1}$; one shared normaliser $\nu$ & other penalties $\lambda$ (sweep); off-target term, Top-K hinge, pull to rank $1$, firing margin; per-term normalisers & \cref{sec:lambda-transfer}, \cref{app:loss-details,app:rank-designs,app:membership-studies} \\
Optimiser & AdamW, $400$ steps, batch $4$; gate temperature annealed $1 \to 0.05$; learning-rate warm-up, then geometric decay & straight-through binarisation; $200$-step and reach-then-hold schedules; larger batches (exploratory) & \cref{tab:protocol}, \cref{app:membership-studies,app:framework-history} \\
Training contexts & $32$ strongest and $16$ mid-band activating contexts, deduplicated and split by whole groups of copies & a split without deduplication, which leaked copies into held-out sets & \cref{sec:setting}, \cref{app:heldout} \\
Evaluation & held-out strongest contexts; Z, A, C in $[0.75, 1.5]$ and necessity $\geq 0.9$; vacuous targets excluded and counted & reported, not required: mid-band contexts, sufficiency to induce, donor-background test, specificity, $\alpha = 1$ rescoring, random-circuit baselines, stable core & \cref{sec:eval}, \cref{app:protocol,app:nulls,app:underdetermination} \\
\end{longtable}}

\paragraph{Training in the evaluation's domain.}
Across these experiments, circuits did best on the evaluations their training resembled.
When a fit fails an evaluation, the failure is often a property of the recipe, and a slot of \cref{tab:framework} can put the fit in that evaluation's domain.
% app:resample (095: 72 targets, 0% -> 24%, CI +14 to +33; criterion 69% -> 62%, CI -17 to +3; 1.18x size, 1.45x time;
% D in place of C: 51%, CI -29 to -7), production paragraph (exp 102: 19 -> 39 and 48 -> 54 of 72, 1.33x size), T10
% (tab:run2026-tl-noR: +5.2, -0.6 to 11.6).
\emph{Real backgrounds.} Circuits trained only against zero and means failed a resample test almost always; adding a resample term over the training contrast contexts raised the resample pass rate on $72$ targets from $0\%$ to $24\%$, at $1.18\times$ the circuit size and $1.45\times$ the fitting time, and with rank-keep raised it took a $72$-target check panel from $19$ to $39$ targets passing (\cref{app:resample}).
Using it in place of the contrast-mean term gained as much on the resample test but cost the pass criterion, so terms should be added for a new evaluation rather than swapped for an old one; on the production sweep the term's effect on the pass criterion itself is $5.2$ points and not significant (\cref{tab:run2026-tl-noR}).
\emph{Ablation methods and fills.} A mask trained under one ablation method learns what that method rewards: without the zero-ablation term, coefficients overshoot zero ablation by orders of magnitude, and circuits trained against dense means can fail under the sparsity-preserving fill that evaluation uses (\cref{app:single-method,app:respect-topk}).
\emph{Top-K survival.} Reproducing the pre-activation does not keep the target in its site's Top-K; the censored read and rank-keep train for it, though neither removes amplifier circuits (\cref{app:loss-details,app:specificity}).
\emph{Other backgrounds.} Contrast faithfulness generalises furthest when training uses random or distant contrast contexts (\cref{app:contrast-source}).
\emph{Size.} Circuits shrink with $\lambda$ at a nearly constant pass rate up to a plateau's edge (\cref{sec:lambda-transfer}).
Two evaluations have no training term yet: sufficiency to induce, which no circuit is trained on (median $0.78$--$0.88$ on held-out contrast contexts, $0.35$ for TuringLLM attention outputs; \cref{app:production}), and transfer to weaker, mid-band contexts (\cref{sec:main-result}).
The framework's prescription for them is a term or a training set in their domain, which we have not tested.
% CANDIDATE EVIDENCE, NOT IN THE PAPER (main.tex line 503: the arena results are reported separately). If Daniel wants
% them here, each needs its own sentence and its own appendix or citation:
%  - arena A tuning (experiments/107-true-baselines/arena_a/TUNING.md): on GPT-2 transcoder + z-SAE substrate the
%    necessity hinge (NEC_W; off in production, docs/wcm-specification.md §2.3) was chosen at 0.3 by a pre-registered
%    9-target rule; NEC_W 1.0 raised nec_a (3.tc.3752 0.32 -> 0.95) but pushed A out of band (1.54). The censored read
%    is inert there by construction (k = d_sae).
%  - arena A scope a1 (z + z2 + tc) vs a2 (tc only) and the wikitext-smoke vs OpenWebText corpus test (DAN-208 text;
%    not traced to an output file in this pass).
%  - logit / maximise / behavioural objectives (arena C/D; memory notes say "don't quote yet").
A fit trained for an evaluation is, for that evaluation, scored on home turf (\cref{app:method-comparison}), so the evaluation stays meaningful only if its contexts and donors are never seen in training, as the donor-background test's donors are not (\cref{tab:protocol}), and the evaluations not trained for are reported alongside.

\paragraph{Choosing the sparsity penalty.}
% sec:lambda-transfer, tab:run2026-tl-sweep, tab:run2026-lambda-pilot.
The penalty sets circuit size and, past a point, the pass rate: on the TuringLLM sweep the pass rate holds at $89$--$93\%$ from $\lambda = 2.5 \times 10^{-5}$ to $5 \times 10^{-4}$ while each doubling of $\lambda$ shrinks circuits $1.4$--$1.9\times$, and it declines beyond (\cref{sec:lambda-transfer}).
We recommend fixing the selection rule before looking at pass rates, as our size-only pilot rule did on the public models, where it kept $\lambda = 10^{-3}$ on both (\cref{tab:run2026-lambda-pilot}), and reporting a curve over $\lambda$ rather than one point, since one setting that transfers is not the best setting for each model.

\paragraph{A new model or dictionary.}
Only the model and dictionary adapters change between our three models (\cref{sec:setting}); four things needed checking.
\emph{Precision}: TuringLLM runs in the bfloat16 its pipeline was validated in, but on the public models a bfloat16 data pass did not reproduce the float32 top-context sets, so they run in float32 (\cref{tab:protocol}); a precision check belongs before any fit, since node sets also move with rounding (\cref{app:underdetermination}).
\emph{The dictionary's sparsity}: $k$ sets the fill's budget and the censored read's cut; on a dictionary that keeps every latent ($k = d_{\mathrm{sae}}$) the censored read has no cut to model, and on dense ReLU dictionaries a random-circuit baseline with fitted coefficients stops discriminating, so necessity, not faithfulness, separates real from random circuits there (\cref{app:dense-relu}).
% The k = d_sae statement follows from the definition of tau^(K) (no K-th other latent); docs/wcm-specification.md §7.6.
\emph{Contrast retrieval}: contrast contexts need a pooled representation for every corpus sequence and a silence check, which rejects $14$--$51\%$ of the closest candidates on the targets we inspected (\cref{app:stores}).
\emph{Corpus copies}: duplicated windows silently leak into held-out sets unless pools are deduplicated and split by groups of copies (\cref{app:heldout}).

\paragraph{What to report.}
Because each slot changes the circuits, a result is interpretable only with its recipe: the ablation methods used in training and in evaluation; faithfulness under each method and necessity, on held-out contexts, by depth band and site kind, with intervals over targets; excluded and vacuous targets counted; circuit size with whether coefficients were fitted, and the $\alpha = 1$ rescoring; random-circuit baselines given the same fitting; and the reported-only evaluations (sufficiency to induce, the donor-background test, specificity, mid-band transfer) beside the pass rate (\cref{sec:discussion}).

\paragraph{Choices that matter, and what stays open.}
The ablation method changes not only scores but rankings: the two mean-ablation variants can order small node sets differently (\cref{app:respect-topk}), and in exploratory comparisons a method tended to win the scores of the ablation method it was trained or ranked under (\cref{app:method-comparison}), so comparisons between recipes need the full matrix of recipe against ablation method.
Node sets are not identified by any recipe: a refit with reordered training contexts shares a median Jaccard of $0.39$--$0.51$ with the original and its shared core alone passes for $44$--$72\%$ of targets (\cref{app:underdetermination}), so changes to a slot should be judged by sizes and scores, not by which latents moved.
Finally, the production choices were calibrated on TuringLLM targets that overlap its evaluation set (\cref{sec:setting}), and on the public models every setting except $\lambda$ carried over unchanged (\cref{tab:protocol}, Provenance), so on a new model they are defaults to be checked, not tuned values.

% ---------------------------------------------------------------------------------------------------------------------
% 2026-11-08 (Daniel, via coordinator): the development path as evidence for where WCM can go next. Every number is
% traced; sources in the comment before each step. Panels and protocols differ between steps, so numbers are compared
% only within a step (each step is a paired test against the recipe before it).
% ---------------------------------------------------------------------------------------------------------------------
\subsection{How the Production Recipe Was Reached}
\label{app:framework-history}
The production recipe is the end of a sequence of changes, each made because a run exposed a failure, and each tested against the recipe before it on the same targets, mostly under a selection rule stated before the comparison.
We list them in order as evidence of which kinds of change have paid off; the panels differ between steps, so numbers compare only within a step.

% app:stores (38% fallback), app:contrast-source; experiments/055-close-contrast, 071-contrast-sweep.
\paragraph{Contrast contexts built per target.}
The first contrast contexts came from a retrieval store whose index covered too few sequences, so $38\%$ of TuringLLM targets fell back to one shared list; WCM now selects contrast contexts per target at fit time, ranked by similarity to its strongest contexts and kept only if the target is silent (\cref{app:stores}), and a later comparison of contrast sources showed how much that choice shapes the C term (\cref{app:contrast-source}).
The lesson: inspect the data each training term actually sees.

% app:specificity, app:rank-designs; experiments/056-specificity.
\paragraph{Rank-keep against off-target lifts.}
Reproducing the pre-activation left some circuits restoring their target by lifting its whole neighbourhood; of the terms tried, rank-keep alone kept the gains without the alternatives' costs, and entered the first protocol at a small weight (\cref{app:rank-designs,app:specificity}).

% app:v1-run (15,046 targets, 46%, attention 14%, layers 8-11 22%); app:heldout (66.8%); app:resample (0.1% vs 43.1%,
% 747 targets); fig:failure-anatomy.
\paragraph{The first run at scale.}
Under that protocol, WCM was fitted to $15{,}046$ TuringLLM targets, and $46\%$ passed, $22\%$ at layers 8--11 and $14\%$ for attention outputs (\cref{app:v1-run}).
Analyses of the run exposed three problems: $66.8\%$ of held-out strongest contexts repeated a training context, which led to deduplicated records split by groups of copies (\cref{app:heldout}); with depth, failing circuits increasingly failed only the two mean ablations (\cref{fig:failure-anatomy}); and $0.1\%$ of circuits passed a resample test, against $43.1\%$ under the pass criterion (\cref{app:resample}).

% experiments/093-v2-validation/README.md (out/tables.md): 72 targets, 8 per band x kind, paired, pass floor 0.8, the
% v1 run's as-run held-out sets (leaky; same for both arms): 31 -> 50 of 72, +26.4 pp [15.3, 38.9]; L8-11 2 -> 11 of 24;
% attn 3 -> 12 of 24; median dZ 0.00; size ratio 0.91; fit time 1.27x. Pilot: experiments/090-sparse-fill-training.
% Also main.tex comment at app:respect-topk ("The same fill in training").
\paragraph{Training under the fill that evaluation uses.}
The mean-ablation failures were a mismatch: training replaced non-members with dense means, a code no Top-K dictionary emits, while evaluation re-imposed the Top-K (\cref{app:respect-topk}).
Training the A and C terms under evaluation's own fill and means raised the pass rate on $72$ stratified targets from $31$ to $50$ (paired difference $+26.4$ points, $95\%$ interval $15.3$--$38.9$), from $2$ to $11$ of $24$ at layers 8--11 and from $3$ to $12$ of $24$ for attention outputs, with zero-ablation faithfulness unchanged (median paired change $0.00$), circuits $0.91\times$ the size, and $1.27\times$ the fitting time.
The gain came from aligning training with evaluation, not from a new term.

% app:resample (095 numbers and the production paragraph: 102's 19 -> 39, 48 -> 54 of 72, 1.33x size);
% experiments/102-pass-improvement/README.md line 89 (diagnosis: 15 of 18 Z failures of the previous version restored
% the pre-activation and lost at the Top-K cut); tab:run2026-tl-noR.
\paragraph{A resample term, and a stronger rank-keep.}
To survive real backgrounds, training gained the resample term (above and \cref{app:resample}).
On the previous version's zero-ablation failures, $15$ of $18$ restored the target's pre-activation but lost it at the Top-K cut, so rank-keep was raised from $3 \times 10^{-3}$ to $10^{-1}$ with the resample term at full weight; chosen by a rule stated before the comparison, the pair took a $72$-target check panel from $48$ to $54$ targets passing and from $19$ to $39$ passing the resample test, at $1.33\times$ the circuit size (\cref{app:resample}).
In the production sweep, the arm trained without the resample term passes $5.2$ points less, not significantly (\cref{tab:run2026-tl-noR}).

% experiments/107-true-baselines/latent_pilot/README.md (9 targets enriched for failures: 3 -> 5 of 9, size ratio ~1.04;
% plain geometric: circuits 2-4x larger, sum(lr) 6.3 vs 20; constant lr makes lr/T grow 20x) and
% experiments/107-true-baselines/adaptive_lambda/README.md, "Check" table (arm "geo x 3.15 + censored term
% (latent_pilot)": 63 of 72 vs 54 for the previous version, +10 / -1, size ratio 1.04 on pass-in-both). VERIFY: whether
% that check arm included the 0.1 warm-up (run_tonight.sh stage 4 fits check_geo_x3.15_warmup_censor).
\paragraph{A learning-rate schedule and the censored read.}
Under a constant learning rate, the gate step grows twentyfold as the gate temperature anneals, and some fits degraded late in training.
A geometric decay alone stabilised them but left gates too little total step to close, so circuits grew $2$--$4\times$; raising its peak $3.15\times$ restored the summed step size, and a short warm-up was added for the early spikes the higher peak brought back (\cref{app:loss-details}).
With the censored read added, so that training also sees whether the target survives its site's Top-K, the check panel went from $54$ to $63$ of $72$ targets passing ($10$ gained, $1$ lost) at about the same circuit size (median ratio $1.04$).

% sec:lambda-transfer, tab:run2026-lambda-pilot; experiments/108-batch-size/out/analysis.md (16 TuringLLM targets;
% B16 at 400 steps: 15 / 16 vs 16 / 16, size 0.75x on pass-in-both, 99.9 vs 29.1 fit-minutes = 3.4x; pre-stated rule:
% +-1 target is noise); experiments/107-true-baselines/adaptive_lambda/README.md (every shrinking policy traded passes
% for size on the check panel; losses were necessity failures).
\paragraph{Settings left as they were.}
The penalty was carried to the public models by a size-only rule fixed before a pilot (\cref{sec:lambda-transfer}).
In exploratory tests, batches of $16$ for the same $400$ steps passed $15$ of $16$ TuringLLM targets against $16$, with circuits $0.75\times$ the size where both passed, at $3.4\times$ the fitting time, so the batch of $4$ was kept; and policies that adapted $\lambda$ per target traded passes for size on the check panel, losing passes mostly to necessity, which no training term covers.
Other rejected designs are in \cref{app:membership-studies}.

% Directions: app:production (induce), sec:main-result (mid-band, deep band, attention), app:specificity (amplifiers,
% v1: 30% flagged), app:rank-designs (off-target term: lifts cut 3-7x, but grew circuits), app:resample (anchors
% only; overfiring away from the anchor), adaptive_lambda README (necessity).
\paragraph{Directions not yet taken.}
The same path suggests where the recipe can go next.
Necessity, sufficiency to induce, and transfer to mid-band contexts are scored but not trained for; each could get a term or training contexts in its domain, as the resample test did.
Deep and attention-output targets still pass least (\cref{sec:main-result}), and in the schedule pilot the attention failures kept zero-ablation faithfulness near $0$ under every schedule, so the learning rate does not explain them.
Amplifier circuits survive rank-keep (\cref{app:specificity}); the off-target term cut lifts three- to sevenfold but grew circuits (\cref{app:rank-designs}), so a cheaper specificity term remains open.
And circuits explain the target at its anchors, not at every position (\cref{app:resample}); training on more positions is untested.
\fi\fi
        % ~1/3 page, Discussion (replaces the paragraph
%                                                                    % "Ablation methods are part of the method")
%   \def\frameworkpart{appendix}% =====================================================================================================================
% DAN-208 (drafted 2026-11-08, agent "framework"): WCM as a framework for solving for circuits, not a fixed algorithm.
% Two versions in one file. Choose with \def\frameworkpart{main} or \def\frameworkpart{appendix} before \input:
%
%   \def\frameworkpart{main}% =====================================================================================================================
% DAN-208 (drafted 2026-11-08, agent "framework"): WCM as a framework for solving for circuits, not a fixed algorithm.
% Two versions in one file. Choose with \def\frameworkpart{main} or \def\frameworkpart{appendix} before \input:
%
%   \def\frameworkpart{main}\input{drafts/framework_section}        % ~1/3 page, Discussion (replaces the paragraph
%                                                                    % "Ablation methods are part of the method")
%   \def\frameworkpart{appendix}\input{drafts/framework_section}    % ~1.3 pages, its own appendix section
%
% New labels: sec:framework (main), app:framework, tab:framework (appendix). Every other \cref target exists in
% main.tex as of 2026-11-08. Every number below is already printed in main.tex (source comments give the place);
% the one exception, experiment 108, is marked NOT-YET-IN-PAPER where it appears.
% Vocabulary: paper/notes/vocabulary.md (no "endpoint", "floor", "drivers"; "fill" only as main.tex already uses it).
% Never-claim list respected: no "first"; nothing here claims novelty.
% =====================================================================================================================
\providecommand{\frameworkpart}{main}
\def\frameworkmain{main}
\def\frameworkappendix{appendix}
\ifx\frameworkpart\frameworkmain
% ---------------------------------------------------------------------------------------------------------------------
% (1) MAIN-TEXT VERSION. Proposed place: sec:discussion, in place of the paragraph "Ablation methods are part of the
% method" (main.tex lines 548-552), whose evidence it repeats and extends; the paragraph "Recommendations" (lines
% 564-566) can then point here. If it must sit in Method instead, it fits after "Pass criterion" (line 341), before
% \section{Experiments and Results}, but there it would cite results not yet shown.
% Numbers: 2-14% = abstract / run2026_alpha1 (sec:weighted); "one to three orders of magnitude" = app:single-method
% (exploratory); contrast sources = app:contrast-source (protocol v1); resample = app:resample.
% ---------------------------------------------------------------------------------------------------------------------
\subsection{Using WCM as a Framework}
\label{sec:framework}
WCM is better read as a framework for solving for circuits than as one fixed algorithm.
Its parts are interchangeable: an objective (the quantity reproduced, the ablation methods it is reproduced under, the regularisers, and their weights; \cref{eq:loss}), a parametrisation (a gate and a coefficient per node), data (training, contrast, and held-out contexts), and an optimiser.
The production recipe is one setting of these parts, each chosen by an experiment (\cref{tab:framework}), and changing a part changes the circuits found.
With every coefficient fixed at $1$ the same nodes pass for $2$--$14\%$ of targets (\cref{sec:weighted}); in exploratory fits, coefficients trained without zero ablation overshoot it by one to three orders of magnitude (\cref{app:single-method}); circuits trained against close contrast contexts are the most faithful on close backgrounds and the least on any other (\cref{app:contrast-source}); and circuits failed a resample test until training included a resample term (\cref{app:resample}).
When WCM fails an evaluation, we therefore suggest first asking whether training ever saw that evaluation's domain, and if not, adding the term, contexts, or setting that puts it there, while keeping the evaluation's own contexts and donors out of training so that it stays a held-out test (\cref{app:framework}).
A pass rate then describes a recipe, a dictionary, and an evaluation together, not WCM alone, which is why the ablation methods used for training and evaluation belong in the method statement~\citep{mueller2025mib, mahale2026explanation}.
%
\else\ifx\frameworkpart\frameworkappendix
% ---------------------------------------------------------------------------------------------------------------------
% (2) APPENDIX VERSION. Proposed place: a new appendix section directly before "Weighted Circuit Masking Design
% Ablations" (main.tex line 2003, app:membership-studies), so that the design ablations it summarises follow it; or
% directly after app:protocol / app:loss-details (before line 853) if Daniel prefers it next to tab:protocol.
% ---------------------------------------------------------------------------------------------------------------------
\section{Building with Weighted Circuit Masking}
\label{app:framework}
Weighted circuit masking fixes a problem, \labelcref{eq:p2}, and a way of relaxing it, but leaves most of the recipe open.
\Cref{tab:framework} lists its slots: what is reproduced, over which nodes, with which parametrisation, under which ablation methods and backgrounds, with which regularisers and optimiser, from which contexts, and how the result is scored.
The production choice in each slot was made by an experiment reported in these appendices, most of them on TuringLLM and carried over unchanged to the public models (\cref{tab:protocol}, Provenance); none is claimed to be the best setting, and several are known to trade one evaluation against another.

{\footnotesize
\setlength{\tabcolsep}{3pt}
\renewcommand{\arraystretch}{1.15}
\begin{longtable}{@{}p{0.13\linewidth}p{0.33\linewidth}p{0.33\linewidth}p{0.16\linewidth}@{}}
\caption{The slots of weighted circuit masking, the production choice in each, and the alternatives this paper examined or that the formulation admits.}
\label{tab:framework} \\
\toprule
\textbf{Slot} & \textbf{Production choice} & \textbf{Alternatives} & \textbf{Evidence} \\
\midrule
\endfirsthead
\toprule
\textbf{Slot} & \textbf{Production choice} & \textbf{Alternatives} & \textbf{Evidence} \\
\midrule
\endhead
\midrule
\multicolumn{4}{r@{}}{\emph{continued on the next page}} \\
\endfoot
\bottomrule
\endlastfoot
Target & the pre-activation $\preact$ of one SAE latent at its anchors; censored read under Z, C, A & an output metric, as in learned-mask circuits for outputs~\citep{bhaskar2024finding, caples2025scaling}; output- and latent-targeted circuits fail each other's objective (exploratory) & \cref{sec:ablmask}, \cref{app:secondary,app:logit-faith} \\
Nodes & every latent at every SAE site upstream of the target; SAE errors passed through; nodes recomputed & dense ReLU dictionaries (the method replicates; the random baseline needs its own design); members clamped (resample diagnostic) & \cref{sec:problem}, \cref{app:dense-relu,app:resample} \\
Parametrisation & gate $\sigma(\theta/T)$ and coefficient $\mathrm{softplus}(\psi)$ per node, $\alpha_0 = 1$ & $\alpha \equiv 1$ (unweighted circuit masking) & \cref{sec:weighted}, \cref{tab:run2026-alpha1} \\
Training ablations & Z (weight $1$), C and A ($0.25$ each) under evaluation's sparsity-preserving fill; resample D over the training contrast contexts ($1$) & one method alone; dense means; D in place of C; no D & \cref{app:single-method,app:respect-topk,app:resample} \\
Contrast background & $64$ close, silence-checked contexts per target ($48$ train, $16$ held out) & random, distant, or fused sources; the unchecked retrieval store & \cref{app:stores,app:contrast-source} \\
Regularisers & flat penalty $\lambda = 10^{-3}$ per open gate; closed-gate charge $\lambda(1-\sigma(\theta))|\alpha - 1|$; rank-keep $\gamma_{\mathrm{K}} = 10^{-1}$; one shared normaliser $\nu$ & other penalties $\lambda$ (sweep); off-target term, Top-K hinge, pull to rank $1$, firing margin; per-term normalisers & \cref{sec:lambda-transfer}, \cref{app:loss-details,app:rank-designs,app:membership-studies} \\
Optimiser & AdamW, $400$ steps, batch $4$; gate temperature annealed $1 \to 0.05$; learning-rate warm-up, then geometric decay & straight-through binarisation; $200$-step and reach-then-hold schedules; larger batches (exploratory) & \cref{tab:protocol}, \cref{app:membership-studies,app:framework-history} \\
Training contexts & $32$ strongest and $16$ mid-band activating contexts, deduplicated and split by whole groups of copies & a split without deduplication, which leaked copies into held-out sets & \cref{sec:setting}, \cref{app:heldout} \\
Evaluation & held-out strongest contexts; Z, A, C in $[0.75, 1.5]$ and necessity $\geq 0.9$; vacuous targets excluded and counted & reported, not required: mid-band contexts, sufficiency to induce, donor-background test, specificity, $\alpha = 1$ rescoring, random-circuit baselines, stable core & \cref{sec:eval}, \cref{app:protocol,app:nulls,app:underdetermination} \\
\end{longtable}}

\paragraph{Training in the evaluation's domain.}
Across these experiments, circuits did best on the evaluations their training resembled.
When a fit fails an evaluation, the failure is often a property of the recipe, and a slot of \cref{tab:framework} can put the fit in that evaluation's domain.
% app:resample (095: 72 targets, 0% -> 24%, CI +14 to +33; criterion 69% -> 62%, CI -17 to +3; 1.18x size, 1.45x time;
% D in place of C: 51%, CI -29 to -7), production paragraph (exp 102: 19 -> 39 and 48 -> 54 of 72, 1.33x size), T10
% (tab:run2026-tl-noR: +5.2, -0.6 to 11.6).
\emph{Real backgrounds.} Circuits trained only against zero and means failed a resample test almost always; adding a resample term over the training contrast contexts raised the resample pass rate on $72$ targets from $0\%$ to $24\%$, at $1.18\times$ the circuit size and $1.45\times$ the fitting time, and with rank-keep raised it took a $72$-target check panel from $19$ to $39$ targets passing (\cref{app:resample}).
Using it in place of the contrast-mean term gained as much on the resample test but cost the pass criterion, so terms should be added for a new evaluation rather than swapped for an old one; on the production sweep the term's effect on the pass criterion itself is $5.2$ points and not significant (\cref{tab:run2026-tl-noR}).
\emph{Ablation methods and fills.} A mask trained under one ablation method learns what that method rewards: without the zero-ablation term, coefficients overshoot zero ablation by orders of magnitude, and circuits trained against dense means can fail under the sparsity-preserving fill that evaluation uses (\cref{app:single-method,app:respect-topk}).
\emph{Top-K survival.} Reproducing the pre-activation does not keep the target in its site's Top-K; the censored read and rank-keep train for it, though neither removes amplifier circuits (\cref{app:loss-details,app:specificity}).
\emph{Other backgrounds.} Contrast faithfulness generalises furthest when training uses random or distant contrast contexts (\cref{app:contrast-source}).
\emph{Size.} Circuits shrink with $\lambda$ at a nearly constant pass rate up to a plateau's edge (\cref{sec:lambda-transfer}).
Two evaluations have no training term yet: sufficiency to induce, which no circuit is trained on (median $0.78$--$0.88$ on held-out contrast contexts, $0.35$ for TuringLLM attention outputs; \cref{app:production}), and transfer to weaker, mid-band contexts (\cref{sec:main-result}).
The framework's prescription for them is a term or a training set in their domain, which we have not tested.
% CANDIDATE EVIDENCE, NOT IN THE PAPER (main.tex line 503: the arena results are reported separately). If Daniel wants
% them here, each needs its own sentence and its own appendix or citation:
%  - arena A tuning (experiments/107-true-baselines/arena_a/TUNING.md): on GPT-2 transcoder + z-SAE substrate the
%    necessity hinge (NEC_W; off in production, docs/wcm-specification.md §2.3) was chosen at 0.3 by a pre-registered
%    9-target rule; NEC_W 1.0 raised nec_a (3.tc.3752 0.32 -> 0.95) but pushed A out of band (1.54). The censored read
%    is inert there by construction (k = d_sae).
%  - arena A scope a1 (z + z2 + tc) vs a2 (tc only) and the wikitext-smoke vs OpenWebText corpus test (DAN-208 text;
%    not traced to an output file in this pass).
%  - logit / maximise / behavioural objectives (arena C/D; memory notes say "don't quote yet").
A fit trained for an evaluation is, for that evaluation, scored on home turf (\cref{app:method-comparison}), so the evaluation stays meaningful only if its contexts and donors are never seen in training, as the donor-background test's donors are not (\cref{tab:protocol}), and the evaluations not trained for are reported alongside.

\paragraph{Choosing the sparsity penalty.}
% sec:lambda-transfer, tab:run2026-tl-sweep, tab:run2026-lambda-pilot.
The penalty sets circuit size and, past a point, the pass rate: on the TuringLLM sweep the pass rate holds at $89$--$93\%$ from $\lambda = 2.5 \times 10^{-5}$ to $5 \times 10^{-4}$ while each doubling of $\lambda$ shrinks circuits $1.4$--$1.9\times$, and it declines beyond (\cref{sec:lambda-transfer}).
We recommend fixing the selection rule before looking at pass rates, as our size-only pilot rule did on the public models, where it kept $\lambda = 10^{-3}$ on both (\cref{tab:run2026-lambda-pilot}), and reporting a curve over $\lambda$ rather than one point, since one setting that transfers is not the best setting for each model.

\paragraph{A new model or dictionary.}
Only the model and dictionary adapters change between our three models (\cref{sec:setting}); four things needed checking.
\emph{Precision}: TuringLLM runs in the bfloat16 its pipeline was validated in, but on the public models a bfloat16 data pass did not reproduce the float32 top-context sets, so they run in float32 (\cref{tab:protocol}); a precision check belongs before any fit, since node sets also move with rounding (\cref{app:underdetermination}).
\emph{The dictionary's sparsity}: $k$ sets the fill's budget and the censored read's cut; on a dictionary that keeps every latent ($k = d_{\mathrm{sae}}$) the censored read has no cut to model, and on dense ReLU dictionaries a random-circuit baseline with fitted coefficients stops discriminating, so necessity, not faithfulness, separates real from random circuits there (\cref{app:dense-relu}).
% The k = d_sae statement follows from the definition of tau^(K) (no K-th other latent); docs/wcm-specification.md §7.6.
\emph{Contrast retrieval}: contrast contexts need a pooled representation for every corpus sequence and a silence check, which rejects $14$--$51\%$ of the closest candidates on the targets we inspected (\cref{app:stores}).
\emph{Corpus copies}: duplicated windows silently leak into held-out sets unless pools are deduplicated and split by groups of copies (\cref{app:heldout}).

\paragraph{What to report.}
Because each slot changes the circuits, a result is interpretable only with its recipe: the ablation methods used in training and in evaluation; faithfulness under each method and necessity, on held-out contexts, by depth band and site kind, with intervals over targets; excluded and vacuous targets counted; circuit size with whether coefficients were fitted, and the $\alpha = 1$ rescoring; random-circuit baselines given the same fitting; and the reported-only evaluations (sufficiency to induce, the donor-background test, specificity, mid-band transfer) beside the pass rate (\cref{sec:discussion}).

\paragraph{Choices that matter, and what stays open.}
The ablation method changes not only scores but rankings: the two mean-ablation variants can order small node sets differently (\cref{app:respect-topk}), and in exploratory comparisons a method tended to win the scores of the ablation method it was trained or ranked under (\cref{app:method-comparison}), so comparisons between recipes need the full matrix of recipe against ablation method.
Node sets are not identified by any recipe: a refit with reordered training contexts shares a median Jaccard of $0.39$--$0.51$ with the original and its shared core alone passes for $44$--$72\%$ of targets (\cref{app:underdetermination}), so changes to a slot should be judged by sizes and scores, not by which latents moved.
Finally, the production choices were calibrated on TuringLLM targets that overlap its evaluation set (\cref{sec:setting}), and on the public models every setting except $\lambda$ carried over unchanged (\cref{tab:protocol}, Provenance), so on a new model they are defaults to be checked, not tuned values.

% ---------------------------------------------------------------------------------------------------------------------
% 2026-11-08 (Daniel, via coordinator): the development path as evidence for where WCM can go next. Every number is
% traced; sources in the comment before each step. Panels and protocols differ between steps, so numbers are compared
% only within a step (each step is a paired test against the recipe before it).
% ---------------------------------------------------------------------------------------------------------------------
\subsection{How the Production Recipe Was Reached}
\label{app:framework-history}
The production recipe is the end of a sequence of changes, each made because a run exposed a failure, and each tested against the recipe before it on the same targets, mostly under a selection rule stated before the comparison.
We list them in order as evidence of which kinds of change have paid off; the panels differ between steps, so numbers compare only within a step.

% app:stores (38% fallback), app:contrast-source; experiments/055-close-contrast, 071-contrast-sweep.
\paragraph{Contrast contexts built per target.}
The first contrast contexts came from a retrieval store whose index covered too few sequences, so $38\%$ of TuringLLM targets fell back to one shared list; WCM now selects contrast contexts per target at fit time, ranked by similarity to its strongest contexts and kept only if the target is silent (\cref{app:stores}), and a later comparison of contrast sources showed how much that choice shapes the C term (\cref{app:contrast-source}).
The lesson: inspect the data each training term actually sees.

% app:specificity, app:rank-designs; experiments/056-specificity.
\paragraph{Rank-keep against off-target lifts.}
Reproducing the pre-activation left some circuits restoring their target by lifting its whole neighbourhood; of the terms tried, rank-keep alone kept the gains without the alternatives' costs, and entered the first protocol at a small weight (\cref{app:rank-designs,app:specificity}).

% app:v1-run (15,046 targets, 46%, attention 14%, layers 8-11 22%); app:heldout (66.8%); app:resample (0.1% vs 43.1%,
% 747 targets); fig:failure-anatomy.
\paragraph{The first run at scale.}
Under that protocol, WCM was fitted to $15{,}046$ TuringLLM targets, and $46\%$ passed, $22\%$ at layers 8--11 and $14\%$ for attention outputs (\cref{app:v1-run}).
Analyses of the run exposed three problems: $66.8\%$ of held-out strongest contexts repeated a training context, which led to deduplicated records split by groups of copies (\cref{app:heldout}); with depth, failing circuits increasingly failed only the two mean ablations (\cref{fig:failure-anatomy}); and $0.1\%$ of circuits passed a resample test, against $43.1\%$ under the pass criterion (\cref{app:resample}).

% experiments/093-v2-validation/README.md (out/tables.md): 72 targets, 8 per band x kind, paired, pass floor 0.8, the
% v1 run's as-run held-out sets (leaky; same for both arms): 31 -> 50 of 72, +26.4 pp [15.3, 38.9]; L8-11 2 -> 11 of 24;
% attn 3 -> 12 of 24; median dZ 0.00; size ratio 0.91; fit time 1.27x. Pilot: experiments/090-sparse-fill-training.
% Also main.tex comment at app:respect-topk ("The same fill in training").
\paragraph{Training under the fill that evaluation uses.}
The mean-ablation failures were a mismatch: training replaced non-members with dense means, a code no Top-K dictionary emits, while evaluation re-imposed the Top-K (\cref{app:respect-topk}).
Training the A and C terms under evaluation's own fill and means raised the pass rate on $72$ stratified targets from $31$ to $50$ (paired difference $+26.4$ points, $95\%$ interval $15.3$--$38.9$), from $2$ to $11$ of $24$ at layers 8--11 and from $3$ to $12$ of $24$ for attention outputs, with zero-ablation faithfulness unchanged (median paired change $0.00$), circuits $0.91\times$ the size, and $1.27\times$ the fitting time.
The gain came from aligning training with evaluation, not from a new term.

% app:resample (095 numbers and the production paragraph: 102's 19 -> 39, 48 -> 54 of 72, 1.33x size);
% experiments/102-pass-improvement/README.md line 89 (diagnosis: 15 of 18 Z failures of the previous version restored
% the pre-activation and lost at the Top-K cut); tab:run2026-tl-noR.
\paragraph{A resample term, and a stronger rank-keep.}
To survive real backgrounds, training gained the resample term (above and \cref{app:resample}).
On the previous version's zero-ablation failures, $15$ of $18$ restored the target's pre-activation but lost it at the Top-K cut, so rank-keep was raised from $3 \times 10^{-3}$ to $10^{-1}$ with the resample term at full weight; chosen by a rule stated before the comparison, the pair took a $72$-target check panel from $48$ to $54$ targets passing and from $19$ to $39$ passing the resample test, at $1.33\times$ the circuit size (\cref{app:resample}).
In the production sweep, the arm trained without the resample term passes $5.2$ points less, not significantly (\cref{tab:run2026-tl-noR}).

% experiments/107-true-baselines/latent_pilot/README.md (9 targets enriched for failures: 3 -> 5 of 9, size ratio ~1.04;
% plain geometric: circuits 2-4x larger, sum(lr) 6.3 vs 20; constant lr makes lr/T grow 20x) and
% experiments/107-true-baselines/adaptive_lambda/README.md, "Check" table (arm "geo x 3.15 + censored term
% (latent_pilot)": 63 of 72 vs 54 for the previous version, +10 / -1, size ratio 1.04 on pass-in-both). VERIFY: whether
% that check arm included the 0.1 warm-up (run_tonight.sh stage 4 fits check_geo_x3.15_warmup_censor).
\paragraph{A learning-rate schedule and the censored read.}
Under a constant learning rate, the gate step grows twentyfold as the gate temperature anneals, and some fits degraded late in training.
A geometric decay alone stabilised them but left gates too little total step to close, so circuits grew $2$--$4\times$; raising its peak $3.15\times$ restored the summed step size, and a short warm-up was added for the early spikes the higher peak brought back (\cref{app:loss-details}).
With the censored read added, so that training also sees whether the target survives its site's Top-K, the check panel went from $54$ to $63$ of $72$ targets passing ($10$ gained, $1$ lost) at about the same circuit size (median ratio $1.04$).

% sec:lambda-transfer, tab:run2026-lambda-pilot; experiments/108-batch-size/out/analysis.md (16 TuringLLM targets;
% B16 at 400 steps: 15 / 16 vs 16 / 16, size 0.75x on pass-in-both, 99.9 vs 29.1 fit-minutes = 3.4x; pre-stated rule:
% +-1 target is noise); experiments/107-true-baselines/adaptive_lambda/README.md (every shrinking policy traded passes
% for size on the check panel; losses were necessity failures).
\paragraph{Settings left as they were.}
The penalty was carried to the public models by a size-only rule fixed before a pilot (\cref{sec:lambda-transfer}).
In exploratory tests, batches of $16$ for the same $400$ steps passed $15$ of $16$ TuringLLM targets against $16$, with circuits $0.75\times$ the size where both passed, at $3.4\times$ the fitting time, so the batch of $4$ was kept; and policies that adapted $\lambda$ per target traded passes for size on the check panel, losing passes mostly to necessity, which no training term covers.
Other rejected designs are in \cref{app:membership-studies}.

% Directions: app:production (induce), sec:main-result (mid-band, deep band, attention), app:specificity (amplifiers,
% v1: 30% flagged), app:rank-designs (off-target term: lifts cut 3-7x, but grew circuits), app:resample (anchors
% only; overfiring away from the anchor), adaptive_lambda README (necessity).
\paragraph{Directions not yet taken.}
The same path suggests where the recipe can go next.
Necessity, sufficiency to induce, and transfer to mid-band contexts are scored but not trained for; each could get a term or training contexts in its domain, as the resample test did.
Deep and attention-output targets still pass least (\cref{sec:main-result}), and in the schedule pilot the attention failures kept zero-ablation faithfulness near $0$ under every schedule, so the learning rate does not explain them.
Amplifier circuits survive rank-keep (\cref{app:specificity}); the off-target term cut lifts three- to sevenfold but grew circuits (\cref{app:rank-designs}), so a cheaper specificity term remains open.
And circuits explain the target at its anchors, not at every position (\cref{app:resample}); training on more positions is untested.
\fi\fi
        % ~1/3 page, Discussion (replaces the paragraph
%                                                                    % "Ablation methods are part of the method")
%   \def\frameworkpart{appendix}% =====================================================================================================================
% DAN-208 (drafted 2026-11-08, agent "framework"): WCM as a framework for solving for circuits, not a fixed algorithm.
% Two versions in one file. Choose with \def\frameworkpart{main} or \def\frameworkpart{appendix} before \input:
%
%   \def\frameworkpart{main}\input{drafts/framework_section}        % ~1/3 page, Discussion (replaces the paragraph
%                                                                    % "Ablation methods are part of the method")
%   \def\frameworkpart{appendix}\input{drafts/framework_section}    % ~1.3 pages, its own appendix section
%
% New labels: sec:framework (main), app:framework, tab:framework (appendix). Every other \cref target exists in
% main.tex as of 2026-11-08. Every number below is already printed in main.tex (source comments give the place);
% the one exception, experiment 108, is marked NOT-YET-IN-PAPER where it appears.
% Vocabulary: paper/notes/vocabulary.md (no "endpoint", "floor", "drivers"; "fill" only as main.tex already uses it).
% Never-claim list respected: no "first"; nothing here claims novelty.
% =====================================================================================================================
\providecommand{\frameworkpart}{main}
\def\frameworkmain{main}
\def\frameworkappendix{appendix}
\ifx\frameworkpart\frameworkmain
% ---------------------------------------------------------------------------------------------------------------------
% (1) MAIN-TEXT VERSION. Proposed place: sec:discussion, in place of the paragraph "Ablation methods are part of the
% method" (main.tex lines 548-552), whose evidence it repeats and extends; the paragraph "Recommendations" (lines
% 564-566) can then point here. If it must sit in Method instead, it fits after "Pass criterion" (line 341), before
% \section{Experiments and Results}, but there it would cite results not yet shown.
% Numbers: 2-14% = abstract / run2026_alpha1 (sec:weighted); "one to three orders of magnitude" = app:single-method
% (exploratory); contrast sources = app:contrast-source (protocol v1); resample = app:resample.
% ---------------------------------------------------------------------------------------------------------------------
\subsection{Using WCM as a Framework}
\label{sec:framework}
WCM is better read as a framework for solving for circuits than as one fixed algorithm.
Its parts are interchangeable: an objective (the quantity reproduced, the ablation methods it is reproduced under, the regularisers, and their weights; \cref{eq:loss}), a parametrisation (a gate and a coefficient per node), data (training, contrast, and held-out contexts), and an optimiser.
The production recipe is one setting of these parts, each chosen by an experiment (\cref{tab:framework}), and changing a part changes the circuits found.
With every coefficient fixed at $1$ the same nodes pass for $2$--$14\%$ of targets (\cref{sec:weighted}); in exploratory fits, coefficients trained without zero ablation overshoot it by one to three orders of magnitude (\cref{app:single-method}); circuits trained against close contrast contexts are the most faithful on close backgrounds and the least on any other (\cref{app:contrast-source}); and circuits failed a resample test until training included a resample term (\cref{app:resample}).
When WCM fails an evaluation, we therefore suggest first asking whether training ever saw that evaluation's domain, and if not, adding the term, contexts, or setting that puts it there, while keeping the evaluation's own contexts and donors out of training so that it stays a held-out test (\cref{app:framework}).
A pass rate then describes a recipe, a dictionary, and an evaluation together, not WCM alone, which is why the ablation methods used for training and evaluation belong in the method statement~\citep{mueller2025mib, mahale2026explanation}.
%
\else\ifx\frameworkpart\frameworkappendix
% ---------------------------------------------------------------------------------------------------------------------
% (2) APPENDIX VERSION. Proposed place: a new appendix section directly before "Weighted Circuit Masking Design
% Ablations" (main.tex line 2003, app:membership-studies), so that the design ablations it summarises follow it; or
% directly after app:protocol / app:loss-details (before line 853) if Daniel prefers it next to tab:protocol.
% ---------------------------------------------------------------------------------------------------------------------
\section{Building with Weighted Circuit Masking}
\label{app:framework}
Weighted circuit masking fixes a problem, \labelcref{eq:p2}, and a way of relaxing it, but leaves most of the recipe open.
\Cref{tab:framework} lists its slots: what is reproduced, over which nodes, with which parametrisation, under which ablation methods and backgrounds, with which regularisers and optimiser, from which contexts, and how the result is scored.
The production choice in each slot was made by an experiment reported in these appendices, most of them on TuringLLM and carried over unchanged to the public models (\cref{tab:protocol}, Provenance); none is claimed to be the best setting, and several are known to trade one evaluation against another.

{\footnotesize
\setlength{\tabcolsep}{3pt}
\renewcommand{\arraystretch}{1.15}
\begin{longtable}{@{}p{0.13\linewidth}p{0.33\linewidth}p{0.33\linewidth}p{0.16\linewidth}@{}}
\caption{The slots of weighted circuit masking, the production choice in each, and the alternatives this paper examined or that the formulation admits.}
\label{tab:framework} \\
\toprule
\textbf{Slot} & \textbf{Production choice} & \textbf{Alternatives} & \textbf{Evidence} \\
\midrule
\endfirsthead
\toprule
\textbf{Slot} & \textbf{Production choice} & \textbf{Alternatives} & \textbf{Evidence} \\
\midrule
\endhead
\midrule
\multicolumn{4}{r@{}}{\emph{continued on the next page}} \\
\endfoot
\bottomrule
\endlastfoot
Target & the pre-activation $\preact$ of one SAE latent at its anchors; censored read under Z, C, A & an output metric, as in learned-mask circuits for outputs~\citep{bhaskar2024finding, caples2025scaling}; output- and latent-targeted circuits fail each other's objective (exploratory) & \cref{sec:ablmask}, \cref{app:secondary,app:logit-faith} \\
Nodes & every latent at every SAE site upstream of the target; SAE errors passed through; nodes recomputed & dense ReLU dictionaries (the method replicates; the random baseline needs its own design); members clamped (resample diagnostic) & \cref{sec:problem}, \cref{app:dense-relu,app:resample} \\
Parametrisation & gate $\sigma(\theta/T)$ and coefficient $\mathrm{softplus}(\psi)$ per node, $\alpha_0 = 1$ & $\alpha \equiv 1$ (unweighted circuit masking) & \cref{sec:weighted}, \cref{tab:run2026-alpha1} \\
Training ablations & Z (weight $1$), C and A ($0.25$ each) under evaluation's sparsity-preserving fill; resample D over the training contrast contexts ($1$) & one method alone; dense means; D in place of C; no D & \cref{app:single-method,app:respect-topk,app:resample} \\
Contrast background & $64$ close, silence-checked contexts per target ($48$ train, $16$ held out) & random, distant, or fused sources; the unchecked retrieval store & \cref{app:stores,app:contrast-source} \\
Regularisers & flat penalty $\lambda = 10^{-3}$ per open gate; closed-gate charge $\lambda(1-\sigma(\theta))|\alpha - 1|$; rank-keep $\gamma_{\mathrm{K}} = 10^{-1}$; one shared normaliser $\nu$ & other penalties $\lambda$ (sweep); off-target term, Top-K hinge, pull to rank $1$, firing margin; per-term normalisers & \cref{sec:lambda-transfer}, \cref{app:loss-details,app:rank-designs,app:membership-studies} \\
Optimiser & AdamW, $400$ steps, batch $4$; gate temperature annealed $1 \to 0.05$; learning-rate warm-up, then geometric decay & straight-through binarisation; $200$-step and reach-then-hold schedules; larger batches (exploratory) & \cref{tab:protocol}, \cref{app:membership-studies,app:framework-history} \\
Training contexts & $32$ strongest and $16$ mid-band activating contexts, deduplicated and split by whole groups of copies & a split without deduplication, which leaked copies into held-out sets & \cref{sec:setting}, \cref{app:heldout} \\
Evaluation & held-out strongest contexts; Z, A, C in $[0.75, 1.5]$ and necessity $\geq 0.9$; vacuous targets excluded and counted & reported, not required: mid-band contexts, sufficiency to induce, donor-background test, specificity, $\alpha = 1$ rescoring, random-circuit baselines, stable core & \cref{sec:eval}, \cref{app:protocol,app:nulls,app:underdetermination} \\
\end{longtable}}

\paragraph{Training in the evaluation's domain.}
Across these experiments, circuits did best on the evaluations their training resembled.
When a fit fails an evaluation, the failure is often a property of the recipe, and a slot of \cref{tab:framework} can put the fit in that evaluation's domain.
% app:resample (095: 72 targets, 0% -> 24%, CI +14 to +33; criterion 69% -> 62%, CI -17 to +3; 1.18x size, 1.45x time;
% D in place of C: 51%, CI -29 to -7), production paragraph (exp 102: 19 -> 39 and 48 -> 54 of 72, 1.33x size), T10
% (tab:run2026-tl-noR: +5.2, -0.6 to 11.6).
\emph{Real backgrounds.} Circuits trained only against zero and means failed a resample test almost always; adding a resample term over the training contrast contexts raised the resample pass rate on $72$ targets from $0\%$ to $24\%$, at $1.18\times$ the circuit size and $1.45\times$ the fitting time, and with rank-keep raised it took a $72$-target check panel from $19$ to $39$ targets passing (\cref{app:resample}).
Using it in place of the contrast-mean term gained as much on the resample test but cost the pass criterion, so terms should be added for a new evaluation rather than swapped for an old one; on the production sweep the term's effect on the pass criterion itself is $5.2$ points and not significant (\cref{tab:run2026-tl-noR}).
\emph{Ablation methods and fills.} A mask trained under one ablation method learns what that method rewards: without the zero-ablation term, coefficients overshoot zero ablation by orders of magnitude, and circuits trained against dense means can fail under the sparsity-preserving fill that evaluation uses (\cref{app:single-method,app:respect-topk}).
\emph{Top-K survival.} Reproducing the pre-activation does not keep the target in its site's Top-K; the censored read and rank-keep train for it, though neither removes amplifier circuits (\cref{app:loss-details,app:specificity}).
\emph{Other backgrounds.} Contrast faithfulness generalises furthest when training uses random or distant contrast contexts (\cref{app:contrast-source}).
\emph{Size.} Circuits shrink with $\lambda$ at a nearly constant pass rate up to a plateau's edge (\cref{sec:lambda-transfer}).
Two evaluations have no training term yet: sufficiency to induce, which no circuit is trained on (median $0.78$--$0.88$ on held-out contrast contexts, $0.35$ for TuringLLM attention outputs; \cref{app:production}), and transfer to weaker, mid-band contexts (\cref{sec:main-result}).
The framework's prescription for them is a term or a training set in their domain, which we have not tested.
% CANDIDATE EVIDENCE, NOT IN THE PAPER (main.tex line 503: the arena results are reported separately). If Daniel wants
% them here, each needs its own sentence and its own appendix or citation:
%  - arena A tuning (experiments/107-true-baselines/arena_a/TUNING.md): on GPT-2 transcoder + z-SAE substrate the
%    necessity hinge (NEC_W; off in production, docs/wcm-specification.md §2.3) was chosen at 0.3 by a pre-registered
%    9-target rule; NEC_W 1.0 raised nec_a (3.tc.3752 0.32 -> 0.95) but pushed A out of band (1.54). The censored read
%    is inert there by construction (k = d_sae).
%  - arena A scope a1 (z + z2 + tc) vs a2 (tc only) and the wikitext-smoke vs OpenWebText corpus test (DAN-208 text;
%    not traced to an output file in this pass).
%  - logit / maximise / behavioural objectives (arena C/D; memory notes say "don't quote yet").
A fit trained for an evaluation is, for that evaluation, scored on home turf (\cref{app:method-comparison}), so the evaluation stays meaningful only if its contexts and donors are never seen in training, as the donor-background test's donors are not (\cref{tab:protocol}), and the evaluations not trained for are reported alongside.

\paragraph{Choosing the sparsity penalty.}
% sec:lambda-transfer, tab:run2026-tl-sweep, tab:run2026-lambda-pilot.
The penalty sets circuit size and, past a point, the pass rate: on the TuringLLM sweep the pass rate holds at $89$--$93\%$ from $\lambda = 2.5 \times 10^{-5}$ to $5 \times 10^{-4}$ while each doubling of $\lambda$ shrinks circuits $1.4$--$1.9\times$, and it declines beyond (\cref{sec:lambda-transfer}).
We recommend fixing the selection rule before looking at pass rates, as our size-only pilot rule did on the public models, where it kept $\lambda = 10^{-3}$ on both (\cref{tab:run2026-lambda-pilot}), and reporting a curve over $\lambda$ rather than one point, since one setting that transfers is not the best setting for each model.

\paragraph{A new model or dictionary.}
Only the model and dictionary adapters change between our three models (\cref{sec:setting}); four things needed checking.
\emph{Precision}: TuringLLM runs in the bfloat16 its pipeline was validated in, but on the public models a bfloat16 data pass did not reproduce the float32 top-context sets, so they run in float32 (\cref{tab:protocol}); a precision check belongs before any fit, since node sets also move with rounding (\cref{app:underdetermination}).
\emph{The dictionary's sparsity}: $k$ sets the fill's budget and the censored read's cut; on a dictionary that keeps every latent ($k = d_{\mathrm{sae}}$) the censored read has no cut to model, and on dense ReLU dictionaries a random-circuit baseline with fitted coefficients stops discriminating, so necessity, not faithfulness, separates real from random circuits there (\cref{app:dense-relu}).
% The k = d_sae statement follows from the definition of tau^(K) (no K-th other latent); docs/wcm-specification.md §7.6.
\emph{Contrast retrieval}: contrast contexts need a pooled representation for every corpus sequence and a silence check, which rejects $14$--$51\%$ of the closest candidates on the targets we inspected (\cref{app:stores}).
\emph{Corpus copies}: duplicated windows silently leak into held-out sets unless pools are deduplicated and split by groups of copies (\cref{app:heldout}).

\paragraph{What to report.}
Because each slot changes the circuits, a result is interpretable only with its recipe: the ablation methods used in training and in evaluation; faithfulness under each method and necessity, on held-out contexts, by depth band and site kind, with intervals over targets; excluded and vacuous targets counted; circuit size with whether coefficients were fitted, and the $\alpha = 1$ rescoring; random-circuit baselines given the same fitting; and the reported-only evaluations (sufficiency to induce, the donor-background test, specificity, mid-band transfer) beside the pass rate (\cref{sec:discussion}).

\paragraph{Choices that matter, and what stays open.}
The ablation method changes not only scores but rankings: the two mean-ablation variants can order small node sets differently (\cref{app:respect-topk}), and in exploratory comparisons a method tended to win the scores of the ablation method it was trained or ranked under (\cref{app:method-comparison}), so comparisons between recipes need the full matrix of recipe against ablation method.
Node sets are not identified by any recipe: a refit with reordered training contexts shares a median Jaccard of $0.39$--$0.51$ with the original and its shared core alone passes for $44$--$72\%$ of targets (\cref{app:underdetermination}), so changes to a slot should be judged by sizes and scores, not by which latents moved.
Finally, the production choices were calibrated on TuringLLM targets that overlap its evaluation set (\cref{sec:setting}), and on the public models every setting except $\lambda$ carried over unchanged (\cref{tab:protocol}, Provenance), so on a new model they are defaults to be checked, not tuned values.

% ---------------------------------------------------------------------------------------------------------------------
% 2026-11-08 (Daniel, via coordinator): the development path as evidence for where WCM can go next. Every number is
% traced; sources in the comment before each step. Panels and protocols differ between steps, so numbers are compared
% only within a step (each step is a paired test against the recipe before it).
% ---------------------------------------------------------------------------------------------------------------------
\subsection{How the Production Recipe Was Reached}
\label{app:framework-history}
The production recipe is the end of a sequence of changes, each made because a run exposed a failure, and each tested against the recipe before it on the same targets, mostly under a selection rule stated before the comparison.
We list them in order as evidence of which kinds of change have paid off; the panels differ between steps, so numbers compare only within a step.

% app:stores (38% fallback), app:contrast-source; experiments/055-close-contrast, 071-contrast-sweep.
\paragraph{Contrast contexts built per target.}
The first contrast contexts came from a retrieval store whose index covered too few sequences, so $38\%$ of TuringLLM targets fell back to one shared list; WCM now selects contrast contexts per target at fit time, ranked by similarity to its strongest contexts and kept only if the target is silent (\cref{app:stores}), and a later comparison of contrast sources showed how much that choice shapes the C term (\cref{app:contrast-source}).
The lesson: inspect the data each training term actually sees.

% app:specificity, app:rank-designs; experiments/056-specificity.
\paragraph{Rank-keep against off-target lifts.}
Reproducing the pre-activation left some circuits restoring their target by lifting its whole neighbourhood; of the terms tried, rank-keep alone kept the gains without the alternatives' costs, and entered the first protocol at a small weight (\cref{app:rank-designs,app:specificity}).

% app:v1-run (15,046 targets, 46%, attention 14%, layers 8-11 22%); app:heldout (66.8%); app:resample (0.1% vs 43.1%,
% 747 targets); fig:failure-anatomy.
\paragraph{The first run at scale.}
Under that protocol, WCM was fitted to $15{,}046$ TuringLLM targets, and $46\%$ passed, $22\%$ at layers 8--11 and $14\%$ for attention outputs (\cref{app:v1-run}).
Analyses of the run exposed three problems: $66.8\%$ of held-out strongest contexts repeated a training context, which led to deduplicated records split by groups of copies (\cref{app:heldout}); with depth, failing circuits increasingly failed only the two mean ablations (\cref{fig:failure-anatomy}); and $0.1\%$ of circuits passed a resample test, against $43.1\%$ under the pass criterion (\cref{app:resample}).

% experiments/093-v2-validation/README.md (out/tables.md): 72 targets, 8 per band x kind, paired, pass floor 0.8, the
% v1 run's as-run held-out sets (leaky; same for both arms): 31 -> 50 of 72, +26.4 pp [15.3, 38.9]; L8-11 2 -> 11 of 24;
% attn 3 -> 12 of 24; median dZ 0.00; size ratio 0.91; fit time 1.27x. Pilot: experiments/090-sparse-fill-training.
% Also main.tex comment at app:respect-topk ("The same fill in training").
\paragraph{Training under the fill that evaluation uses.}
The mean-ablation failures were a mismatch: training replaced non-members with dense means, a code no Top-K dictionary emits, while evaluation re-imposed the Top-K (\cref{app:respect-topk}).
Training the A and C terms under evaluation's own fill and means raised the pass rate on $72$ stratified targets from $31$ to $50$ (paired difference $+26.4$ points, $95\%$ interval $15.3$--$38.9$), from $2$ to $11$ of $24$ at layers 8--11 and from $3$ to $12$ of $24$ for attention outputs, with zero-ablation faithfulness unchanged (median paired change $0.00$), circuits $0.91\times$ the size, and $1.27\times$ the fitting time.
The gain came from aligning training with evaluation, not from a new term.

% app:resample (095 numbers and the production paragraph: 102's 19 -> 39, 48 -> 54 of 72, 1.33x size);
% experiments/102-pass-improvement/README.md line 89 (diagnosis: 15 of 18 Z failures of the previous version restored
% the pre-activation and lost at the Top-K cut); tab:run2026-tl-noR.
\paragraph{A resample term, and a stronger rank-keep.}
To survive real backgrounds, training gained the resample term (above and \cref{app:resample}).
On the previous version's zero-ablation failures, $15$ of $18$ restored the target's pre-activation but lost it at the Top-K cut, so rank-keep was raised from $3 \times 10^{-3}$ to $10^{-1}$ with the resample term at full weight; chosen by a rule stated before the comparison, the pair took a $72$-target check panel from $48$ to $54$ targets passing and from $19$ to $39$ passing the resample test, at $1.33\times$ the circuit size (\cref{app:resample}).
In the production sweep, the arm trained without the resample term passes $5.2$ points less, not significantly (\cref{tab:run2026-tl-noR}).

% experiments/107-true-baselines/latent_pilot/README.md (9 targets enriched for failures: 3 -> 5 of 9, size ratio ~1.04;
% plain geometric: circuits 2-4x larger, sum(lr) 6.3 vs 20; constant lr makes lr/T grow 20x) and
% experiments/107-true-baselines/adaptive_lambda/README.md, "Check" table (arm "geo x 3.15 + censored term
% (latent_pilot)": 63 of 72 vs 54 for the previous version, +10 / -1, size ratio 1.04 on pass-in-both). VERIFY: whether
% that check arm included the 0.1 warm-up (run_tonight.sh stage 4 fits check_geo_x3.15_warmup_censor).
\paragraph{A learning-rate schedule and the censored read.}
Under a constant learning rate, the gate step grows twentyfold as the gate temperature anneals, and some fits degraded late in training.
A geometric decay alone stabilised them but left gates too little total step to close, so circuits grew $2$--$4\times$; raising its peak $3.15\times$ restored the summed step size, and a short warm-up was added for the early spikes the higher peak brought back (\cref{app:loss-details}).
With the censored read added, so that training also sees whether the target survives its site's Top-K, the check panel went from $54$ to $63$ of $72$ targets passing ($10$ gained, $1$ lost) at about the same circuit size (median ratio $1.04$).

% sec:lambda-transfer, tab:run2026-lambda-pilot; experiments/108-batch-size/out/analysis.md (16 TuringLLM targets;
% B16 at 400 steps: 15 / 16 vs 16 / 16, size 0.75x on pass-in-both, 99.9 vs 29.1 fit-minutes = 3.4x; pre-stated rule:
% +-1 target is noise); experiments/107-true-baselines/adaptive_lambda/README.md (every shrinking policy traded passes
% for size on the check panel; losses were necessity failures).
\paragraph{Settings left as they were.}
The penalty was carried to the public models by a size-only rule fixed before a pilot (\cref{sec:lambda-transfer}).
In exploratory tests, batches of $16$ for the same $400$ steps passed $15$ of $16$ TuringLLM targets against $16$, with circuits $0.75\times$ the size where both passed, at $3.4\times$ the fitting time, so the batch of $4$ was kept; and policies that adapted $\lambda$ per target traded passes for size on the check panel, losing passes mostly to necessity, which no training term covers.
Other rejected designs are in \cref{app:membership-studies}.

% Directions: app:production (induce), sec:main-result (mid-band, deep band, attention), app:specificity (amplifiers,
% v1: 30% flagged), app:rank-designs (off-target term: lifts cut 3-7x, but grew circuits), app:resample (anchors
% only; overfiring away from the anchor), adaptive_lambda README (necessity).
\paragraph{Directions not yet taken.}
The same path suggests where the recipe can go next.
Necessity, sufficiency to induce, and transfer to mid-band contexts are scored but not trained for; each could get a term or training contexts in its domain, as the resample test did.
Deep and attention-output targets still pass least (\cref{sec:main-result}), and in the schedule pilot the attention failures kept zero-ablation faithfulness near $0$ under every schedule, so the learning rate does not explain them.
Amplifier circuits survive rank-keep (\cref{app:specificity}); the off-target term cut lifts three- to sevenfold but grew circuits (\cref{app:rank-designs}), so a cheaper specificity term remains open.
And circuits explain the target at its anchors, not at every position (\cref{app:resample}); training on more positions is untested.
\fi\fi
    % ~1.3 pages, its own appendix section
%
% New labels: sec:framework (main), app:framework, tab:framework (appendix). Every other \cref target exists in
% main.tex as of 2026-11-08. Every number below is already printed in main.tex (source comments give the place);
% the one exception, experiment 108, is marked NOT-YET-IN-PAPER where it appears.
% Vocabulary: paper/notes/vocabulary.md (no "endpoint", "floor", "drivers"; "fill" only as main.tex already uses it).
% Never-claim list respected: no "first"; nothing here claims novelty.
% =====================================================================================================================
\providecommand{\frameworkpart}{main}
\def\frameworkmain{main}
\def\frameworkappendix{appendix}
\ifx\frameworkpart\frameworkmain
% ---------------------------------------------------------------------------------------------------------------------
% (1) MAIN-TEXT VERSION. Proposed place: sec:discussion, in place of the paragraph "Ablation methods are part of the
% method" (main.tex lines 548-552), whose evidence it repeats and extends; the paragraph "Recommendations" (lines
% 564-566) can then point here. If it must sit in Method instead, it fits after "Pass criterion" (line 341), before
% \section{Experiments and Results}, but there it would cite results not yet shown.
% Numbers: 2-14% = abstract / run2026_alpha1 (sec:weighted); "one to three orders of magnitude" = app:single-method
% (exploratory); contrast sources = app:contrast-source (protocol v1); resample = app:resample.
% ---------------------------------------------------------------------------------------------------------------------
\subsection{Using WCM as a Framework}
\label{sec:framework}
WCM is better read as a framework for solving for circuits than as one fixed algorithm.
Its parts are interchangeable: an objective (the quantity reproduced, the ablation methods it is reproduced under, the regularisers, and their weights; \cref{eq:loss}), a parametrisation (a gate and a coefficient per node), data (training, contrast, and held-out contexts), and an optimiser.
The production recipe is one setting of these parts, each chosen by an experiment (\cref{tab:framework}), and changing a part changes the circuits found.
With every coefficient fixed at $1$ the same nodes pass for $2$--$14\%$ of targets (\cref{sec:weighted}); in exploratory fits, coefficients trained without zero ablation overshoot it by one to three orders of magnitude (\cref{app:single-method}); circuits trained against close contrast contexts are the most faithful on close backgrounds and the least on any other (\cref{app:contrast-source}); and circuits failed a resample test until training included a resample term (\cref{app:resample}).
When WCM fails an evaluation, we therefore suggest first asking whether training ever saw that evaluation's domain, and if not, adding the term, contexts, or setting that puts it there, while keeping the evaluation's own contexts and donors out of training so that it stays a held-out test (\cref{app:framework}).
A pass rate then describes a recipe, a dictionary, and an evaluation together, not WCM alone, which is why the ablation methods used for training and evaluation belong in the method statement~\citep{mueller2025mib, mahale2026explanation}.
%
\else\ifx\frameworkpart\frameworkappendix
% ---------------------------------------------------------------------------------------------------------------------
% (2) APPENDIX VERSION. Proposed place: a new appendix section directly before "Weighted Circuit Masking Design
% Ablations" (main.tex line 2003, app:membership-studies), so that the design ablations it summarises follow it; or
% directly after app:protocol / app:loss-details (before line 853) if Daniel prefers it next to tab:protocol.
% ---------------------------------------------------------------------------------------------------------------------
\section{Building with Weighted Circuit Masking}
\label{app:framework}
Weighted circuit masking fixes a problem, \labelcref{eq:p2}, and a way of relaxing it, but leaves most of the recipe open.
\Cref{tab:framework} lists its slots: what is reproduced, over which nodes, with which parametrisation, under which ablation methods and backgrounds, with which regularisers and optimiser, from which contexts, and how the result is scored.
The production choice in each slot was made by an experiment reported in these appendices, most of them on TuringLLM and carried over unchanged to the public models (\cref{tab:protocol}, Provenance); none is claimed to be the best setting, and several are known to trade one evaluation against another.

{\footnotesize
\setlength{\tabcolsep}{3pt}
\renewcommand{\arraystretch}{1.15}
\begin{longtable}{@{}p{0.13\linewidth}p{0.33\linewidth}p{0.33\linewidth}p{0.16\linewidth}@{}}
\caption{The slots of weighted circuit masking, the production choice in each, and the alternatives this paper examined or that the formulation admits.}
\label{tab:framework} \\
\toprule
\textbf{Slot} & \textbf{Production choice} & \textbf{Alternatives} & \textbf{Evidence} \\
\midrule
\endfirsthead
\toprule
\textbf{Slot} & \textbf{Production choice} & \textbf{Alternatives} & \textbf{Evidence} \\
\midrule
\endhead
\midrule
\multicolumn{4}{r@{}}{\emph{continued on the next page}} \\
\endfoot
\bottomrule
\endlastfoot
Target & the pre-activation $\preact$ of one SAE latent at its anchors; censored read under Z, C, A & an output metric, as in learned-mask circuits for outputs~\citep{bhaskar2024finding, caples2025scaling}; output- and latent-targeted circuits fail each other's objective (exploratory) & \cref{sec:ablmask}, \cref{app:secondary,app:logit-faith} \\
Nodes & every latent at every SAE site upstream of the target; SAE errors passed through; nodes recomputed & dense ReLU dictionaries (the method replicates; the random baseline needs its own design); members clamped (resample diagnostic) & \cref{sec:problem}, \cref{app:dense-relu,app:resample} \\
Parametrisation & gate $\sigma(\theta/T)$ and coefficient $\mathrm{softplus}(\psi)$ per node, $\alpha_0 = 1$ & $\alpha \equiv 1$ (unweighted circuit masking) & \cref{sec:weighted}, \cref{tab:run2026-alpha1} \\
Training ablations & Z (weight $1$), C and A ($0.25$ each) under evaluation's sparsity-preserving fill; resample D over the training contrast contexts ($1$) & one method alone; dense means; D in place of C; no D & \cref{app:single-method,app:respect-topk,app:resample} \\
Contrast background & $64$ close, silence-checked contexts per target ($48$ train, $16$ held out) & random, distant, or fused sources; the unchecked retrieval store & \cref{app:stores,app:contrast-source} \\
Regularisers & flat penalty $\lambda = 10^{-3}$ per open gate; closed-gate charge $\lambda(1-\sigma(\theta))|\alpha - 1|$; rank-keep $\gamma_{\mathrm{K}} = 10^{-1}$; one shared normaliser $\nu$ & other penalties $\lambda$ (sweep); off-target term, Top-K hinge, pull to rank $1$, firing margin; per-term normalisers & \cref{sec:lambda-transfer}, \cref{app:loss-details,app:rank-designs,app:membership-studies} \\
Optimiser & AdamW, $400$ steps, batch $4$; gate temperature annealed $1 \to 0.05$; learning-rate warm-up, then geometric decay & straight-through binarisation; $200$-step and reach-then-hold schedules; larger batches (exploratory) & \cref{tab:protocol}, \cref{app:membership-studies,app:framework-history} \\
Training contexts & $32$ strongest and $16$ mid-band activating contexts, deduplicated and split by whole groups of copies & a split without deduplication, which leaked copies into held-out sets & \cref{sec:setting}, \cref{app:heldout} \\
Evaluation & held-out strongest contexts; Z, A, C in $[0.75, 1.5]$ and necessity $\geq 0.9$; vacuous targets excluded and counted & reported, not required: mid-band contexts, sufficiency to induce, donor-background test, specificity, $\alpha = 1$ rescoring, random-circuit baselines, stable core & \cref{sec:eval}, \cref{app:protocol,app:nulls,app:underdetermination} \\
\end{longtable}}

\paragraph{Training in the evaluation's domain.}
Across these experiments, circuits did best on the evaluations their training resembled.
When a fit fails an evaluation, the failure is often a property of the recipe, and a slot of \cref{tab:framework} can put the fit in that evaluation's domain.
% app:resample (095: 72 targets, 0% -> 24%, CI +14 to +33; criterion 69% -> 62%, CI -17 to +3; 1.18x size, 1.45x time;
% D in place of C: 51%, CI -29 to -7), production paragraph (exp 102: 19 -> 39 and 48 -> 54 of 72, 1.33x size), T10
% (tab:run2026-tl-noR: +5.2, -0.6 to 11.6).
\emph{Real backgrounds.} Circuits trained only against zero and means failed a resample test almost always; adding a resample term over the training contrast contexts raised the resample pass rate on $72$ targets from $0\%$ to $24\%$, at $1.18\times$ the circuit size and $1.45\times$ the fitting time, and with rank-keep raised it took a $72$-target check panel from $19$ to $39$ targets passing (\cref{app:resample}).
Using it in place of the contrast-mean term gained as much on the resample test but cost the pass criterion, so terms should be added for a new evaluation rather than swapped for an old one; on the production sweep the term's effect on the pass criterion itself is $5.2$ points and not significant (\cref{tab:run2026-tl-noR}).
\emph{Ablation methods and fills.} A mask trained under one ablation method learns what that method rewards: without the zero-ablation term, coefficients overshoot zero ablation by orders of magnitude, and circuits trained against dense means can fail under the sparsity-preserving fill that evaluation uses (\cref{app:single-method,app:respect-topk}).
\emph{Top-K survival.} Reproducing the pre-activation does not keep the target in its site's Top-K; the censored read and rank-keep train for it, though neither removes amplifier circuits (\cref{app:loss-details,app:specificity}).
\emph{Other backgrounds.} Contrast faithfulness generalises furthest when training uses random or distant contrast contexts (\cref{app:contrast-source}).
\emph{Size.} Circuits shrink with $\lambda$ at a nearly constant pass rate up to a plateau's edge (\cref{sec:lambda-transfer}).
Two evaluations have no training term yet: sufficiency to induce, which no circuit is trained on (median $0.78$--$0.88$ on held-out contrast contexts, $0.35$ for TuringLLM attention outputs; \cref{app:production}), and transfer to weaker, mid-band contexts (\cref{sec:main-result}).
The framework's prescription for them is a term or a training set in their domain, which we have not tested.
% CANDIDATE EVIDENCE, NOT IN THE PAPER (main.tex line 503: the arena results are reported separately). If Daniel wants
% them here, each needs its own sentence and its own appendix or citation:
%  - arena A tuning (experiments/107-true-baselines/arena_a/TUNING.md): on GPT-2 transcoder + z-SAE substrate the
%    necessity hinge (NEC_W; off in production, docs/wcm-specification.md §2.3) was chosen at 0.3 by a pre-registered
%    9-target rule; NEC_W 1.0 raised nec_a (3.tc.3752 0.32 -> 0.95) but pushed A out of band (1.54). The censored read
%    is inert there by construction (k = d_sae).
%  - arena A scope a1 (z + z2 + tc) vs a2 (tc only) and the wikitext-smoke vs OpenWebText corpus test (DAN-208 text;
%    not traced to an output file in this pass).
%  - logit / maximise / behavioural objectives (arena C/D; memory notes say "don't quote yet").
A fit trained for an evaluation is, for that evaluation, scored on home turf (\cref{app:method-comparison}), so the evaluation stays meaningful only if its contexts and donors are never seen in training, as the donor-background test's donors are not (\cref{tab:protocol}), and the evaluations not trained for are reported alongside.

\paragraph{Choosing the sparsity penalty.}
% sec:lambda-transfer, tab:run2026-tl-sweep, tab:run2026-lambda-pilot.
The penalty sets circuit size and, past a point, the pass rate: on the TuringLLM sweep the pass rate holds at $89$--$93\%$ from $\lambda = 2.5 \times 10^{-5}$ to $5 \times 10^{-4}$ while each doubling of $\lambda$ shrinks circuits $1.4$--$1.9\times$, and it declines beyond (\cref{sec:lambda-transfer}).
We recommend fixing the selection rule before looking at pass rates, as our size-only pilot rule did on the public models, where it kept $\lambda = 10^{-3}$ on both (\cref{tab:run2026-lambda-pilot}), and reporting a curve over $\lambda$ rather than one point, since one setting that transfers is not the best setting for each model.

\paragraph{A new model or dictionary.}
Only the model and dictionary adapters change between our three models (\cref{sec:setting}); four things needed checking.
\emph{Precision}: TuringLLM runs in the bfloat16 its pipeline was validated in, but on the public models a bfloat16 data pass did not reproduce the float32 top-context sets, so they run in float32 (\cref{tab:protocol}); a precision check belongs before any fit, since node sets also move with rounding (\cref{app:underdetermination}).
\emph{The dictionary's sparsity}: $k$ sets the fill's budget and the censored read's cut; on a dictionary that keeps every latent ($k = d_{\mathrm{sae}}$) the censored read has no cut to model, and on dense ReLU dictionaries a random-circuit baseline with fitted coefficients stops discriminating, so necessity, not faithfulness, separates real from random circuits there (\cref{app:dense-relu}).
% The k = d_sae statement follows from the definition of tau^(K) (no K-th other latent); docs/wcm-specification.md §7.6.
\emph{Contrast retrieval}: contrast contexts need a pooled representation for every corpus sequence and a silence check, which rejects $14$--$51\%$ of the closest candidates on the targets we inspected (\cref{app:stores}).
\emph{Corpus copies}: duplicated windows silently leak into held-out sets unless pools are deduplicated and split by groups of copies (\cref{app:heldout}).

\paragraph{What to report.}
Because each slot changes the circuits, a result is interpretable only with its recipe: the ablation methods used in training and in evaluation; faithfulness under each method and necessity, on held-out contexts, by depth band and site kind, with intervals over targets; excluded and vacuous targets counted; circuit size with whether coefficients were fitted, and the $\alpha = 1$ rescoring; random-circuit baselines given the same fitting; and the reported-only evaluations (sufficiency to induce, the donor-background test, specificity, mid-band transfer) beside the pass rate (\cref{sec:discussion}).

\paragraph{Choices that matter, and what stays open.}
The ablation method changes not only scores but rankings: the two mean-ablation variants can order small node sets differently (\cref{app:respect-topk}), and in exploratory comparisons a method tended to win the scores of the ablation method it was trained or ranked under (\cref{app:method-comparison}), so comparisons between recipes need the full matrix of recipe against ablation method.
Node sets are not identified by any recipe: a refit with reordered training contexts shares a median Jaccard of $0.39$--$0.51$ with the original and its shared core alone passes for $44$--$72\%$ of targets (\cref{app:underdetermination}), so changes to a slot should be judged by sizes and scores, not by which latents moved.
Finally, the production choices were calibrated on TuringLLM targets that overlap its evaluation set (\cref{sec:setting}), and on the public models every setting except $\lambda$ carried over unchanged (\cref{tab:protocol}, Provenance), so on a new model they are defaults to be checked, not tuned values.

% ---------------------------------------------------------------------------------------------------------------------
% 2026-11-08 (Daniel, via coordinator): the development path as evidence for where WCM can go next. Every number is
% traced; sources in the comment before each step. Panels and protocols differ between steps, so numbers are compared
% only within a step (each step is a paired test against the recipe before it).
% ---------------------------------------------------------------------------------------------------------------------
\subsection{How the Production Recipe Was Reached}
\label{app:framework-history}
The production recipe is the end of a sequence of changes, each made because a run exposed a failure, and each tested against the recipe before it on the same targets, mostly under a selection rule stated before the comparison.
We list them in order as evidence of which kinds of change have paid off; the panels differ between steps, so numbers compare only within a step.

% app:stores (38% fallback), app:contrast-source; experiments/055-close-contrast, 071-contrast-sweep.
\paragraph{Contrast contexts built per target.}
The first contrast contexts came from a retrieval store whose index covered too few sequences, so $38\%$ of TuringLLM targets fell back to one shared list; WCM now selects contrast contexts per target at fit time, ranked by similarity to its strongest contexts and kept only if the target is silent (\cref{app:stores}), and a later comparison of contrast sources showed how much that choice shapes the C term (\cref{app:contrast-source}).
The lesson: inspect the data each training term actually sees.

% app:specificity, app:rank-designs; experiments/056-specificity.
\paragraph{Rank-keep against off-target lifts.}
Reproducing the pre-activation left some circuits restoring their target by lifting its whole neighbourhood; of the terms tried, rank-keep alone kept the gains without the alternatives' costs, and entered the first protocol at a small weight (\cref{app:rank-designs,app:specificity}).

% app:v1-run (15,046 targets, 46%, attention 14%, layers 8-11 22%); app:heldout (66.8%); app:resample (0.1% vs 43.1%,
% 747 targets); fig:failure-anatomy.
\paragraph{The first run at scale.}
Under that protocol, WCM was fitted to $15{,}046$ TuringLLM targets, and $46\%$ passed, $22\%$ at layers 8--11 and $14\%$ for attention outputs (\cref{app:v1-run}).
Analyses of the run exposed three problems: $66.8\%$ of held-out strongest contexts repeated a training context, which led to deduplicated records split by groups of copies (\cref{app:heldout}); with depth, failing circuits increasingly failed only the two mean ablations (\cref{fig:failure-anatomy}); and $0.1\%$ of circuits passed a resample test, against $43.1\%$ under the pass criterion (\cref{app:resample}).

% experiments/093-v2-validation/README.md (out/tables.md): 72 targets, 8 per band x kind, paired, pass floor 0.8, the
% v1 run's as-run held-out sets (leaky; same for both arms): 31 -> 50 of 72, +26.4 pp [15.3, 38.9]; L8-11 2 -> 11 of 24;
% attn 3 -> 12 of 24; median dZ 0.00; size ratio 0.91; fit time 1.27x. Pilot: experiments/090-sparse-fill-training.
% Also main.tex comment at app:respect-topk ("The same fill in training").
\paragraph{Training under the fill that evaluation uses.}
The mean-ablation failures were a mismatch: training replaced non-members with dense means, a code no Top-K dictionary emits, while evaluation re-imposed the Top-K (\cref{app:respect-topk}).
Training the A and C terms under evaluation's own fill and means raised the pass rate on $72$ stratified targets from $31$ to $50$ (paired difference $+26.4$ points, $95\%$ interval $15.3$--$38.9$), from $2$ to $11$ of $24$ at layers 8--11 and from $3$ to $12$ of $24$ for attention outputs, with zero-ablation faithfulness unchanged (median paired change $0.00$), circuits $0.91\times$ the size, and $1.27\times$ the fitting time.
The gain came from aligning training with evaluation, not from a new term.

% app:resample (095 numbers and the production paragraph: 102's 19 -> 39, 48 -> 54 of 72, 1.33x size);
% experiments/102-pass-improvement/README.md line 89 (diagnosis: 15 of 18 Z failures of the previous version restored
% the pre-activation and lost at the Top-K cut); tab:run2026-tl-noR.
\paragraph{A resample term, and a stronger rank-keep.}
To survive real backgrounds, training gained the resample term (above and \cref{app:resample}).
On the previous version's zero-ablation failures, $15$ of $18$ restored the target's pre-activation but lost it at the Top-K cut, so rank-keep was raised from $3 \times 10^{-3}$ to $10^{-1}$ with the resample term at full weight; chosen by a rule stated before the comparison, the pair took a $72$-target check panel from $48$ to $54$ targets passing and from $19$ to $39$ passing the resample test, at $1.33\times$ the circuit size (\cref{app:resample}).
In the production sweep, the arm trained without the resample term passes $5.2$ points less, not significantly (\cref{tab:run2026-tl-noR}).

% experiments/107-true-baselines/latent_pilot/README.md (9 targets enriched for failures: 3 -> 5 of 9, size ratio ~1.04;
% plain geometric: circuits 2-4x larger, sum(lr) 6.3 vs 20; constant lr makes lr/T grow 20x) and
% experiments/107-true-baselines/adaptive_lambda/README.md, "Check" table (arm "geo x 3.15 + censored term
% (latent_pilot)": 63 of 72 vs 54 for the previous version, +10 / -1, size ratio 1.04 on pass-in-both). VERIFY: whether
% that check arm included the 0.1 warm-up (run_tonight.sh stage 4 fits check_geo_x3.15_warmup_censor).
\paragraph{A learning-rate schedule and the censored read.}
Under a constant learning rate, the gate step grows twentyfold as the gate temperature anneals, and some fits degraded late in training.
A geometric decay alone stabilised them but left gates too little total step to close, so circuits grew $2$--$4\times$; raising its peak $3.15\times$ restored the summed step size, and a short warm-up was added for the early spikes the higher peak brought back (\cref{app:loss-details}).
With the censored read added, so that training also sees whether the target survives its site's Top-K, the check panel went from $54$ to $63$ of $72$ targets passing ($10$ gained, $1$ lost) at about the same circuit size (median ratio $1.04$).

% sec:lambda-transfer, tab:run2026-lambda-pilot; experiments/108-batch-size/out/analysis.md (16 TuringLLM targets;
% B16 at 400 steps: 15 / 16 vs 16 / 16, size 0.75x on pass-in-both, 99.9 vs 29.1 fit-minutes = 3.4x; pre-stated rule:
% +-1 target is noise); experiments/107-true-baselines/adaptive_lambda/README.md (every shrinking policy traded passes
% for size on the check panel; losses were necessity failures).
\paragraph{Settings left as they were.}
The penalty was carried to the public models by a size-only rule fixed before a pilot (\cref{sec:lambda-transfer}).
In exploratory tests, batches of $16$ for the same $400$ steps passed $15$ of $16$ TuringLLM targets against $16$, with circuits $0.75\times$ the size where both passed, at $3.4\times$ the fitting time, so the batch of $4$ was kept; and policies that adapted $\lambda$ per target traded passes for size on the check panel, losing passes mostly to necessity, which no training term covers.
Other rejected designs are in \cref{app:membership-studies}.

% Directions: app:production (induce), sec:main-result (mid-band, deep band, attention), app:specificity (amplifiers,
% v1: 30% flagged), app:rank-designs (off-target term: lifts cut 3-7x, but grew circuits), app:resample (anchors
% only; overfiring away from the anchor), adaptive_lambda README (necessity).
\paragraph{Directions not yet taken.}
The same path suggests where the recipe can go next.
Necessity, sufficiency to induce, and transfer to mid-band contexts are scored but not trained for; each could get a term or training contexts in its domain, as the resample test did.
Deep and attention-output targets still pass least (\cref{sec:main-result}), and in the schedule pilot the attention failures kept zero-ablation faithfulness near $0$ under every schedule, so the learning rate does not explain them.
Amplifier circuits survive rank-keep (\cref{app:specificity}); the off-target term cut lifts three- to sevenfold but grew circuits (\cref{app:rank-designs}), so a cheaper specificity term remains open.
And circuits explain the target at its anchors, not at every position (\cref{app:resample}); training on more positions is untested.
\fi\fi
    % ~1.3 pages, its own appendix section
%
% New labels: sec:framework (main), app:framework, tab:framework (appendix). Every other \cref target exists in
% main.tex as of 2026-11-08. Every number below is already printed in main.tex (source comments give the place);
% the one exception, experiment 108, is marked NOT-YET-IN-PAPER where it appears.
% Vocabulary: paper/notes/vocabulary.md (no "endpoint", "floor", "drivers"; "fill" only as main.tex already uses it).
% Never-claim list respected: no "first"; nothing here claims novelty.
% =====================================================================================================================
\providecommand{\frameworkpart}{main}
\def\frameworkmain{main}
\def\frameworkappendix{appendix}
\ifx\frameworkpart\frameworkmain
% ---------------------------------------------------------------------------------------------------------------------
% (1) MAIN-TEXT VERSION. Proposed place: sec:discussion, in place of the paragraph "Ablation methods are part of the
% method" (main.tex lines 548-552), whose evidence it repeats and extends; the paragraph "Recommendations" (lines
% 564-566) can then point here. If it must sit in Method instead, it fits after "Pass criterion" (line 341), before
% \section{Experiments and Results}, but there it would cite results not yet shown.
% Numbers: 2-14% = abstract / run2026_alpha1 (sec:weighted); "one to three orders of magnitude" = app:single-method
% (exploratory); contrast sources = app:contrast-source (protocol v1); resample = app:resample.
% ---------------------------------------------------------------------------------------------------------------------
\subsection{Using WCM as a Framework}
\label{sec:framework}
WCM is better read as a framework for solving for circuits than as one fixed algorithm.
Its parts are interchangeable: an objective (the quantity reproduced, the ablation methods it is reproduced under, the regularisers, and their weights; \cref{eq:loss}), a parametrisation (a gate and a coefficient per node), data (training, contrast, and held-out contexts), and an optimiser.
The production recipe is one setting of these parts, each chosen by an experiment (\cref{tab:framework}), and changing a part changes the circuits found.
With every coefficient fixed at $1$ the same nodes pass for $2$--$14\%$ of targets (\cref{sec:weighted}); in exploratory fits, coefficients trained without zero ablation overshoot it by one to three orders of magnitude (\cref{app:single-method}); circuits trained against close contrast contexts are the most faithful on close backgrounds and the least on any other (\cref{app:contrast-source}); and circuits failed a resample test until training included a resample term (\cref{app:resample}).
When WCM fails an evaluation, we therefore suggest first asking whether training ever saw that evaluation's domain, and if not, adding the term, contexts, or setting that puts it there, while keeping the evaluation's own contexts and donors out of training so that it stays a held-out test (\cref{app:framework}).
A pass rate then describes a recipe, a dictionary, and an evaluation together, not WCM alone, which is why the ablation methods used for training and evaluation belong in the method statement~\citep{mueller2025mib, mahale2026explanation}.
%
\else\ifx\frameworkpart\frameworkappendix
% ---------------------------------------------------------------------------------------------------------------------
% (2) APPENDIX VERSION. Proposed place: a new appendix section directly before "Weighted Circuit Masking Design
% Ablations" (main.tex line 2003, app:membership-studies), so that the design ablations it summarises follow it; or
% directly after app:protocol / app:loss-details (before line 853) if Daniel prefers it next to tab:protocol.
% ---------------------------------------------------------------------------------------------------------------------
\section{Building with Weighted Circuit Masking}
\label{app:framework}
Weighted circuit masking fixes a problem, \labelcref{eq:p2}, and a way of relaxing it, but leaves most of the recipe open.
\Cref{tab:framework} lists its slots: what is reproduced, over which nodes, with which parametrisation, under which ablation methods and backgrounds, with which regularisers and optimiser, from which contexts, and how the result is scored.
The production choice in each slot was made by an experiment reported in these appendices, most of them on TuringLLM and carried over unchanged to the public models (\cref{tab:protocol}, Provenance); none is claimed to be the best setting, and several are known to trade one evaluation against another.

{\footnotesize
\setlength{\tabcolsep}{3pt}
\renewcommand{\arraystretch}{1.15}
\begin{longtable}{@{}p{0.13\linewidth}p{0.33\linewidth}p{0.33\linewidth}p{0.16\linewidth}@{}}
\caption{The slots of weighted circuit masking, the production choice in each, and the alternatives this paper examined or that the formulation admits.}
\label{tab:framework} \\
\toprule
\textbf{Slot} & \textbf{Production choice} & \textbf{Alternatives} & \textbf{Evidence} \\
\midrule
\endfirsthead
\toprule
\textbf{Slot} & \textbf{Production choice} & \textbf{Alternatives} & \textbf{Evidence} \\
\midrule
\endhead
\midrule
\multicolumn{4}{r@{}}{\emph{continued on the next page}} \\
\endfoot
\bottomrule
\endlastfoot
Target & the pre-activation $\preact$ of one SAE latent at its anchors; censored read under Z, C, A & an output metric, as in learned-mask circuits for outputs~\citep{bhaskar2024finding, caples2025scaling}; output- and latent-targeted circuits fail each other's objective (exploratory) & \cref{sec:ablmask}, \cref{app:secondary,app:logit-faith} \\
Nodes & every latent at every SAE site upstream of the target; SAE errors passed through; nodes recomputed & dense ReLU dictionaries (the method replicates; the random baseline needs its own design); members clamped (resample diagnostic) & \cref{sec:problem}, \cref{app:dense-relu,app:resample} \\
Parametrisation & gate $\sigma(\theta/T)$ and coefficient $\mathrm{softplus}(\psi)$ per node, $\alpha_0 = 1$ & $\alpha \equiv 1$ (unweighted circuit masking) & \cref{sec:weighted}, \cref{tab:run2026-alpha1} \\
Training ablations & Z (weight $1$), C and A ($0.25$ each) under evaluation's sparsity-preserving fill; resample D over the training contrast contexts ($1$) & one method alone; dense means; D in place of C; no D & \cref{app:single-method,app:respect-topk,app:resample} \\
Contrast background & $64$ close, silence-checked contexts per target ($48$ train, $16$ held out) & random, distant, or fused sources; the unchecked retrieval store & \cref{app:stores,app:contrast-source} \\
Regularisers & flat penalty $\lambda = 10^{-3}$ per open gate; closed-gate charge $\lambda(1-\sigma(\theta))|\alpha - 1|$; rank-keep $\gamma_{\mathrm{K}} = 10^{-1}$; one shared normaliser $\nu$ & other penalties $\lambda$ (sweep); off-target term, Top-K hinge, pull to rank $1$, firing margin; per-term normalisers & \cref{sec:lambda-transfer}, \cref{app:loss-details,app:rank-designs,app:membership-studies} \\
Optimiser & AdamW, $400$ steps, batch $4$; gate temperature annealed $1 \to 0.05$; learning-rate warm-up, then geometric decay & straight-through binarisation; $200$-step and reach-then-hold schedules; larger batches (exploratory) & \cref{tab:protocol}, \cref{app:membership-studies,app:framework-history} \\
Training contexts & $32$ strongest and $16$ mid-band activating contexts, deduplicated and split by whole groups of copies & a split without deduplication, which leaked copies into held-out sets & \cref{sec:setting}, \cref{app:heldout} \\
Evaluation & held-out strongest contexts; Z, A, C in $[0.75, 1.5]$ and necessity $\geq 0.9$; vacuous targets excluded and counted & reported, not required: mid-band contexts, sufficiency to induce, donor-background test, specificity, $\alpha = 1$ rescoring, random-circuit baselines, stable core & \cref{sec:eval}, \cref{app:protocol,app:nulls,app:underdetermination} \\
\end{longtable}}

\paragraph{Training in the evaluation's domain.}
Across these experiments, circuits did best on the evaluations their training resembled.
When a fit fails an evaluation, the failure is often a property of the recipe, and a slot of \cref{tab:framework} can put the fit in that evaluation's domain.
% app:resample (095: 72 targets, 0% -> 24%, CI +14 to +33; criterion 69% -> 62%, CI -17 to +3; 1.18x size, 1.45x time;
% D in place of C: 51%, CI -29 to -7), production paragraph (exp 102: 19 -> 39 and 48 -> 54 of 72, 1.33x size), T10
% (tab:run2026-tl-noR: +5.2, -0.6 to 11.6).
\emph{Real backgrounds.} Circuits trained only against zero and means failed a resample test almost always; adding a resample term over the training contrast contexts raised the resample pass rate on $72$ targets from $0\%$ to $24\%$, at $1.18\times$ the circuit size and $1.45\times$ the fitting time, and with rank-keep raised it took a $72$-target check panel from $19$ to $39$ targets passing (\cref{app:resample}).
Using it in place of the contrast-mean term gained as much on the resample test but cost the pass criterion, so terms should be added for a new evaluation rather than swapped for an old one; on the production sweep the term's effect on the pass criterion itself is $5.2$ points and not significant (\cref{tab:run2026-tl-noR}).
\emph{Ablation methods and fills.} A mask trained under one ablation method learns what that method rewards: without the zero-ablation term, coefficients overshoot zero ablation by orders of magnitude, and circuits trained against dense means can fail under the sparsity-preserving fill that evaluation uses (\cref{app:single-method,app:respect-topk}).
\emph{Top-K survival.} Reproducing the pre-activation does not keep the target in its site's Top-K; the censored read and rank-keep train for it, though neither removes amplifier circuits (\cref{app:loss-details,app:specificity}).
\emph{Other backgrounds.} Contrast faithfulness generalises furthest when training uses random or distant contrast contexts (\cref{app:contrast-source}).
\emph{Size.} Circuits shrink with $\lambda$ at a nearly constant pass rate up to a plateau's edge (\cref{sec:lambda-transfer}).
Two evaluations have no training term yet: sufficiency to induce, which no circuit is trained on (median $0.78$--$0.88$ on held-out contrast contexts, $0.35$ for TuringLLM attention outputs; \cref{app:production}), and transfer to weaker, mid-band contexts (\cref{sec:main-result}).
The framework's prescription for them is a term or a training set in their domain, which we have not tested.
% CANDIDATE EVIDENCE, NOT IN THE PAPER (main.tex line 503: the arena results are reported separately). If Daniel wants
% them here, each needs its own sentence and its own appendix or citation:
%  - arena A tuning (experiments/107-true-baselines/arena_a/TUNING.md): on GPT-2 transcoder + z-SAE substrate the
%    necessity hinge (NEC_W; off in production, docs/wcm-specification.md §2.3) was chosen at 0.3 by a pre-registered
%    9-target rule; NEC_W 1.0 raised nec_a (3.tc.3752 0.32 -> 0.95) but pushed A out of band (1.54). The censored read
%    is inert there by construction (k = d_sae).
%  - arena A scope a1 (z + z2 + tc) vs a2 (tc only) and the wikitext-smoke vs OpenWebText corpus test (DAN-208 text;
%    not traced to an output file in this pass).
%  - logit / maximise / behavioural objectives (arena C/D; memory notes say "don't quote yet").
A fit trained for an evaluation is, for that evaluation, scored on home turf (\cref{app:method-comparison}), so the evaluation stays meaningful only if its contexts and donors are never seen in training, as the donor-background test's donors are not (\cref{tab:protocol}), and the evaluations not trained for are reported alongside.

\paragraph{Choosing the sparsity penalty.}
% sec:lambda-transfer, tab:run2026-tl-sweep, tab:run2026-lambda-pilot.
The penalty sets circuit size and, past a point, the pass rate: on the TuringLLM sweep the pass rate holds at $89$--$93\%$ from $\lambda = 2.5 \times 10^{-5}$ to $5 \times 10^{-4}$ while each doubling of $\lambda$ shrinks circuits $1.4$--$1.9\times$, and it declines beyond (\cref{sec:lambda-transfer}).
We recommend fixing the selection rule before looking at pass rates, as our size-only pilot rule did on the public models, where it kept $\lambda = 10^{-3}$ on both (\cref{tab:run2026-lambda-pilot}), and reporting a curve over $\lambda$ rather than one point, since one setting that transfers is not the best setting for each model.

\paragraph{A new model or dictionary.}
Only the model and dictionary adapters change between our three models (\cref{sec:setting}); four things needed checking.
\emph{Precision}: TuringLLM runs in the bfloat16 its pipeline was validated in, but on the public models a bfloat16 data pass did not reproduce the float32 top-context sets, so they run in float32 (\cref{tab:protocol}); a precision check belongs before any fit, since node sets also move with rounding (\cref{app:underdetermination}).
\emph{The dictionary's sparsity}: $k$ sets the fill's budget and the censored read's cut; on a dictionary that keeps every latent ($k = d_{\mathrm{sae}}$) the censored read has no cut to model, and on dense ReLU dictionaries a random-circuit baseline with fitted coefficients stops discriminating, so necessity, not faithfulness, separates real from random circuits there (\cref{app:dense-relu}).
% The k = d_sae statement follows from the definition of tau^(K) (no K-th other latent); docs/wcm-specification.md §7.6.
\emph{Contrast retrieval}: contrast contexts need a pooled representation for every corpus sequence and a silence check, which rejects $14$--$51\%$ of the closest candidates on the targets we inspected (\cref{app:stores}).
\emph{Corpus copies}: duplicated windows silently leak into held-out sets unless pools are deduplicated and split by groups of copies (\cref{app:heldout}).

\paragraph{What to report.}
Because each slot changes the circuits, a result is interpretable only with its recipe: the ablation methods used in training and in evaluation; faithfulness under each method and necessity, on held-out contexts, by depth band and site kind, with intervals over targets; excluded and vacuous targets counted; circuit size with whether coefficients were fitted, and the $\alpha = 1$ rescoring; random-circuit baselines given the same fitting; and the reported-only evaluations (sufficiency to induce, the donor-background test, specificity, mid-band transfer) beside the pass rate (\cref{sec:discussion}).

\paragraph{Choices that matter, and what stays open.}
The ablation method changes not only scores but rankings: the two mean-ablation variants can order small node sets differently (\cref{app:respect-topk}), and in exploratory comparisons a method tended to win the scores of the ablation method it was trained or ranked under (\cref{app:method-comparison}), so comparisons between recipes need the full matrix of recipe against ablation method.
Node sets are not identified by any recipe: a refit with reordered training contexts shares a median Jaccard of $0.39$--$0.51$ with the original and its shared core alone passes for $44$--$72\%$ of targets (\cref{app:underdetermination}), so changes to a slot should be judged by sizes and scores, not by which latents moved.
Finally, the production choices were calibrated on TuringLLM targets that overlap its evaluation set (\cref{sec:setting}), and on the public models every setting except $\lambda$ carried over unchanged (\cref{tab:protocol}, Provenance), so on a new model they are defaults to be checked, not tuned values.

% ---------------------------------------------------------------------------------------------------------------------
% 2026-11-08 (Daniel, via coordinator): the development path as evidence for where WCM can go next. Every number is
% traced; sources in the comment before each step. Panels and protocols differ between steps, so numbers are compared
% only within a step (each step is a paired test against the recipe before it).
% ---------------------------------------------------------------------------------------------------------------------
\subsection{How the Production Recipe Was Reached}
\label{app:framework-history}
The production recipe is the end of a sequence of changes, each made because a run exposed a failure, and each tested against the recipe before it on the same targets, mostly under a selection rule stated before the comparison.
We list them in order as evidence of which kinds of change have paid off; the panels differ between steps, so numbers compare only within a step.

% app:stores (38% fallback), app:contrast-source; experiments/055-close-contrast, 071-contrast-sweep.
\paragraph{Contrast contexts built per target.}
The first contrast contexts came from a retrieval store whose index covered too few sequences, so $38\%$ of TuringLLM targets fell back to one shared list; WCM now selects contrast contexts per target at fit time, ranked by similarity to its strongest contexts and kept only if the target is silent (\cref{app:stores}), and a later comparison of contrast sources showed how much that choice shapes the C term (\cref{app:contrast-source}).
The lesson: inspect the data each training term actually sees.

% app:specificity, app:rank-designs; experiments/056-specificity.
\paragraph{Rank-keep against off-target lifts.}
Reproducing the pre-activation left some circuits restoring their target by lifting its whole neighbourhood; of the terms tried, rank-keep alone kept the gains without the alternatives' costs, and entered the first protocol at a small weight (\cref{app:rank-designs,app:specificity}).

% app:v1-run (15,046 targets, 46%, attention 14%, layers 8-11 22%); app:heldout (66.8%); app:resample (0.1% vs 43.1%,
% 747 targets); fig:failure-anatomy.
\paragraph{The first run at scale.}
Under that protocol, WCM was fitted to $15{,}046$ TuringLLM targets, and $46\%$ passed, $22\%$ at layers 8--11 and $14\%$ for attention outputs (\cref{app:v1-run}).
Analyses of the run exposed three problems: $66.8\%$ of held-out strongest contexts repeated a training context, which led to deduplicated records split by groups of copies (\cref{app:heldout}); with depth, failing circuits increasingly failed only the two mean ablations (\cref{fig:failure-anatomy}); and $0.1\%$ of circuits passed a resample test, against $43.1\%$ under the pass criterion (\cref{app:resample}).

% experiments/093-v2-validation/README.md (out/tables.md): 72 targets, 8 per band x kind, paired, pass floor 0.8, the
% v1 run's as-run held-out sets (leaky; same for both arms): 31 -> 50 of 72, +26.4 pp [15.3, 38.9]; L8-11 2 -> 11 of 24;
% attn 3 -> 12 of 24; median dZ 0.00; size ratio 0.91; fit time 1.27x. Pilot: experiments/090-sparse-fill-training.
% Also main.tex comment at app:respect-topk ("The same fill in training").
\paragraph{Training under the fill that evaluation uses.}
The mean-ablation failures were a mismatch: training replaced non-members with dense means, a code no Top-K dictionary emits, while evaluation re-imposed the Top-K (\cref{app:respect-topk}).
Training the A and C terms under evaluation's own fill and means raised the pass rate on $72$ stratified targets from $31$ to $50$ (paired difference $+26.4$ points, $95\%$ interval $15.3$--$38.9$), from $2$ to $11$ of $24$ at layers 8--11 and from $3$ to $12$ of $24$ for attention outputs, with zero-ablation faithfulness unchanged (median paired change $0.00$), circuits $0.91\times$ the size, and $1.27\times$ the fitting time.
The gain came from aligning training with evaluation, not from a new term.

% app:resample (095 numbers and the production paragraph: 102's 19 -> 39, 48 -> 54 of 72, 1.33x size);
% experiments/102-pass-improvement/README.md line 89 (diagnosis: 15 of 18 Z failures of the previous version restored
% the pre-activation and lost at the Top-K cut); tab:run2026-tl-noR.
\paragraph{A resample term, and a stronger rank-keep.}
To survive real backgrounds, training gained the resample term (above and \cref{app:resample}).
On the previous version's zero-ablation failures, $15$ of $18$ restored the target's pre-activation but lost it at the Top-K cut, so rank-keep was raised from $3 \times 10^{-3}$ to $10^{-1}$ with the resample term at full weight; chosen by a rule stated before the comparison, the pair took a $72$-target check panel from $48$ to $54$ targets passing and from $19$ to $39$ passing the resample test, at $1.33\times$ the circuit size (\cref{app:resample}).
In the production sweep, the arm trained without the resample term passes $5.2$ points less, not significantly (\cref{tab:run2026-tl-noR}).

% experiments/107-true-baselines/latent_pilot/README.md (9 targets enriched for failures: 3 -> 5 of 9, size ratio ~1.04;
% plain geometric: circuits 2-4x larger, sum(lr) 6.3 vs 20; constant lr makes lr/T grow 20x) and
% experiments/107-true-baselines/adaptive_lambda/README.md, "Check" table (arm "geo x 3.15 + censored term
% (latent_pilot)": 63 of 72 vs 54 for the previous version, +10 / -1, size ratio 1.04 on pass-in-both). VERIFY: whether
% that check arm included the 0.1 warm-up (run_tonight.sh stage 4 fits check_geo_x3.15_warmup_censor).
\paragraph{A learning-rate schedule and the censored read.}
Under a constant learning rate, the gate step grows twentyfold as the gate temperature anneals, and some fits degraded late in training.
A geometric decay alone stabilised them but left gates too little total step to close, so circuits grew $2$--$4\times$; raising its peak $3.15\times$ restored the summed step size, and a short warm-up was added for the early spikes the higher peak brought back (\cref{app:loss-details}).
With the censored read added, so that training also sees whether the target survives its site's Top-K, the check panel went from $54$ to $63$ of $72$ targets passing ($10$ gained, $1$ lost) at about the same circuit size (median ratio $1.04$).

% sec:lambda-transfer, tab:run2026-lambda-pilot; experiments/108-batch-size/out/analysis.md (16 TuringLLM targets;
% B16 at 400 steps: 15 / 16 vs 16 / 16, size 0.75x on pass-in-both, 99.9 vs 29.1 fit-minutes = 3.4x; pre-stated rule:
% +-1 target is noise); experiments/107-true-baselines/adaptive_lambda/README.md (every shrinking policy traded passes
% for size on the check panel; losses were necessity failures).
\paragraph{Settings left as they were.}
The penalty was carried to the public models by a size-only rule fixed before a pilot (\cref{sec:lambda-transfer}).
In exploratory tests, batches of $16$ for the same $400$ steps passed $15$ of $16$ TuringLLM targets against $16$, with circuits $0.75\times$ the size where both passed, at $3.4\times$ the fitting time, so the batch of $4$ was kept; and policies that adapted $\lambda$ per target traded passes for size on the check panel, losing passes mostly to necessity, which no training term covers.
Other rejected designs are in \cref{app:membership-studies}.

% Directions: app:production (induce), sec:main-result (mid-band, deep band, attention), app:specificity (amplifiers,
% v1: 30% flagged), app:rank-designs (off-target term: lifts cut 3-7x, but grew circuits), app:resample (anchors
% only; overfiring away from the anchor), adaptive_lambda README (necessity).
\paragraph{Directions not yet taken.}
The same path suggests where the recipe can go next.
Necessity, sufficiency to induce, and transfer to mid-band contexts are scored but not trained for; each could get a term or training contexts in its domain, as the resample test did.
Deep and attention-output targets still pass least (\cref{sec:main-result}), and in the schedule pilot the attention failures kept zero-ablation faithfulness near $0$ under every schedule, so the learning rate does not explain them.
Amplifier circuits survive rank-keep (\cref{app:specificity}); the off-target term cut lifts three- to sevenfold but grew circuits (\cref{app:rank-designs}), so a cheaper specificity term remains open.
And circuits explain the target at its anchors, not at every position (\cref{app:resample}); training on more positions is untested.
\fi\fi
    % ~1.3 pages, its own appendix section
%
% New labels: sec:framework (main), app:framework, tab:framework (appendix). Every other \cref target exists in
% main.tex as of 2026-11-08. Every number below is already printed in main.tex (source comments give the place);
% the one exception, experiment 108, is marked NOT-YET-IN-PAPER where it appears.
% Vocabulary: paper/notes/vocabulary.md (no "endpoint", "floor", "drivers"; "fill" only as main.tex already uses it).
% Never-claim list respected: no "first"; nothing here claims novelty.
% =====================================================================================================================
\providecommand{\frameworkpart}{main}
\def\frameworkmain{main}
\def\frameworkappendix{appendix}
\ifx\frameworkpart\frameworkmain
% ---------------------------------------------------------------------------------------------------------------------
% (1) MAIN-TEXT VERSION. Proposed place: sec:discussion, in place of the paragraph "Ablation methods are part of the
% method" (main.tex lines 548-552), whose evidence it repeats and extends; the paragraph "Recommendations" (lines
% 564-566) can then point here. If it must sit in Method instead, it fits after "Pass criterion" (line 341), before
% \section{Experiments and Results}, but there it would cite results not yet shown.
% Numbers: 2-14% = abstract / run2026_alpha1 (sec:weighted); "one to three orders of magnitude" = app:single-method
% (exploratory); contrast sources = app:contrast-source (protocol v1); resample = app:resample.
% ---------------------------------------------------------------------------------------------------------------------
\subsection{Using WCM as a Framework}
\label{sec:framework}
WCM is better read as a framework for solving for circuits than as one fixed algorithm.
Its parts are interchangeable: an objective (the quantity reproduced, the ablation methods it is reproduced under, the regularisers, and their weights; \cref{eq:loss}), a parametrisation (a gate and a coefficient per node), data (training, contrast, and held-out contexts), and an optimiser.
The production recipe is one setting of these parts, each chosen by an experiment (\cref{tab:framework}), and changing a part changes the circuits found.
With every coefficient fixed at $1$ the same nodes pass for $2$--$14\%$ of targets (\cref{sec:weighted}); in exploratory fits, coefficients trained without zero ablation overshoot it by one to three orders of magnitude (\cref{app:single-method}); circuits trained against close contrast contexts are the most faithful on close backgrounds and the least on any other (\cref{app:contrast-source}); and circuits failed a resample test until training included a resample term (\cref{app:resample}).
When WCM fails an evaluation, we therefore suggest first asking whether training ever saw that evaluation's domain, and if not, adding the term, contexts, or setting that puts it there, while keeping the evaluation's own contexts and donors out of training so that it stays a held-out test (\cref{app:framework}).
A pass rate then describes a recipe, a dictionary, and an evaluation together, not WCM alone, which is why the ablation methods used for training and evaluation belong in the method statement~\citep{mueller2025mib, mahale2026explanation}.
%
\else\ifx\frameworkpart\frameworkappendix
% ---------------------------------------------------------------------------------------------------------------------
% (2) APPENDIX VERSION. Proposed place: a new appendix section directly before "Weighted Circuit Masking Design
% Ablations" (main.tex line 2003, app:membership-studies), so that the design ablations it summarises follow it; or
% directly after app:protocol / app:loss-details (before line 853) if Daniel prefers it next to tab:protocol.
% ---------------------------------------------------------------------------------------------------------------------
\section{Building with Weighted Circuit Masking}
\label{app:framework}
Weighted circuit masking fixes a problem, \labelcref{eq:p2}, and a way of relaxing it, but leaves most of the recipe open.
\Cref{tab:framework} lists its slots: what is reproduced, over which nodes, with which parametrisation, under which ablation methods and backgrounds, with which regularisers and optimiser, from which contexts, and how the result is scored.
The production choice in each slot was made by an experiment reported in these appendices, most of them on TuringLLM and carried over unchanged to the public models (\cref{tab:protocol}, Provenance); none is claimed to be the best setting, and several are known to trade one evaluation against another.

{\footnotesize
\setlength{\tabcolsep}{3pt}
\renewcommand{\arraystretch}{1.15}
\begin{longtable}{@{}p{0.13\linewidth}p{0.33\linewidth}p{0.33\linewidth}p{0.16\linewidth}@{}}
\caption{The slots of weighted circuit masking, the production choice in each, and the alternatives this paper examined or that the formulation admits.}
\label{tab:framework} \\
\toprule
\textbf{Slot} & \textbf{Production choice} & \textbf{Alternatives} & \textbf{Evidence} \\
\midrule
\endfirsthead
\toprule
\textbf{Slot} & \textbf{Production choice} & \textbf{Alternatives} & \textbf{Evidence} \\
\midrule
\endhead
\midrule
\multicolumn{4}{r@{}}{\emph{continued on the next page}} \\
\endfoot
\bottomrule
\endlastfoot
Target & the pre-activation $\preact$ of one SAE latent at its anchors; censored read under Z, C, A & an output metric, as in learned-mask circuits for outputs~\citep{bhaskar2024finding, caples2025scaling}; output- and latent-targeted circuits fail each other's objective (exploratory) & \cref{sec:ablmask}, \cref{app:secondary,app:logit-faith} \\
Nodes & every latent at every SAE site upstream of the target; SAE errors passed through; nodes recomputed & dense ReLU dictionaries (the method replicates; the random baseline needs its own design); members clamped (resample diagnostic) & \cref{sec:problem}, \cref{app:dense-relu,app:resample} \\
Parametrisation & gate $\sigma(\theta/T)$ and coefficient $\mathrm{softplus}(\psi)$ per node, $\alpha_0 = 1$ & $\alpha \equiv 1$ (unweighted circuit masking) & \cref{sec:weighted}, \cref{tab:run2026-alpha1} \\
Training ablations & Z (weight $1$), C and A ($0.25$ each) under evaluation's sparsity-preserving fill; resample D over the training contrast contexts ($1$) & one method alone; dense means; D in place of C; no D & \cref{app:single-method,app:respect-topk,app:resample} \\
Contrast background & $64$ close, silence-checked contexts per target ($48$ train, $16$ held out) & random, distant, or fused sources; the unchecked retrieval store & \cref{app:stores,app:contrast-source} \\
Regularisers & flat penalty $\lambda = 10^{-3}$ per open gate; closed-gate charge $\lambda(1-\sigma(\theta))|\alpha - 1|$; rank-keep $\gamma_{\mathrm{K}} = 10^{-1}$; one shared normaliser $\nu$ & other penalties $\lambda$ (sweep); off-target term, Top-K hinge, pull to rank $1$, firing margin; per-term normalisers & \cref{sec:lambda-transfer}, \cref{app:loss-details,app:rank-designs,app:membership-studies} \\
Optimiser & AdamW, $400$ steps, batch $4$; gate temperature annealed $1 \to 0.05$; learning-rate warm-up, then geometric decay & straight-through binarisation; $200$-step and reach-then-hold schedules; larger batches (exploratory) & \cref{tab:protocol}, \cref{app:membership-studies,app:framework-history} \\
Training contexts & $32$ strongest and $16$ mid-band activating contexts, deduplicated and split by whole groups of copies & a split without deduplication, which leaked copies into held-out sets & \cref{sec:setting}, \cref{app:heldout} \\
Evaluation & held-out strongest contexts; Z, A, C in $[0.75, 1.5]$ and necessity $\geq 0.9$; vacuous targets excluded and counted & reported, not required: mid-band contexts, sufficiency to induce, donor-background test, specificity, $\alpha = 1$ rescoring, random-circuit baselines, stable core & \cref{sec:eval}, \cref{app:protocol,app:nulls,app:underdetermination} \\
\end{longtable}}

\paragraph{Training in the evaluation's domain.}
Across these experiments, circuits did best on the evaluations their training resembled.
When a fit fails an evaluation, the failure is often a property of the recipe, and a slot of \cref{tab:framework} can put the fit in that evaluation's domain.
% app:resample (095: 72 targets, 0% -> 24%, CI +14 to +33; criterion 69% -> 62%, CI -17 to +3; 1.18x size, 1.45x time;
% D in place of C: 51%, CI -29 to -7), production paragraph (exp 102: 19 -> 39 and 48 -> 54 of 72, 1.33x size), T10
% (tab:run2026-tl-noR: +5.2, -0.6 to 11.6).
\emph{Real backgrounds.} Circuits trained only against zero and means failed a resample test almost always; adding a resample term over the training contrast contexts raised the resample pass rate on $72$ targets from $0\%$ to $24\%$, at $1.18\times$ the circuit size and $1.45\times$ the fitting time, and with rank-keep raised it took a $72$-target check panel from $19$ to $39$ targets passing (\cref{app:resample}).
Using it in place of the contrast-mean term gained as much on the resample test but cost the pass criterion, so terms should be added for a new evaluation rather than swapped for an old one; on the production sweep the term's effect on the pass criterion itself is $5.2$ points and not significant (\cref{tab:run2026-tl-noR}).
\emph{Ablation methods and fills.} A mask trained under one ablation method learns what that method rewards: without the zero-ablation term, coefficients overshoot zero ablation by orders of magnitude, and circuits trained against dense means can fail under the sparsity-preserving fill that evaluation uses (\cref{app:single-method,app:respect-topk}).
\emph{Top-K survival.} Reproducing the pre-activation does not keep the target in its site's Top-K; the censored read and rank-keep train for it, though neither removes amplifier circuits (\cref{app:loss-details,app:specificity}).
\emph{Other backgrounds.} Contrast faithfulness generalises furthest when training uses random or distant contrast contexts (\cref{app:contrast-source}).
\emph{Size.} Circuits shrink with $\lambda$ at a nearly constant pass rate up to a plateau's edge (\cref{sec:lambda-transfer}).
Two evaluations have no training term yet: sufficiency to induce, which no circuit is trained on (median $0.78$--$0.88$ on held-out contrast contexts, $0.35$ for TuringLLM attention outputs; \cref{app:production}), and transfer to weaker, mid-band contexts (\cref{sec:main-result}).
The framework's prescription for them is a term or a training set in their domain, which we have not tested.
% CANDIDATE EVIDENCE, NOT IN THE PAPER (main.tex line 503: the arena results are reported separately). If Daniel wants
% them here, each needs its own sentence and its own appendix or citation:
%  - arena A tuning (experiments/107-true-baselines/arena_a/TUNING.md): on GPT-2 transcoder + z-SAE substrate the
%    necessity hinge (NEC_W; off in production, docs/wcm-specification.md §2.3) was chosen at 0.3 by a pre-registered
%    9-target rule; NEC_W 1.0 raised nec_a (3.tc.3752 0.32 -> 0.95) but pushed A out of band (1.54). The censored read
%    is inert there by construction (k = d_sae).
%  - arena A scope a1 (z + z2 + tc) vs a2 (tc only) and the wikitext-smoke vs OpenWebText corpus test (DAN-208 text;
%    not traced to an output file in this pass).
%  - logit / maximise / behavioural objectives (arena C/D; memory notes say "don't quote yet").
A fit trained for an evaluation is, for that evaluation, scored on home turf (\cref{app:method-comparison}), so the evaluation stays meaningful only if its contexts and donors are never seen in training, as the donor-background test's donors are not (\cref{tab:protocol}), and the evaluations not trained for are reported alongside.

\paragraph{Choosing the sparsity penalty.}
% sec:lambda-transfer, tab:run2026-tl-sweep, tab:run2026-lambda-pilot.
The penalty sets circuit size and, past a point, the pass rate: on the TuringLLM sweep the pass rate holds at $89$--$93\%$ from $\lambda = 2.5 \times 10^{-5}$ to $5 \times 10^{-4}$ while each doubling of $\lambda$ shrinks circuits $1.4$--$1.9\times$, and it declines beyond (\cref{sec:lambda-transfer}).
We recommend fixing the selection rule before looking at pass rates, as our size-only pilot rule did on the public models, where it kept $\lambda = 10^{-3}$ on both (\cref{tab:run2026-lambda-pilot}), and reporting a curve over $\lambda$ rather than one point, since one setting that transfers is not the best setting for each model.

\paragraph{A new model or dictionary.}
Only the model and dictionary adapters change between our three models (\cref{sec:setting}); four things needed checking.
\emph{Precision}: TuringLLM runs in the bfloat16 its pipeline was validated in, but on the public models a bfloat16 data pass did not reproduce the float32 top-context sets, so they run in float32 (\cref{tab:protocol}); a precision check belongs before any fit, since node sets also move with rounding (\cref{app:underdetermination}).
\emph{The dictionary's sparsity}: $k$ sets the fill's budget and the censored read's cut; on a dictionary that keeps every latent ($k = d_{\mathrm{sae}}$) the censored read has no cut to model, and on dense ReLU dictionaries a random-circuit baseline with fitted coefficients stops discriminating, so necessity, not faithfulness, separates real from random circuits there (\cref{app:dense-relu}).
% The k = d_sae statement follows from the definition of tau^(K) (no K-th other latent); docs/wcm-specification.md §7.6.
\emph{Contrast retrieval}: contrast contexts need a pooled representation for every corpus sequence and a silence check, which rejects $14$--$51\%$ of the closest candidates on the targets we inspected (\cref{app:stores}).
\emph{Corpus copies}: duplicated windows silently leak into held-out sets unless pools are deduplicated and split by groups of copies (\cref{app:heldout}).

\paragraph{What to report.}
Because each slot changes the circuits, a result is interpretable only with its recipe: the ablation methods used in training and in evaluation; faithfulness under each method and necessity, on held-out contexts, by depth band and site kind, with intervals over targets; excluded and vacuous targets counted; circuit size with whether coefficients were fitted, and the $\alpha = 1$ rescoring; random-circuit baselines given the same fitting; and the reported-only evaluations (sufficiency to induce, the donor-background test, specificity, mid-band transfer) beside the pass rate (\cref{sec:discussion}).

\paragraph{Choices that matter, and what stays open.}
The ablation method changes not only scores but rankings: the two mean-ablation variants can order small node sets differently (\cref{app:respect-topk}), and in exploratory comparisons a method tended to win the scores of the ablation method it was trained or ranked under (\cref{app:method-comparison}), so comparisons between recipes need the full matrix of recipe against ablation method.
Node sets are not identified by any recipe: a refit with reordered training contexts shares a median Jaccard of $0.39$--$0.51$ with the original and its shared core alone passes for $44$--$72\%$ of targets (\cref{app:underdetermination}), so changes to a slot should be judged by sizes and scores, not by which latents moved.
Finally, the production choices were calibrated on TuringLLM targets that overlap its evaluation set (\cref{sec:setting}), and on the public models every setting except $\lambda$ carried over unchanged (\cref{tab:protocol}, Provenance), so on a new model they are defaults to be checked, not tuned values.

% ---------------------------------------------------------------------------------------------------------------------
% 2026-11-08 (Daniel, via coordinator): the development path as evidence for where WCM can go next. Every number is
% traced; sources in the comment before each step. Panels and protocols differ between steps, so numbers are compared
% only within a step (each step is a paired test against the recipe before it).
% ---------------------------------------------------------------------------------------------------------------------
\subsection{How the Production Recipe Was Reached}
\label{app:framework-history}
The production recipe is the end of a sequence of changes, each made because a run exposed a failure, and each tested against the recipe before it on the same targets, mostly under a selection rule stated before the comparison.
We list them in order as evidence of which kinds of change have paid off; the panels differ between steps, so numbers compare only within a step.

% app:stores (38% fallback), app:contrast-source; experiments/055-close-contrast, 071-contrast-sweep.
\paragraph{Contrast contexts built per target.}
The first contrast contexts came from a retrieval store whose index covered too few sequences, so $38\%$ of TuringLLM targets fell back to one shared list; WCM now selects contrast contexts per target at fit time, ranked by similarity to its strongest contexts and kept only if the target is silent (\cref{app:stores}), and a later comparison of contrast sources showed how much that choice shapes the C term (\cref{app:contrast-source}).
The lesson: inspect the data each training term actually sees.

% app:specificity, app:rank-designs; experiments/056-specificity.
\paragraph{Rank-keep against off-target lifts.}
Reproducing the pre-activation left some circuits restoring their target by lifting its whole neighbourhood; of the terms tried, rank-keep alone kept the gains without the alternatives' costs, and entered the first protocol at a small weight (\cref{app:rank-designs,app:specificity}).

% app:v1-run (15,046 targets, 46%, attention 14%, layers 8-11 22%); app:heldout (66.8%); app:resample (0.1% vs 43.1%,
% 747 targets); fig:failure-anatomy.
\paragraph{The first run at scale.}
Under that protocol, WCM was fitted to $15{,}046$ TuringLLM targets, and $46\%$ passed, $22\%$ at layers 8--11 and $14\%$ for attention outputs (\cref{app:v1-run}).
Analyses of the run exposed three problems: $66.8\%$ of held-out strongest contexts repeated a training context, which led to deduplicated records split by groups of copies (\cref{app:heldout}); with depth, failing circuits increasingly failed only the two mean ablations (\cref{fig:failure-anatomy}); and $0.1\%$ of circuits passed a resample test, against $43.1\%$ under the pass criterion (\cref{app:resample}).

% experiments/093-v2-validation/README.md (out/tables.md): 72 targets, 8 per band x kind, paired, pass floor 0.8, the
% v1 run's as-run held-out sets (leaky; same for both arms): 31 -> 50 of 72, +26.4 pp [15.3, 38.9]; L8-11 2 -> 11 of 24;
% attn 3 -> 12 of 24; median dZ 0.00; size ratio 0.91; fit time 1.27x. Pilot: experiments/090-sparse-fill-training.
% Also main.tex comment at app:respect-topk ("The same fill in training").
\paragraph{Training under the fill that evaluation uses.}
The mean-ablation failures were a mismatch: training replaced non-members with dense means, a code no Top-K dictionary emits, while evaluation re-imposed the Top-K (\cref{app:respect-topk}).
Training the A and C terms under evaluation's own fill and means raised the pass rate on $72$ stratified targets from $31$ to $50$ (paired difference $+26.4$ points, $95\%$ interval $15.3$--$38.9$), from $2$ to $11$ of $24$ at layers 8--11 and from $3$ to $12$ of $24$ for attention outputs, with zero-ablation faithfulness unchanged (median paired change $0.00$), circuits $0.91\times$ the size, and $1.27\times$ the fitting time.
The gain came from aligning training with evaluation, not from a new term.

% app:resample (095 numbers and the production paragraph: 102's 19 -> 39, 48 -> 54 of 72, 1.33x size);
% experiments/102-pass-improvement/README.md line 89 (diagnosis: 15 of 18 Z failures of the previous version restored
% the pre-activation and lost at the Top-K cut); tab:run2026-tl-noR.
\paragraph{A resample term, and a stronger rank-keep.}
To survive real backgrounds, training gained the resample term (above and \cref{app:resample}).
On the previous version's zero-ablation failures, $15$ of $18$ restored the target's pre-activation but lost it at the Top-K cut, so rank-keep was raised from $3 \times 10^{-3}$ to $10^{-1}$ with the resample term at full weight; chosen by a rule stated before the comparison, the pair took a $72$-target check panel from $48$ to $54$ targets passing and from $19$ to $39$ passing the resample test, at $1.33\times$ the circuit size (\cref{app:resample}).
In the production sweep, the arm trained without the resample term passes $5.2$ points less, not significantly (\cref{tab:run2026-tl-noR}).

% experiments/107-true-baselines/latent_pilot/README.md (9 targets enriched for failures: 3 -> 5 of 9, size ratio ~1.04;
% plain geometric: circuits 2-4x larger, sum(lr) 6.3 vs 20; constant lr makes lr/T grow 20x) and
% experiments/107-true-baselines/adaptive_lambda/README.md, "Check" table (arm "geo x 3.15 + censored term
% (latent_pilot)": 63 of 72 vs 54 for the previous version, +10 / -1, size ratio 1.04 on pass-in-both). VERIFY: whether
% that check arm included the 0.1 warm-up (run_tonight.sh stage 4 fits check_geo_x3.15_warmup_censor).
\paragraph{A learning-rate schedule and the censored read.}
Under a constant learning rate, the gate step grows twentyfold as the gate temperature anneals, and some fits degraded late in training.
A geometric decay alone stabilised them but left gates too little total step to close, so circuits grew $2$--$4\times$; raising its peak $3.15\times$ restored the summed step size, and a short warm-up was added for the early spikes the higher peak brought back (\cref{app:loss-details}).
With the censored read added, so that training also sees whether the target survives its site's Top-K, the check panel went from $54$ to $63$ of $72$ targets passing ($10$ gained, $1$ lost) at about the same circuit size (median ratio $1.04$).

% sec:lambda-transfer, tab:run2026-lambda-pilot; experiments/108-batch-size/out/analysis.md (16 TuringLLM targets;
% B16 at 400 steps: 15 / 16 vs 16 / 16, size 0.75x on pass-in-both, 99.9 vs 29.1 fit-minutes = 3.4x; pre-stated rule:
% +-1 target is noise); experiments/107-true-baselines/adaptive_lambda/README.md (every shrinking policy traded passes
% for size on the check panel; losses were necessity failures).
\paragraph{Settings left as they were.}
The penalty was carried to the public models by a size-only rule fixed before a pilot (\cref{sec:lambda-transfer}).
In exploratory tests, batches of $16$ for the same $400$ steps passed $15$ of $16$ TuringLLM targets against $16$, with circuits $0.75\times$ the size where both passed, at $3.4\times$ the fitting time, so the batch of $4$ was kept; and policies that adapted $\lambda$ per target traded passes for size on the check panel, losing passes mostly to necessity, which no training term covers.
Other rejected designs are in \cref{app:membership-studies}.

% Directions: app:production (induce), sec:main-result (mid-band, deep band, attention), app:specificity (amplifiers,
% v1: 30% flagged), app:rank-designs (off-target term: lifts cut 3-7x, but grew circuits), app:resample (anchors
% only; overfiring away from the anchor), adaptive_lambda README (necessity).
\paragraph{Directions not yet taken.}
The same path suggests where the recipe can go next.
Necessity, sufficiency to induce, and transfer to mid-band contexts are scored but not trained for; each could get a term or training contexts in its domain, as the resample test did.
Deep and attention-output targets still pass least (\cref{sec:main-result}), and in the schedule pilot the attention failures kept zero-ablation faithfulness near $0$ under every schedule, so the learning rate does not explain them.
Amplifier circuits survive rank-keep (\cref{app:specificity}); the off-target term cut lifts three- to sevenfold but grew circuits (\cref{app:rank-designs}), so a cheaper specificity term remains open.
And circuits explain the target at its anchors, not at every position (\cref{app:resample}); training on more positions is untested.
\fi\fi
